\documentclass[
  justified,
  openright,
]{tufte-book}

% --------------------------------------------------
% Engine sanity (XeLaTeX or LuaLaTeX required)
% --------------------------------------------------
\usepackage{iftex}
\ifPDFTeX
  \PackageError{template}{This template requires XeLaTeX or LuaLaTeX}{}
\fi

% --------------------------------------------------
% Core math
% --------------------------------------------------
\usepackage{amsmath,amssymb,amsfonts}

% --------------------------------------------------
% Tufte: subsubsection is forbidden --- remap safely
% --------------------------------------------------
\makeatletter
\let\subsubsection\paragraph
\let\paragraph\paragraph
\makeatother

% --------------------------------------------------
% Fix Pandoc tabularray counter=none bug
% --------------------------------------------------
\makeatletter
\ExplSyntaxOn
\cs_if_exist:NF \c@none { \newcounter{none} }
\ExplSyntaxOff
\makeatother

\usepackage{fontspec}

\IfFontExistsTF{TeX Gyre Pagella}
  {\setmainfont{TeX Gyre Pagella}}
  {\setmainfont{Latin Modern Roman}}

\IfFontExistsTF{TeX Gyre Heros}
  {\setsansfont[Scale=0.92]{TeX Gyre Heros}}
  {\setsansfont{Latin Modern Sans}}

\IfFontExistsTF{TeX Gyre Cursor}
  {\setmonofont[Scale=0.82]{TeX Gyre Cursor}}
  {\setmonofont{Latin Modern Mono}}
% --------------------------------------------------
% Graphics and tables
% --------------------------------------------------
\usepackage{graphicx}
\usepackage{booktabs}
\usepackage{longtable}
\usepackage{array}
\usepackage{tabularray}
\UseTblrLibrary{booktabs}

\setkeys{Gin}{width=\maxwidth,height=\maxheight,keepaspectratio}

% --------------------------------------------------
% Colors
% --------------------------------------------------
\usepackage[table,xcdraw,svgnames]{xcolor}

\definecolor{DarkGray}{RGB}{55,60,70}
\definecolor{CalloutBlue}{HTML}{DCEEFF}
\definecolor{CalloutBorder}{HTML}{5B8BD9}
\definecolor{BrandBlue}{RGB}{18,65,124}
\definecolor{BrandAccent}{RGB}{0,153,204}
\definecolor{LightGray}{RGB}{245,246,248}
\definecolor{MidGray}{RGB}{175,180,190}
\definecolor{WarningAmber}{RGB}{230,100,0}
\definecolor{WarningBg}{RGB}{255,248,235}
\definecolor{CautionYellow}{RGB}{180,140,0}
\definecolor{CautionBg}{RGB}{255,252,230}
\definecolor{TipGreen}{RGB}{36,110,44}
\definecolor{TipBg}{RGB}{237,248,239}
\definecolor{NoteBg}{RGB}{235,243,255}
\definecolor{NoteBlue}{RGB}{18,90,160}
\definecolor{DangerRed}{RGB}{185,28,28}
\definecolor{DangerBg}{RGB}{255,237,237}
\definecolor{TableHeader}{RGB}{18,65,124}
\definecolor{TableRowAlt}{RGB}{235,241,252}
\definecolor{MarginRule}{RGB}{210,215,225}
\definecolor{SpineAccent}{RGB}{0,153,204}
% --------------------------------------------------
% Pandoc syntax highlighting macros
% --------------------------------------------------
\newcommand{\AlertTok}[1]{\textcolor{Red}{#1}}
\newcommand{\AnnotationTok}[1]{\textcolor{Gray}{#1}}
\newcommand{\AttributeTok}[1]{\textcolor{DarkSlateGray}{#1}}
\newcommand{\BaseNTok}[1]{\textcolor{Purple}{#1}}
\newcommand{\BuiltInTok}[1]{#1}
\newcommand{\CharTok}[1]{\textcolor{Teal}{#1}}
\newcommand{\CommentTok}[1]{\textcolor{MidGray}{#1}}
\newcommand{\CommentVarTok}[1]{\textcolor{MidGray}{#1}}
\newcommand{\ConstantTok}[1]{\textcolor{Purple}{#1}}
\newcommand{\ControlFlowTok}[1]{\textcolor{BrandBlue}{#1}}
\newcommand{\DataTypeTok}[1]{\textcolor{Purple}{#1}}
\newcommand{\DecValTok}[1]{\textcolor{Purple}{#1}}
\newcommand{\DocumentationTok}[1]{\textcolor{MidGray}{#1}}
\newcommand{\ErrorTok}[1]{\textcolor{Red}{#1}}
\newcommand{\ExtensionTok}[1]{#1}
\newcommand{\FloatTok}[1]{\textcolor{Purple}{#1}}
\newcommand{\FunctionTok}[1]{\textcolor{BrandBlue}{#1}}
\newcommand{\ImportTok}[1]{\textcolor{BrandAccent}{#1}}
\newcommand{\InformationTok}[1]{\textcolor{BrandAccent}{#1}}
\newcommand{\KeywordTok}[1]{\textcolor{BrandBlue}{#1}}
\newcommand{\NormalTok}[1]{#1}
\newcommand{\OperatorTok}[1]{\textcolor{DarkGray}{#1}}
\newcommand{\OtherTok}[1]{#1}
\newcommand{\PreprocessorTok}[1]{\textcolor{BrandAccent}{#1}}
\newcommand{\RegionMarkerTok}[1]{#1}
\newcommand{\SpecialCharTok}[1]{\textcolor{Teal}{#1}}
\newcommand{\SpecialStringTok}[1]{\textcolor{Teal}{#1}}
\newcommand{\StringTok}[1]{\textcolor{Teal}{#1}}
\newcommand{\VariableTok}[1]{#1}
\newcommand{\VerbatimStringTok}[1]{\textcolor{Teal}{#1}}
\newcommand{\WarningTok}[1]{\textcolor{WarningAmber}{#1}}
% --------------------------------------------------
% Ensure ALL Pandoc *Tok macros are defined (robust)
% --------------------------------------------------
\makeatletter
\@for\tok:=Alert,Annotation,Attribute,BaseN,BuiltIn,Char,Comment,CommentVar,%
Constant,ControlFlow,DataType,DecVal,Documentation,Error,Extension,Float,%
Function,Import,Information,Keyword,Normal,Operator,Other,Preprocessor,%
RegionMarker,SpecialChar,SpecialString,String,Variable,VerbatimString,Warning%
\do{%
  \expandafter\providecommand\csname\tok Tok\endcsname[1]{#1}%
}
\makeatother

% --------------------------------------------------
% Pandoc code block support (Shaded environment)
% --------------------------------------------------
\usepackage{fancyvrb}
\usepackage{framed}

\definecolor{shadecolor}{RGB}{248,248,248}

\newenvironment{Shaded}{%
  \begin{snugshade}
}{%
  \end{snugshade}
}

% --------------------------------------------------
% Pandoc Highlighting environment
% --------------------------------------------------
\DefineVerbatimEnvironment{Highlighting}{Verbatim}{
  commandchars=\\\{\},
  fontsize=\small
}

% --------------------------------------------------
% Tufte: subsubsection is forbidden --- remap safely
% --------------------------------------------------
\makeatletter
\let\subsubsection\paragraph
\let\paragraph\paragraph
\makeatother

% --------------------------------------------------
% Utilities
% --------------------------------------------------
\usepackage{etoolbox}
\usepackage{xparse}
\usepackage{ifthen}
\usepackage{textcase}
\usepackage{titlesec}
% --------------------------------------------------
% Fix Pandoc tabularray counter=none bug
% --------------------------------------------------
\makeatletter
\ExplSyntaxOn
\cs_if_exist:NF \c@none { \newcounter{none} }
\ExplSyntaxOff
\makeatother
% --------------------------------------------------
% Paragraph & section spacing
% --------------------------------------------------
\setlength{\parindent}{1em}
\setlength{\parskip}{0pt}

\titlespacing*{\section}{0pt}{2.0ex plus .2ex}{1.0ex}
\titlespacing*{\subsection}{0pt}{1.5ex plus .2ex}{0.8ex}

% --------------------------------------------------
% Margin notes (tufte)
% --------------------------------------------------
\setlength{\marginparwidth}{2.2in}

\newcommand{\mn}[1]{%
  \marginnote{\raggedright\footnotesize\color{DarkGray}#1}
}

\NewDocumentCommand{\margincallout}{m}{%
  \marginnote{%
    \setlength{\fboxsep}{6pt}%
    \colorbox{CalloutBlue}{%
      \begin{minipage}{\marginparwidth}
      \footnotesize #1
      \end{minipage}
    }
  }
}

% --------------------------------------------------
% Soul / tabularray / marking fixes
% --------------------------------------------------
\usepackage{soul}
\AfterEndEnvironment{longtblr}{\soulsuspend}

\makeatletter
\providecommand{\soulsuspend}{}
\providecommand{\soulsuspend*}{}
\providecommand{\soulsuspend@}{}
\makeatother

\DeclareRobustCommand{\checkbox}{\(\square\)}

\soulregister\checkbox{0}
\soulregister\textbf{1}
\soulregister\emph{1}
\soulregister\href{2}
\soulregister\label{0}
\soulregister\ref{0}
\soulregister\pageref{0}
\soulregister\footnote{1}
\soulregister\cite{0}

% --------------------------------------------------
% Headers / marks (no case mangling)
% --------------------------------------------------
\makeatletter
\renewcommand{\chaptermark}[1]{\markboth{#1}{}}
\renewcommand{\sectionmark}[1]{\markright{#1}}

\pretocmd\markboth{\let\allcapsspacing\relax\let\smallcapsspacing\relax}{}{}
\pretocmd\markright{\let\allcapsspacing\relax\let\smallcapsspacing\relax}{}{}
\makeatother

% --------------------------------------------------
% Blank page helper
% --------------------------------------------------
\newcommand{\blankpage}{%
  \clearpage
  \thispagestyle{empty}
  \null\vfill
  \begin{center}
    \color{MidGray}\sffamily\small\itshape
    (This page intentionally left blank.)
  \end{center}
  \vfill
  \clearpage
}

% --------------------------------------------------
% Fancyhdr
% --------------------------------------------------
\usepackage{fancyhdr}
\pagestyle{fancy}
\fancyhf{}
\fancyhead[LE]{\small\scshape\leftmark}
\fancyhead[RO]{\small\scshape\rightmark}
\fancyfoot[C]{\thepage}
\AtBeginDocument{\pagestyle{fancy}}

% --------------------------------------------------
% Hyperlinks
% --------------------------------------------------
\usepackage{hyperref}
\hypersetup{
  colorlinks=true,
  linkcolor=blue,
  urlcolor=blue,
  citecolor=blue,
  hypertexnames=false
}

% --------------------------------------------------
% Pandoc helpers
% --------------------------------------------------
\providecommand{\tightlist}{%
  \setlength{\itemsep}{0pt}\setlength{\parskip}{0pt}}

% Show only parts and chapters in TOC; number only parts
\setcounter{tocdepth}{1}
\setcounter{secnumdepth}{-1}

% --------------------------------------------------
% Metadata
% --------------------------------------------------



% --------------------------------------------------
% ============================================================
%  DOCUMENT
% ============================================================
\begin{document}
% -- Front matter -------------------------------------------
\frontmatter
\pagestyle{fancy}

% 1) Blank half title page
\thispagestyle{empty}
\null
\clearpage

% 2) Frontispiece (placeholder)
\thispagestyle{empty}
\begin{center}
  \vspace*{0.35\textheight}
  {\sffamily\large Frontispiece}
\end{center}
\clearpage

% 3) Full title page
\begin{titlepage}
    \vspace*{1.6in}
    {\sffamily\fontsize{9}{11}\selectfont\MakeUppercase{Training Manual}}
    \vspace{0.3in}
    {\color{BrandAccent}\rule{4.5in}{2pt}}
    \vspace{0.25in}
    % ADD YOUR TITLE TEXT HERE:
    {\sffamily\bfseries\fontsize{32}{38}\selectfont Coding With //ScreenReaders}
    \vspace{0.2in}
    {\color{BrandAccent}\rule{4.5in}{0.6pt}}
    \vspace{0.5in}
    \vfill
    {\sffamily\fontsize{7.5}{9}\selectfont\color{MidGray}%
      CONFIDENTIAL --- For Internal Use Only%
    }
  \clearpage
  \nopagecolor
\end{titlepage}

% 4) Copyright page (placeholder)
\thispagestyle{empty}
\vspace*{0.25\textheight}
{\small
\noindent Copyright \textcopyright\ \the\year\ Michael Ryan Hunsaker, M.Ed., Ph.D. Coding With ScreenReaders is licensed under a Creative Commons Attribution-NoNCommerical-NoDerivatives 4.0 International: https://creativecommons.org/licenses/by-nc-nd/4.0/
\noindent All rights reserved.
\noindent NO AI TRAINING: Without in any way limiting the author's rights under copyright, any use of this publication to “train” generative artificial intelligence (AI) technologies to generate text is expressly prohibited.
}
\clearpage

% 5) Contents
{
\hypersetup{linkcolor=BrandBlue}
\tableofcontents
\clearpage
}

% 6) Dedication (placeholder)
\chapter*{Dedication}
\thispagestyle{empty}
\begin{center}
  \vspace*{0.3\textheight}
  {\itshape To the countless individuals who have been told they lack the capability to code, yet remain determined to learn and create through accessible coding practices, this book is dedicated.}\\[1em]
\end{center}
\clearpage

% 7) Epigraph (placeholder)
\chapter*{Epigraph}
\thispagestyle{empty}
\begin{flushright}
  \vspace*{0.25\textheight}
  {\itshape The power of the Web is in its universality. Access by everyone regardless of disability is an essential aspect}\\[1em]
  --- Tim Berners-Lee, Inventor of the World Wide Web
\end{flushright}
\clearpage

% 8) Introduction (placeholder)
\newpage

% -- Main content -------------------------------------------
\mainmatter
\pagenumbering{arabic} % Switch to numbers (1, 2, 3...) for the main content


\section{Teacher's Guide: Introduction to Computer Science with
Python}\label{teachers-guide-introduction-to-computer-science-with-python}

\subsection{For Students with Visual Impairments --- Complete Instructor
Edition}\label{for-students-with-visual-impairments-complete-instructor-edition}

\begin{center}\rule{0.5\linewidth}{0.5pt}\end{center}

\subsection{Table of Contents}\label{table-of-contents}

\begin{enumerate}
\def\labelenumi{\arabic{enumi}.}
\tightlist
\item
  \hyperref[philosophy-and-foundational-principles]{Philosophy and
  Foundational Principles}
\item
  \hyperref[understanding-your-students]{Understanding Your Students}
\item
  \hyperref[classroom-setup-and-technology-configuration]{Classroom
  Setup and Technology Configuration}
\item
  \hyperref[universal-spoken-code-conventions]{Universal Spoken Code
  Conventions}
\item
  \hyperref[unit-0-setup-and-environment]{Unit 0 --- Setup and
  Environment}
\item
  \hyperref[unit-1-foundations-and-getting-started]{Unit 1 ---
  Foundations and Getting Started}
\item
  \hyperref[unit-2-numeric-computations-and-io]{Unit 2 --- Numeric
  Computations and I/O}
\item
  \hyperref[unit-3-control-flow-decisions-and-loops]{Unit 3 --- Control
  Flow, Decisions, and Loops}
\item
  \hyperref[unit-4-data-structures-lists-and-dictionaries]{Unit 4 ---
  Data Structures: Lists and Dictionaries}
\item
  \hyperref[unit-5-user-input-and-while-loops]{Unit 5 --- User Input and
  While Loops}
\item
  \hyperref[unit-6-functions]{Unit 6 --- Functions}
\item
  \hyperref[unit-7-classes-and-object-oriented-programming]{Unit 7 ---
  Classes and Object-Oriented Programming}
\item
  \hyperref[unit-8-files-and-exceptions]{Unit 8 --- Files and
  Exceptions}
\item
  \hyperref[unit-9-testing-and-data-analysis]{Unit 9 --- Testing and
  Data Analysis}
\item
  \hyperref[unit-10-project-based-learning]{Unit 10 --- Project-Based
  Learning}
\item
  \hyperref[unit-11-capstone-project]{Unit 11 --- Capstone Project}
\item
  \hyperref[assessment-strategies]{Assessment Strategies}
\item
  \hyperref[supporting-diverse-needs-within-the-visually-impaired-population]{Supporting
  Diverse Needs Within the Visually Impaired Population}
\item
  \hyperref[teaching-ai-tools-responsibly]{Teaching AI Tools
  Responsibly}
\item
  \hyperref[alternative-and-supplementary-resources]{Alternative and
  Supplementary Resources}
\item
  \hyperref[quick-reference-sheets]{Quick Reference Sheets}
\end{enumerate}

\begin{center}\rule{0.5\linewidth}{0.5pt}\end{center}

\subsection{1. Philosophy and Foundational
Principles}\label{philosophy-and-foundational-principles}

\subsubsection{Why Python for Blind
Students?}\label{why-python-for-blind-students}

Python is widely recognized as one of the most accessible first
programming languages for visually impaired learners for several
concrete reasons:

\begin{itemize}
\tightlist
\item
  \textbf{Text-based I/O is the default.} Unlike languages that push
  beginners toward GUI frameworks or graphical plotting, Python's most
  natural outputs are plain text to the terminal --- fully accessible to
  any screen reader.
\item
  \textbf{Indentation enforces structure.} While indentation introduces
  a learning challenge (addressed in detail below), it also means that
  code read aloud by a screen reader carries structural meaning in the
  form of announced indent levels --- something curly-brace languages do
  not provide.
\item
  \textbf{Error messages are descriptive.} Python's tracebacks include
  the file name, line number, error type, and a message. A screen reader
  user can navigate directly to the problem without visual scanning.
\item
  \textbf{The standard library is enormous.} Python ships with modules
  for file I/O, JSON, math, random numbers, CSV parsing, unit testing,
  and more --- reducing the complexity of dependency management in a
  learning context.
\item
  \textbf{The REPL is an accessible sandbox.} Python's interactive
  interpreter provides immediate spoken feedback for any expression,
  enabling rapid experimentation without the overhead of creating and
  running files.
\end{itemize}

\subsubsection{Core Pedagogical
Commitments}\label{core-pedagogical-commitments}

\textbf{Commit to ears-first instruction.} Every demonstration, every
example, every debugging session should be conducted as if the teacher
is also using a screen reader. Ask yourself before every lesson: ``If I
could not see the screen, how would I know what was happening right
now?'' If the answer is unclear, redesign the demonstration.

\textbf{Commit to explicit navigation.} Sighted students can glance at
code and immediately locate structures. Blind students must navigate
deliberately using keyboard shortcuts and screen reader commands. Teach
navigation skills as primary content, not auxiliary support.

\textbf{Commit to verbal precision.} Vague instructions like ``look at
the next line'' or ``you can see here that\ldots{}'' create invisible
barriers. Replace all visual language with navigation language: ``Move
down two lines,'' ``Jump to line 14 with Ctrl+G,'' ``Navigate to the end
of the file with Ctrl+End.''

\textbf{Commit to error-positive culture.} For a sighted student, an
error is a red underline they can ignore or address at leisure. For a
screen reader user, an unaddressed error is noise that clutters the
output of every subsequent navigation. Treat errors as immediate
priorities and build the habit of reading and resolving them before
moving on.

\textbf{Commit to redundancy.} Provide concepts in multiple formats ---
spoken explanation, written pseudocode, physical analogy, and code.
Blind students cannot fall back on visual re-reading; redundant
representation builds retention.

\begin{center}\rule{0.5\linewidth}{0.5pt}\end{center}

\subsection{2. Understanding Your
Students}\label{understanding-your-students}

\subsubsection{The Heterogeneity of Visual
Impairment}\label{the-heterogeneity-of-visual-impairment}

``Visually impaired'' encompasses a wide spectrum. Your students may
include:

{\def\LTcaptype{none} % do not increment counter
\begin{longtable}[]{@{}
  >{\raggedright\arraybackslash}p{(\linewidth - 2\tabcolsep) * \real{0.2533}}
  >{\raggedright\arraybackslash}p{(\linewidth - 2\tabcolsep) * \real{0.7467}}@{}}
\toprule\noalign{}
\begin{minipage}[b]{\linewidth}\raggedright
Profile
\end{minipage} & \begin{minipage}[b]{\linewidth}\raggedright
Implications for Teaching
\end{minipage} \\
\midrule\noalign{}
\endhead
\bottomrule\noalign{}
\endlastfoot
\textbf{Congenitally blind} (blind from birth) & May have no visual
spatial metaphors. Analogies must be tactile, auditory, or conceptual.
May be faster with screen readers from years of experience. \\
\textbf{Late-onset blindness} (lost vision after sighted years) & May
retain strong visual mental models. Can use these as bridges but should
not be required to. May still be learning screen reader skills. \\
\textbf{Low vision} (functional but impaired vision) & May use a
combination of magnification and screen reading. May be able to see some
screen content but cannot rely on fine detail. Do not assume screen-only
or vision-only. \\
\textbf{Color blindness} & Primarily relevant for syntax highlighting.
Ensure themes are not relied upon for meaning. \\
\textbf{Progressive conditions} & Vision may change during the course.
Build habits that remain useful at any level of vision. \\
\end{longtable}
}

Never assume a student's profile without asking. Some students are
private about their diagnoses; frame questions around what tools they
use and what works best, not what their condition is.

\subsubsection{Screen Reader Literacy
Varies}\label{screen-reader-literacy-varies}

Some students arrive with years of fluent screen reader experience.
Others are relatively new to their screen reader and may still be
developing efficiency. Assess screen reader fluency separately from
programming aptitude --- a student who struggles to navigate a file is
not struggling with Python; they need more time with their tool.

Suggest that students spend time with their screen reader in
non-programming contexts (browsing documentation, reading formatted
text) before or alongside the first programming units.

\subsubsection{Braille Display Users}\label{braille-display-users}

Some students use a refreshable braille display in addition to or
instead of audio output. Key considerations:

\begin{itemize}
\tightlist
\item
  Braille displays have limited width (typically 40 or 80 cells). Long
  lines of code may need to be panned across. Encourage keeping line
  length under 79 characters (PEP 8 standard) for this reason.
\item
  Code rendered in braille uses Grade 1 (uncontracted) braille because
  contractions would misrepresent code characters. Confirm this setting
  is active.
\item
  Navigation with a braille display typically involves routing cursor
  buttons --- physical buttons above each braille cell that move the
  cursor to that position. Students using these may navigate differently
  than audio-only users.
\item
  Braille display users often prefer lower verbosity audio output since
  they read braille silently. Do not require audio announcements as the
  only feedback mechanism.
\end{itemize}

\subsubsection{Students Who Are Also
Deaf-Blind}\label{students-who-are-also-deaf-blind}

Students who are deaf-blind rely exclusively on braille for all
information. Audio feedback from the screen reader is not available.
Considerations:

\begin{itemize}
\tightlist
\item
  All instructional communication must be via tactile means, text-based
  communication, or tactile sign language interpreters.
\item
  Braille output of code is the primary interface. Pay special attention
  to line length and symbol representation.
\item
  Terminal output must be navigable via braille display --- avoid
  programs that produce output faster than a braille display can be
  read.
\item
  Refer to specialists in deaf-blind education for specific
  accommodation strategies beyond the scope of this guide.
\end{itemize}

\begin{center}\rule{0.5\linewidth}{0.5pt}\end{center}

\subsection{3. Classroom Setup and Technology
Configuration}\label{classroom-setup-and-technology-configuration}

\subsubsection{Physical Environment}\label{physical-environment}

\begin{itemize}
\tightlist
\item
  Ensure the student's workstation does not have excessive background
  noise that interferes with screen reader audio. Headphones are often
  preferable in a shared classroom.
\item
  Position the workstation so the teacher can sit beside (not in front
  of) the student, allowing them to listen to the screen reader output
  together.
\item
  If a classroom assistant or orientation and mobility specialist is
  present, coordinate roles in advance. The programming teacher leads on
  programming content; the VI specialist leads on accessibility tool
  configuration.
\end{itemize}

\subsubsection{VS Code Configuration
Checklist}\label{vs-code-configuration-checklist}

Work through this checklist with the student on day one:

\textbf{Screen Reader}

\begin{itemize}
\tightlist
\item
  \begin{itemize}
  \tightlist
  \item[$\square$]
    Screen Reader Optimized Mode ON (Ctrl+Shift+P -\textgreater{} Toggle
    Screen Reader Optimized)
  \end{itemize}
\item
  \begin{itemize}
  \tightlist
  \item[$\square$]
    Accessibility help panel reviewed (Shift+Alt+F1)
  \end{itemize}
\end{itemize}

\textbf{Editor Settings (Ctrl+,)}

\begin{itemize}
\tightlist
\item
  \begin{itemize}
  \tightlist
  \item[$\square$]
    editor.lineNumbers = ``on''
  \end{itemize}
\item
  \begin{itemize}
  \tightlist
  \item[$\square$]
    editor.renderWhitespace = ``all'' or ``boundary''
  \end{itemize}
\item
  \begin{itemize}
  \tightlist
  \item[$\square$]
    editor.wordWrap = ``on'' (for long lines)
  \end{itemize}
\item
  \begin{itemize}
  \tightlist
  \item[$\square$]
    workbench.reduceMotion = ``on'' (reduces distracting animations)
  \end{itemize}
\item
  \begin{itemize}
  \tightlist
  \item[$\square$]
    editor.tabSize = 4
  \end{itemize}
\item
  \begin{itemize}
  \tightlist
  \item[$\square$]
    editor.insertSpaces = true (spaces not tabs for Python)
  \end{itemize}
\end{itemize}

\textbf{Screen Reader Indentation (in screen reader settings, not VS
Code)}

\begin{itemize}
\tightlist
\item
  \begin{itemize}
  \tightlist
  \item[$\square$]
    NVDA: Document Formatting -\textgreater{} Indentation
    -\textgreater{} ``Say all'' or ``Beeps''
  \end{itemize}
\item
  \begin{itemize}
  \tightlist
  \item[$\square$]
    JAWS: Configure speech for VS Code environment
  \end{itemize}
\item
  \begin{itemize}
  \tightlist
  \item[$\square$]
    VoiceOver: Verbosity -\textgreater{} Text -\textgreater{} announce
    indentation
  \end{itemize}
\item
  \begin{itemize}
  \tightlist
  \item[$\square$]
    Narrator: No built-in indent announcement; use Beep plugin or
    extension
  \end{itemize}
\end{itemize}

\textbf{Extensions Installed}

\begin{itemize}
\tightlist
\item
  \begin{itemize}
  \tightlist
  \item[$\square$]
    Python (Microsoft)
  \end{itemize}
\item
  \begin{itemize}
  \tightlist
  \item[$\square$]
    Pylance
  \end{itemize}
\item
  \begin{itemize}
  \tightlist
  \item[$\square$]
    Python Indent (for consistent auto-indentation)
  \end{itemize}
\item
  \begin{itemize}
  \tightlist
  \item[$\square$]
    GitLens or Git Graph (for version control visibility)
  \end{itemize}
\item
  \begin{itemize}
  \tightlist
  \item[$\square$]
    Error Lens (inline error annotations --- screen-reader-readable)
  \end{itemize}
\end{itemize}

\textbf{Keyboard Shortcuts Practiced}

\begin{itemize}
\tightlist
\item
  \begin{itemize}
  \tightlist
  \item[$\square$]
    Ctrl+G --- Go to line
  \end{itemize}
\item
  \begin{itemize}
  \tightlist
  \item[$\square$]
    F8 --- Next error
  \end{itemize}
\item
  \begin{itemize}
  \tightlist
  \item[$\square$]
    Shift+F8 --- Previous error
  \end{itemize}
\item
  \begin{itemize}
  \tightlist
  \item[$\square$]
    Ctrl+Shift+M --- Problems panel
  \end{itemize}
\item
  \begin{itemize}
  \tightlist
  \item[$\square$]
    Shift+Alt+F2 --- Accessible View
  \end{itemize}
\item
  \begin{itemize}
  \tightlist
  \item[$\square$]
    Ctrl+Shift+P --- Command Palette
  \end{itemize}
\item
  \begin{itemize}
  \tightlist
  \item[$\square$]
    Ctrl+` --- Toggle terminal
  \end{itemize}
\item
  \begin{itemize}
  \tightlist
  \item[$\square$]
    Ctrl+Shift+E --- Explorer
  \end{itemize}
\item
  \begin{itemize}
  \tightlist
  \item[$\square$]
    Ctrl+Shift+O --- Go to symbol (function/class navigation)
  \end{itemize}
\item
  \begin{itemize}
  \tightlist
  \item[$\square$]
    Alt+F12 --- Peek Definition
  \end{itemize}
\item
  \begin{itemize}
  \tightlist
  \item[$\square$]
    F12 --- Go to Definition
  \end{itemize}
\item
  \begin{itemize}
  \tightlist
  \item[$\square$]
    Shift+F12 --- Find all References
  \end{itemize}
\end{itemize}

\subsubsection{Terminal Configuration}\label{terminal-configuration}

\begin{itemize}
\tightlist
\item
  \textbf{Windows Terminal} (preferred over legacy cmd.exe on Windows)
  has better accessibility with NVDA and JAWS.
\item
  \textbf{NVDA + Windows Terminal:} Enable ``Use UI Automation'' in NVDA
  settings for better terminal reading.
\item
  \textbf{VoiceOver + macOS Terminal:} Works well by default. Use
  \texttt{VO+Shift+Down} to interact with terminal content.
\item
  \textbf{Orca + Linux terminal:} Orca reads terminal output well. Use
  \texttt{Orca+T} (flat review) to re-read recent output.
\item
  Configure the terminal font to a monospace font at a size comfortable
  for any residual vision.
\end{itemize}

\subsubsection{Python Virtual Environment Setup for
Accessibility}\label{python-virtual-environment-setup-for-accessibility}

\begin{Shaded}
\begin{Highlighting}[]
\CommentTok{\# Create environment with a consistent name the student will always type}
\ExtensionTok{python3} \AttributeTok{{-}m}\NormalTok{ venv .venv}

\CommentTok{\# Activation {-}{-}{-} save this as a text file the student can open and copy from}
\CommentTok{\# Windows: .venv\textbackslash{}Scripts\textbackslash{}activate}
\CommentTok{\# macOS/Linux: source .venv/bin/activate}

\CommentTok{\# Verify with a spoken command}
\ExtensionTok{python3} \AttributeTok{{-}{-}version}
\ExtensionTok{pip}\NormalTok{ list}
\end{Highlighting}
\end{Shaded}

Consider creating a startup script (\texttt{start.bat} on Windows,
\texttt{start.sh} on macOS/Linux) that activates the environment and
opens VS Code. Students can run one command to get their environment
ready, reducing setup cognitive load.

\begin{center}\rule{0.5\linewidth}{0.5pt}\end{center}

\subsection{4. Universal Spoken Code
Conventions}\label{universal-spoken-code-conventions}

Establish this vocabulary on the first day and reinforce it
consistently. When a student is reading code aloud (for debugging, peer
review, or just narrating their work), they should use these terms
without hesitation.

\subsubsection{Symbol Pronunciation
Table}\label{symbol-pronunciation-table}

{\def\LTcaptype{none} % do not increment counter
\begin{longtable}[]{@{}
  >{\raggedright\arraybackslash}p{(\linewidth - 4\tabcolsep) * \real{0.2500}}
  >{\raggedright\arraybackslash}p{(\linewidth - 4\tabcolsep) * \real{0.3696}}
  >{\raggedright\arraybackslash}p{(\linewidth - 4\tabcolsep) * \real{0.3804}}@{}}
\toprule\noalign{}
\begin{minipage}[b]{\linewidth}\raggedright
Written Symbol
\end{minipage} & \begin{minipage}[b]{\linewidth}\raggedright
Spoken Form
\end{minipage} & \begin{minipage}[b]{\linewidth}\raggedright
Notes
\end{minipage} \\
\midrule\noalign{}
\endhead
\bottomrule\noalign{}
\endlastfoot
\texttt{=} & ``gets'' & Assignment only \\
\texttt{==} & ``equals equals'' or ``is equal to'' & Comparison \\
\texttt{!=} & ``not equal to'' or ``bang equals'' & \\
\texttt{\textgreater{}=} & ``greater than or equal to'' & \\
\texttt{\textless{}=} & ``less than or equal to'' & \\
\texttt{\textgreater{}} & ``greater than'' & \\
\texttt{\textless{}} & ``less than'' & \\
\texttt{+=} & ``plus equals'' or ``add and assign'' & \\
\texttt{-=} & ``minus equals'' & \\
\texttt{*=} & ``times equals'' & \\
\texttt{**} & ``to the power of'' or ``double star'' & \\
\texttt{//} & ``floor divide'' or ``integer divide'' & \\
\texttt{\%} & ``modulo'' or ``remainder'' & \\
\texttt{()} & ``open paren \ldots{} close paren'' & Always say both \\
\texttt{{[}{]}} & ``open bracket \ldots{} close bracket'' & \\
\texttt{\{\}} & ``open brace \ldots{} close brace'' & \\
\texttt{:} & ``colon'' & End of if/for/def/class line \\
\texttt{"text"} & ``quote text quote'' & \\
\texttt{\textquotesingle{}text\textquotesingle{}} & ``single-quote text
single-quote'' & \\
\texttt{f"..."} & ``f-string'' & Mention the f prefix \\
\texttt{\{var\}} inside f-string & ``open brace var close brace'' & \\
\texttt{\#} & ``hash'' & Followed by comment text \\
\texttt{\textbackslash{}n} & ``backslash n'' or ``newline'' & \\
\texttt{\textbackslash{}t} & ``backslash t'' or ``tab'' & \\
\texttt{\_} & ``underscore'' & In names: ``my underscore name'' \\
\texttt{\_\_init\_\_} & ``dunder init'' & Double-underscore methods \\
\texttt{\_\_str\_\_} & ``dunder str'' & \\
\texttt{.} & ``dot'' & In \texttt{math.sqrt}, say ``math dot sqrt'' \\
\texttt{-\textgreater{}} & ``arrow'' & Type hints \\
\texttt{*args} & ``star args'' & \\
\texttt{**kwargs} & ``double-star kwargs'' & \\
\texttt{None} & ``None'' (capital N) & Distinct from the word
``none'' \\
\texttt{True} / \texttt{False} & ``True'' / ``False'' (capital) & \\
\texttt{import} & ``import'' & \\
\texttt{from\ x\ import\ y} & ``from x import y'' & Read as English
sentence \\
\end{longtable}
}

\subsubsection{Reading a Line of Code
Aloud}\label{reading-a-line-of-code-aloud}

Teach students to read full lines this way, left to right, then explain:

\begin{Shaded}
\begin{Highlighting}[]
\NormalTok{total\_price }\OperatorTok{=}\NormalTok{ unit\_price }\OperatorTok{*}\NormalTok{ quantity }\OperatorTok{+}\NormalTok{ tax}
\end{Highlighting}
\end{Shaded}

\emph{``total underscore price gets unit underscore price times quantity
plus tax''}

Then explain: \emph{``This line calculates the total price by
multiplying the unit price by the quantity, then adding tax, and stores
the result in total\_price.''}

The two-pass approach --- first read the symbols, then explain the
meaning --- builds both syntax recognition and semantic understanding.

\begin{center}\rule{0.5\linewidth}{0.5pt}\end{center}

\subsection{5. Unit 0 --- Setup and
Environment}\label{unit-0-setup-and-environment}

\subsubsection{Lesson 0.1: Installing Python and VS
Code}\label{lesson-0.1-installing-python-and-vs-code}

\paragraph{Before This Lesson}\label{before-this-lesson}

Survey the student's current technology setup:

\begin{itemize}
\tightlist
\item
  What screen reader do they use, and what version?
\item
  Do they use a braille display?
\item
  What operating system?
\item
  Have they used a terminal before?
\item
  Do they have administrator access to their computer?
\end{itemize}

Confirm administrator access before scheduling this lesson. Installation
without admin rights will fail.

\paragraph{Instructional Cues}\label{instructional-cues}

Walk through every installer step verbally as the student performs it.
Do not say ``click the next button'' --- say ``Tab until you hear `Next
button,' then press Enter.'' Treat this as teaching screen reader
navigation, not just software installation.

\textbf{Critical checkpoint --- PATH configuration:} After installation,
before anything else: \texttt{python3\ -\/-version}. If this fails, stop
and address the PATH issue before proceeding. A student who continues
without Python on PATH will encounter mysterious failures throughout the
course.

\paragraph{Non-Obvious Accessibility
Strategies}\label{non-obvious-accessibility-strategies}

\begin{itemize}
\tightlist
\item
  \textbf{Windows Defender SmartScreen:} On Windows, running a
  downloaded \texttt{.exe} from Python.org may trigger a SmartScreen
  dialog. This dialog can be navigated with Tab and Enter but is not
  always announced immediately by NVDA. Tell the student what to expect
  before they run the installer.
\item
  \textbf{macOS Gatekeeper:} Downloaded packages may require right-click
  -\textgreater{} Open to bypass security. On a screen reader,
  right-click is VO+Shift+M (VoiceOver) or App Key (NVDA with keyboard
  access).
\item
  \textbf{Linux package managers as an alternative:} For Linux students,
  \texttt{sudo\ apt\ install\ python3\ python3-pip\ python3-venv} is
  often more reliable than the Python.org download and is fully
  scriptable.
\item
  \textbf{Portable Python for school computers:} Some schools restrict
  software installation. Python.org offers portable/embeddable
  distributions. Replit.com is a web-based Python IDE that requires no
  installation and is reasonably accessible.
\end{itemize}

\paragraph{Alternative Installation
Methods}\label{alternative-installation-methods}

{\def\LTcaptype{none} % do not increment counter
\begin{longtable}[]{@{}
  >{\raggedright\arraybackslash}p{(\linewidth - 4\tabcolsep) * \real{0.2115}}
  >{\raggedright\arraybackslash}p{(\linewidth - 4\tabcolsep) * \real{0.4327}}
  >{\raggedright\arraybackslash}p{(\linewidth - 4\tabcolsep) * \real{0.3558}}@{}}
\toprule\noalign{}
\begin{minipage}[b]{\linewidth}\raggedright
Method
\end{minipage} & \begin{minipage}[b]{\linewidth}\raggedright
Accessibility
\end{minipage} & \begin{minipage}[b]{\linewidth}\raggedright
Best For
\end{minipage} \\
\midrule\noalign{}
\endhead
\bottomrule\noalign{}
\endlastfoot
Python.org installer & High --- standard dialog flow & Most
installations \\
Chocolatey (Windows) & High --- command-line, screen-reader friendly &
Windows power users \\
Homebrew (macOS) & High --- command-line & macOS users comfortable with
terminal \\
apt/dnf/pacman (Linux) & High --- command-line & Linux users \\
Replit.com & Medium --- web-based, some accessibility gaps & No local
install possible \\
Thonny IDE & Medium --- simpler UI, bundled Python & Beginners who
struggle with VS Code \\
\end{longtable}
}

\textbf{On Thonny:} Thonny is a Python IDE specifically designed for
beginners. It has a simpler interface than VS Code and is sometimes more
accessible for students who are new to both screen readers and
programming. The trade-off is fewer advanced accessibility features
long-term. Consider Thonny for the first 2--3 lessons to reduce setup
burden, then migrate to VS Code.

\paragraph{Thought Experiments}\label{thought-experiments}

\begin{enumerate}
\def\labelenumi{\arabic{enumi}.}
\item
  \emph{``Why do you think the Python installer has a checkbox called
  `Add Python to PATH'? What is a PATH, and what would happen if Python
  weren't on it? Think about it like a phonebook --- your operating
  system looks up programs by name. If Python isn't in the phonebook,
  what happens when you type \texttt{python3}?''}
\item
  \emph{``We're installing two things: Python and VS Code. What role
  does each one play? If you had only Python and no VS Code, could you
  still write programs? How? What would you lose?''}
\item
  \emph{``Virtual environments keep each project's dependencies
  separate. Imagine two school projects --- one requiring a specific
  edition of a textbook that contradicts the edition another class uses.
  What would happen if you tried to use both editions at the same time?
  How does a virtual environment solve the programming equivalent of
  this problem?''}
\item
  \emph{``Why do we verify every installation immediately after
  completing it? What could go wrong between downloading something and
  confirming it works? Name three things that could silently fail.''}
\end{enumerate}

\begin{center}\rule{0.5\linewidth}{0.5pt}\end{center}

\subsubsection{Lesson 0.2: Configuring Your Screen Reader for VS
Code}\label{lesson-0.2-configuring-your-screen-reader-for-vs-code}

\paragraph{Instructional Cues}\label{instructional-cues-1}

Before showing any configuration screens, do an unguided exploration:
``Open VS Code, navigate around with only your keyboard and screen
reader, and come back in ten minutes. Tell me what you found, what
worked, what was confusing.'' This builds genuine familiarity and
identifies that student's specific challenges.

Map each configuration step to a problem the student already identified.
This makes the configuration feel purposeful rather than arbitrary.

\paragraph{Screen Reader Specific
Guidance}\label{screen-reader-specific-guidance}

\textbf{NVDA users:}

\begin{itemize}
\tightlist
\item
  Enable ``Use UI Automation to access VS Code'' in NVDA preferences
  (Tools -\textgreater{} Manage Plugins, or via NVDA+N -\textgreater{}
  Preferences -\textgreater{} Settings -\textgreater{} Advanced
  -\textgreater{} Use UI Automation where available). This dramatically
  improves reading accuracy in VS Code.
\item
  Set ``Report dynamic content changes'' to ``All'' to hear error
  highlights and IntelliSense popups.
\item
  Configure Browse Mode to be OFF by default in VS Code --- programming
  requires forms mode (interactive mode) at all times.
\item
  Speech rate: many experienced NVDA users run at 70--90\% of maximum
  speed. For new students, 40--60\% allows better comprehension of code
  read aloud.
\end{itemize}

\textbf{JAWS users:}

\begin{itemize}
\tightlist
\item
  JAWS has VS Code-specific scripts maintained by the JAWS community.
  Check \texttt{FreedomScientific.com} and GitHub for the latest VS Code
  JAWS scripts.
\item
  PlaceMarkers in JAWS can mark frequently visited locations in a file
  --- teach \texttt{Ctrl+Win+K} to set and \texttt{Ctrl+Win+J} to jump.
\item
  JAWS Verbosity settings: set ``Announce line numbers'' to ON for code
  files.
\end{itemize}

\textbf{VoiceOver (macOS) users:}

\begin{itemize}
\tightlist
\item
  VoiceOver interacts with VS Code via the macOS Accessibility API.
  Ensure VS Code is not run with elevated permissions (which can break
  accessibility API access).
\item
  Use \texttt{VO+Shift+Down} to enter interactive elements (the editor,
  terminal), then navigate normally.
\item
  Enable ``Speak text under mouse'' in VoiceOver Utility -\textgreater{}
  Verbosity if the student uses residual vision with a mouse.
\item
  VoiceOver's ``item chooser'' (VO+I) can provide a quick list of
  interactive elements in the current window.
\end{itemize}

\textbf{Windows Narrator users:}

\begin{itemize}
\tightlist
\item
  Narrator is less feature-rich than NVDA or JAWS for programming. It
  works but may require more keystrokes.
\item
  Narrator scan mode must be OFF in VS Code for normal keyboard input to
  work. Toggle with Caps Lock+Space.
\item
  Consider recommending the free NVDA screen reader as an upgrade for
  students using Narrator, as NVDA's VS Code support is significantly
  better.
\end{itemize}

\textbf{Orca (Linux) users:}

\begin{itemize}
\tightlist
\item
  Orca works well with VS Code on Linux via the AT-SPI2 accessibility
  interface.
\item
  Use \texttt{Orca+H} for ``say all'' from cursor position.
\item
  Ensure the \texttt{at-spi2-core} package is installed:
  \texttt{sudo\ apt\ install\ at-spi2-core}.
\item
  Orca's ``flat review'' mode (\texttt{Orca+KP\_Enter} or
  \texttt{Orca+Insert+KP\_Enter}) is useful for reading lines that are
  not focus-accessible.
\end{itemize}

\paragraph{Indentation Configuration --- Step by
Step}\label{indentation-configuration-step-by-step}

This is the most important configuration in the entire course. Walk
through it during the first session and verify it is working before
writing any Python.

\textbf{Test procedure:}

\begin{enumerate}
\def\labelenumi{\arabic{enumi}.}
\item
  Create a temporary file \texttt{test\_indent.py}.
\item
  Write:

\begin{Shaded}
\begin{Highlighting}[]
\ControlFlowTok{if} \VariableTok{True}\NormalTok{:}
    \BuiltInTok{print}\NormalTok{(}\StringTok{"indented"}\NormalTok{)}
\BuiltInTok{print}\NormalTok{(}\StringTok{"not indented"}\NormalTok{)}
\end{Highlighting}
\end{Shaded}
\item
  Navigate line by line with Down arrow.
\item
  The screen reader should announce the indent level on line 2 before
  reading its content.
\item
  On line 3, the indent should return to zero.
\end{enumerate}

If announcement is not happening, troubleshoot screen reader settings
before proceeding.

\paragraph{Thought Experiments}\label{thought-experiments-1}

\begin{enumerate}
\def\labelenumi{\arabic{enumi}.}
\item
  \emph{``A sighted programmer can look at a Python file and immediately
  understand its structure from indentation --- like an outline. How
  does your screen reader translate that same 2D visual structure into a
  1D audio stream? What information might be lost in that
  translation?''}
\item
  \emph{``VS Code has hundreds of keyboard shortcuts. No sighted
  programmer uses all of them --- they use the 15 or 20 that matter for
  their workflow. If you could only keep 10 shortcuts, which 10 would
  make the biggest difference for writing Python? Why?''}
\item
  \emph{``The Accessible View (Shift+Alt+F2) converts complex visual
  elements into flat text. What kind of information in VS Code might
  require this conversion? Can you think of a scenario where the
  `visual' presentation of something would be unusable without
  translation?''}
\end{enumerate}

\begin{center}\rule{0.5\linewidth}{0.5pt}\end{center}

\subsubsection{Lesson 0.3: Setting Up Your VS Code
Workspace}\label{lesson-0.3-setting-up-your-vs-code-workspace}

\paragraph{Alternative Workspace
Tools}\label{alternative-workspace-tools}

If VS Code proves too complex initially, consider these alternatives as
temporary or permanent solutions:

{\def\LTcaptype{none} % do not increment counter
\begin{longtable}[]{@{}
  >{\raggedright\arraybackslash}p{(\linewidth - 4\tabcolsep) * \real{0.2263}}
  >{\raggedright\arraybackslash}p{(\linewidth - 4\tabcolsep) * \real{0.0584}}
  >{\raggedright\arraybackslash}p{(\linewidth - 4\tabcolsep) * \real{0.7153}}@{}}
\toprule\noalign{}
\begin{minipage}[b]{\linewidth}\raggedright
Tool
\end{minipage} & \begin{minipage}[b]{\linewidth}\raggedright
Platform
\end{minipage} & \begin{minipage}[b]{\linewidth}\raggedright
Accessibility Notes
\end{minipage} \\
\midrule\noalign{}
\endhead
\bottomrule\noalign{}
\endlastfoot
\textbf{IDLE} (bundled with Python) & All & Very simple UI. Limited
features but highly keyboard accessible. Good for first programs. \\
\textbf{Thonny} & All & Beginner-focused. Has a ``simple mode.'' Some
screen reader users find it more approachable. \\
\textbf{Mu Editor} & All & Designed for education. Simpler than VS Code.
Python 3 only. \\
\textbf{Jupyter Notebook (terminal)} & All & Text-based notebook.
Accessible for structured exploration. \\
\textbf{Braille Note Apex/Touch} & Device & For students using these
notetakers, Python can be run via their built-in terminals. \\
\textbf{Replit} & Web & No install required. Accessibility varies; test
before recommending. \\
\textbf{NVDA Remote + VS Code} & Windows & Student uses their own screen
reader on their own machine while teacher controls a shared session. \\
\end{longtable}
}

\paragraph{Thought Experiments}\label{thought-experiments-2}

\begin{enumerate}
\def\labelenumi{\arabic{enumi}.}
\item
  \emph{``Programmers organize their code into folders by project. What
  happens if every programmer at a company kept all their files in one
  shared folder with no organization? How does this analogy apply to
  organizing your own Python projects?''}
\item
  \emph{``VS Code Extensions add features to the editor the same way
  apps add features to a phone. If you could design one VS Code
  extension specifically for blind Python programmers that doesn't exist
  yet, what would it do?''}
\end{enumerate}

\begin{center}\rule{0.5\linewidth}{0.5pt}\end{center}

\subsubsection{Lesson 0.4: Python Environments and Terminal
Usage}\label{lesson-0.4-python-environments-and-terminal-usage}

\paragraph{Terminal Proficiency as a
Prerequisite}\label{terminal-proficiency-as-a-prerequisite}

Before writing any Python, ensure the student can fluently:

\begin{itemize}
\tightlist
\item
  Navigate to a directory with \texttt{cd}
\item
  List files with \texttt{ls} (macOS/Linux) or \texttt{dir} (Windows)
\item
  Create a directory with \texttt{mkdir}
\item
  Print the current directory with \texttt{pwd} (macOS/Linux) or
  \texttt{cd} with no arguments (Windows)
\item
  Use Tab completion for filenames
\item
  Use Up/Down arrow for command history
\item
  Use \texttt{Ctrl+C} to stop a running program
\end{itemize}

If these are not automatic, spend additional time on terminal basics.
Terminal fluency is foundational to everything else.

\paragraph{Alternative Terminals for
Accessibility}\label{alternative-terminals-for-accessibility}

{\def\LTcaptype{none} % do not increment counter
\begin{longtable}[]{@{}
  >{\raggedright\arraybackslash}p{(\linewidth - 4\tabcolsep) * \real{0.2597}}
  >{\raggedright\arraybackslash}p{(\linewidth - 4\tabcolsep) * \real{0.1039}}
  >{\raggedright\arraybackslash}p{(\linewidth - 4\tabcolsep) * \real{0.6364}}@{}}
\toprule\noalign{}
\begin{minipage}[b]{\linewidth}\raggedright
Terminal
\end{minipage} & \begin{minipage}[b]{\linewidth}\raggedright
Platform
\end{minipage} & \begin{minipage}[b]{\linewidth}\raggedright
Screen Reader Compatibility
\end{minipage} \\
\midrule\noalign{}
\endhead
\bottomrule\noalign{}
\endlastfoot
\textbf{Windows Terminal} & Windows & Excellent with NVDA/JAWS (use UI
Automation mode) \\
\textbf{cmd.exe} & Windows & Good; legacy but very reliable \\
\textbf{PowerShell} & Windows & Good with NVDA; some JAWS quirks \\
\textbf{macOS Terminal} & macOS & Excellent with VoiceOver \\
\textbf{iTerm2} & macOS & Good; more features, slightly more complex \\
\textbf{GNOME Terminal} & Linux & Good with Orca \\
\textbf{Konsole} & Linux & Good with Orca \\
\textbf{Terminator} & Linux & Good; split-pane support \\
\end{longtable}
}

\paragraph{Thought Experiments}\label{thought-experiments-3}

\begin{enumerate}
\def\labelenumi{\arabic{enumi}.}
\item
  \emph{``Virtual environments isolate each project's dependencies. Why
  would this matter for a professional programmer with 20 active
  projects? What could go wrong without them?''}
\item
  \emph{``When you type \texttt{python3\ hello.py} and press Enter, what
  is the operating system doing? What is Python doing? Walk through it
  step by step as if explaining to someone who has never used a
  computer.''}
\item
  \emph{``\texttt{pip\ install\ requests} downloads and installs code
  written by other people. What are the risks of running code written by
  strangers? What does PyPI (the Python Package Index) do to reduce
  these risks?''}
\end{enumerate}

\begin{center}\rule{0.5\linewidth}{0.5pt}\end{center}

\subsection{6. Unit 1 --- Foundations and Getting
Started}\label{unit-1-foundations-and-getting-started}

\subsubsection{Lesson 1.1: Welcome to Computer
Science}\label{lesson-1.1-welcome-to-computer-science}

\paragraph{Instructional Cues}\label{instructional-cues-2}

Open with \textbf{the ``Human Computer'' activity} --- this is the most
valuable exercise in the entire unit for blind students. The student
gives verbal, step-by-step instructions for a simple physical task. The
teacher follows them literally, making no inferences. This demonstrates
the core challenge of programming: computers are perfectly literal and
do not fill in gaps.

Suggest tasks that connect to the student's daily experience:

\begin{itemize}
\tightlist
\item
  Instructions for making tea or coffee
\item
  Instructions for navigating from the classroom door to their desk
\item
  Instructions for finding a contact in their phone using a screen
  reader
\end{itemize}

After the activity, debrief: \emph{``Where did your instructions assume
I could see? Where did they assume I already knew something? A computer
makes neither assumption.''}

\paragraph{Connecting Computing to Accessibility
Technology}\label{connecting-computing-to-accessibility-technology}

Begin with what the student already knows:

\begin{itemize}
\tightlist
\item
  Their screen reader is software. It runs on hardware. It was written
  by programmers. It converts visual information (text rendering) into
  audio output.
\item
  Ask: ``What are the inputs to your screen reader? What are the
  outputs? What would happen if a website used images of text instead of
  real text --- how would that change what your screen reader can do?''
\item
  Connect to ARIA attributes, alt text, and semantic HTML as programming
  concepts that directly affect the student's daily experience.
\end{itemize}

\paragraph{Thought Experiments}\label{thought-experiments-4}

\begin{enumerate}
\def\labelenumi{\arabic{enumi}.}
\item
  \emph{``Computer science is described as the study of problem-solving,
  not just programming. What non-programming problems have you solved
  using computational thinking --- breaking a problem down into steps,
  looking for patterns, and abstracting the details? Describe one
  example from daily life.''}
\item
  \emph{``Your screen reader converts text to audio in real time. What
  `algorithm' do you think it uses? What decisions does it make --- when
  to pause, how to handle punctuation, what to do with symbols? What
  would a flowchart of those decisions look like?''}
\item
  \emph{``Early computers could only run one program at a time. Modern
  computers run hundreds simultaneously. Without changing the physical
  hardware, what concept makes this possible? (Hint: the answer is
  software --- specifically, an operating system.)''}
\item
  \emph{``Software has often been described as `making the hardware do
  something useful.' But hardware without software is just metal and
  silicon. In what sense is software more important than hardware? In
  what sense is hardware more important?''}
\item
  \emph{``Accessibility technology like screen readers, braille
  displays, and audio description didn't exist when computers were first
  invented. Who decided these should be built? What does this history
  tell you about who gets to shape what technology does?''}
\end{enumerate}

\paragraph{Additional Resources for This
Lesson}\label{additional-resources-for-this-lesson}

\begin{itemize}
\tightlist
\item
  \textbf{``The Computer Boys Take Over'' by Nathan Ensmenger} ---
  accessible history of how programming became a profession.
\item
  \textbf{``Technically Wrong'' by Sara Wachter-Boettcher} --- explores
  how biases in tech affect users, including disabled users. Good for
  older students.
\item
  \textbf{American Foundation for the Blind (AFB) CareerConnect} ---
  profiles of blind professionals in STEM fields, including software
  engineers. Good for motivation and role models.
\item
  \textbf{``Haben: The Deafblind Woman Who Conquered Harvard Law'' by
  Haben Girma} --- autobiography of a deafblind lawyer and disability
  rights advocate. Haben's work directly involves technology access.
\item
  \textbf{Code.org's ``How Computers Work'' video series} ---
  audio-described versions available. Covers hardware and software
  basics accessibly.
\end{itemize}

\begin{center}\rule{0.5\linewidth}{0.5pt}\end{center}

\subsubsection{Lesson 1.2: Your First Python
Program}\label{lesson-1.2-your-first-python-program}

\paragraph{Instructional Cues}\label{instructional-cues-3}

\textbf{Before any code:} Ask the student to predict. ``If I write
\texttt{print("Hello,\ World!")} and run it, what do you think will
appear in the terminal? Where exactly will it appear?'' Let them answer,
then verify together.

\textbf{Make the first error deliberate.} Remove a closing parenthesis,
run the program, and navigate to the error message together. Establish
the error-reading protocol:

\begin{enumerate}
\def\labelenumi{\arabic{enumi}.}
\tightlist
\item
  Terminal output appears after the code.
\item
  Read from the bottom of the error upward --- the bottom is the most
  specific message.
\item
  Line number tells you where to go. Use \texttt{Ctrl+G}.
\item
  Error type (\texttt{SyntaxError}, \texttt{NameError}) narrows the
  category of problem.
\end{enumerate}

\textbf{The REPL as the first environment.} Before creating any files,
start with \texttt{python3} in the terminal and type
\texttt{print("Hello,\ World!")} directly. The immediate audio feedback
of the REPL establishes cause and effect more directly than the file
-\textgreater{} run -\textgreater{} find output workflow.

\paragraph{Common First Errors and How to Explain
Them}\label{common-first-errors-and-how-to-explain-them}

{\def\LTcaptype{none} % do not increment counter
\begin{longtable}[]{@{}
  >{\raggedright\arraybackslash}p{(\linewidth - 4\tabcolsep) * \real{0.2038}}
  >{\raggedright\arraybackslash}p{(\linewidth - 4\tabcolsep) * \real{0.1754}}
  >{\raggedright\arraybackslash}p{(\linewidth - 4\tabcolsep) * \real{0.6209}}@{}}
\toprule\noalign{}
\begin{minipage}[b]{\linewidth}\raggedright
Error
\end{minipage} & \begin{minipage}[b]{\linewidth}\raggedright
Likely Cause
\end{minipage} & \begin{minipage}[b]{\linewidth}\raggedright
Explanation for Student
\end{minipage} \\
\midrule\noalign{}
\endhead
\bottomrule\noalign{}
\endlastfoot
\texttt{SyntaxError:\ unexpected\ EOF\ while\ parsing} & Missing closing
\texttt{"}, \texttt{)}, or \texttt{{]}} & ``Python reached the end of
the file before it expected to. It was waiting for something ---
probably a closing character.'' \\
\texttt{NameError:\ name\ \textquotesingle{}print\textquotesingle{}\ is\ not\ defined}
& Capital P in \texttt{Print} & ``Python is case-sensitive.
\texttt{Print} and \texttt{print} are completely different names. Only
\texttt{print} (lowercase) is the built-in function.'' \\
\texttt{IndentationError} & Extra space before first line & ``Python's
first rule: top-level code cannot be indented. Every line that is not
inside a function or class must start at column 0.'' \\
\texttt{SyntaxError:\ invalid\ syntax} at \texttt{:} & Using = instead
of == in a comparison & ``The assignment operator = doesn't belong
inside an if condition. For comparison, you need the double equals
==.'' \\
\end{longtable}
}

\paragraph{Thought Experiments}\label{thought-experiments-5}

\begin{enumerate}
\def\labelenumi{\arabic{enumi}.}
\item
  \emph{``\texttt{print()} produces output in the terminal. What would a
  program be like with NO output at all? Could such a program still do
  something useful? (Hint: think about programs that read files, send
  network requests, or modify a database.)''}
\item
  \emph{``Python is an interpreted language --- each line is read and
  executed one at a time. Other languages are compiled --- the entire
  program is translated before any of it runs. How does this affect the
  kind of errors you encounter, and when you encounter them?''}
\item
  \emph{``The \texttt{\#} symbol starts a comment --- text that Python
  ignores entirely. Why do comments exist? If the computer ignores them,
  who are they for? Could you write good programs with no comments?
  Should you?''}
\item
  \emph{``\,`Hello, World!' is the traditional first program in almost
  every programming language. Why do you think this tradition exists?
  What does this simple program prove about a working environment?''}
\end{enumerate}

\paragraph{Alternative Resources for This
Lesson}\label{alternative-resources-for-this-lesson}

\begin{itemize}
\tightlist
\item
  \textbf{``Invent Your Own Computer Games with Python'' by Al Sweigart
  (free online):} Text-based games that produce only terminal output ---
  perfectly accessible. Chapter 1 is a good companion to this lesson.
\item
  \textbf{CS Unplugged --- ``Parity Magic'' activity:} Teaches binary
  representation and algorithms without a computer. Can precede or
  follow this lesson.
\item
  \textbf{Python REPL via iOS/Android with screen reader:} Several apps
  provide a Python REPL on mobile. For students whose phone screen
  reader is more familiar than their computer setup, this may be a
  comfortable way to experiment.
\end{itemize}

\begin{center}\rule{0.5\linewidth}{0.5pt}\end{center}

\subsubsection{Lesson 1.3: Variables and Simple Data
Types}\label{lesson-1.3-variables-and-simple-data-types}

\paragraph{Instructional Cues}\label{instructional-cues-4}

\textbf{The labeled box metaphor} works well across modalities ---
introduce it verbally and reinforce with tactile analogy if helpful.
Provide a physical box (a small container) with a label stuck to it and
a slip of paper inside representing the value. The box remains; the
paper changes. The label stays the same; the contents change with
reassignment.

\textbf{\texttt{type()} as a diagnostic reflex:} Begin every data type
lesson with \texttt{type()}. \texttt{print(type(42))} -\textgreater{}
\texttt{int}. \texttt{print(type("hello"))} -\textgreater{}
\texttt{str}. Make this a reflex the student uses whenever they're
unsure what kind of data they have. For a student who cannot visually
scan code for context clues, \texttt{type()} is often the only reliable
way to confirm what a variable contains.

\textbf{The string/number audio confusion:} When a screen reader reads
\texttt{"42"} and \texttt{42}, they sound identical. This is a critical
accessibility issue. Address it directly:

\begin{itemize}
\tightlist
\item
  \emph{``Your screen reader will read these the same way. The only
  reliable way to tell them apart is \texttt{type()}. Use it.
  Constantly.''}
\item
  Demonstrate the consequence: \texttt{"42"\ +\ 1} raises
  \texttt{TypeError}. \texttt{42\ +\ 1} gives \texttt{43}.
\end{itemize}

\paragraph{Multiple Explanations for Variable Naming
Rules}\label{multiple-explanations-for-variable-naming-rules}

Provide all of these; let the student choose what sticks:

\begin{itemize}
\tightlist
\item
  \emph{Rule-based:} Must start with letter or underscore. May contain
  letters, digits, underscores. May not be a Python keyword.
\item
  \emph{Analogy:} ``Variable names are like legal names. They can't
  start with a number (no one's legal name is `2Fast'), can't contain
  hyphens (the name Mc-Donald would confuse paperwork), and can't be
  reserved words like `True' or `if' that already mean something
  official.''
\item
  \emph{Mnemonic for invalid starts:} ``Numbers first? Never. Hyphens?
  Horrible. Spaces? Surely not.''
\end{itemize}

\paragraph{Thought Experiments}\label{thought-experiments-6}

\begin{enumerate}
\def\labelenumi{\arabic{enumi}.}
\item
  \emph{``If variables are labeled boxes, what happens when you do
  \texttt{x\ =\ 5} and then \texttt{y\ =\ x} and then \texttt{x\ =\ 10}?
  What is \texttt{y} now? Does reassigning \texttt{x} affect \texttt{y}?
  Why or why not?''}
\item
  \emph{``Python is `dynamically typed,' meaning a variable can change
  its type at any time: \texttt{x\ =\ 5} then
  \texttt{x\ =\ \textquotesingle{}hello\textquotesingle{}}. Many
  languages forbid this. What bugs could dynamic typing accidentally
  create? What flexibility does it offer?''}
\item
  \emph{``The integer \texttt{42} and the string
  \texttt{\textquotesingle{}42\textquotesingle{}} look identical when
  spoken aloud. How many other pairs of values can you think of that
  might be ambiguous when a screen reader reads them? What coding habits
  help prevent confusion?''}
\item
  \emph{``Constants in programming represent values that never change
  --- like \texttt{SPEED\_OF\_LIGHT} or \texttt{MAX\_USERS}. Python has
  no true constants (you can always reassign them). Why do programmers
  use \texttt{CAPS\_SNAKE\_CASE} as a convention for constants even
  though Python can't enforce them?''}
\end{enumerate}

\paragraph{Alternative Resources}\label{alternative-resources}

\begin{itemize}
\tightlist
\item
  \textbf{CS50 Python (Harvard, edX):} David Malan's accessible video
  lectures. Closed-captions and transcripts available for all lessons.
  Covers variables and types in Week 0.
\item
  \textbf{Python.org official tutorial:} Plain HTML, navigable by screen
  reader. Section 3 covers numbers, strings, and basic types.
\item
  \textbf{``Python for Everybody'' by Dr.~Chuck Severance (free
  online):} Audio-described video lectures available. Chapter 2 covers
  variables and expressions.
\item
  \textbf{Codecademy ``Learn Python 3'':} Web-based coding environment.
  Accessibility is partial --- code editor is accessible but some UI
  elements are not. Worth testing with the student's specific setup.
\end{itemize}

\begin{center}\rule{0.5\linewidth}{0.5pt}\end{center}

\subsection{7. Unit 2 --- Numeric Computations and
I/O}\label{unit-2-numeric-computations-and-io}

\subsubsection{Lesson 2.1: The Software Development
Process}\label{lesson-2.1-the-software-development-process}

\paragraph{Instructional Cues}\label{instructional-cues-5}

This is a conceptual lesson. Resist the urge to rush to code. Spend the
full session on design and pseudocode.

\textbf{The ``Five Whys'' technique:} Before writing pseudocode, ask the
student to state the problem in one sentence. Then ask ``why?'' five
times to reach the root of what the program really needs to do. This
technique from Lean manufacturing translates well to programming.

Example: ``My program calculates a restaurant tip.''

\begin{itemize}
\tightlist
\item
  Why? ``So the user knows how much to pay.''
\item
  Why do they need to calculate it? ``The tip amount isn't printed on
  the bill.''
\item
  Why? ``The restaurant leaves the amount to the customer's
  discretion.''
\item
  So the real problem is: ``Given a bill amount, calculate 15\%, 18\%,
  and 20\% tip options, plus the total.''
\end{itemize}

This more complete problem statement produces a better program design.

\textbf{Pseudocode format for screen reader users:} Plain \texttt{.txt}
files work best for pseudocode. Avoid word processors or special
formatting tools that add accessibility overhead. The student should be
able to read their pseudocode fluently while writing their Python.

\paragraph{Thought Experiments}\label{thought-experiments-7}

\begin{enumerate}
\def\labelenumi{\arabic{enumi}.}
\item
  \emph{``In engineering, a bridge is designed many times before a
  single beam is placed. The cost of redesigning a blueprint is near
  zero; the cost of redesigning a half-built bridge is enormous. How
  does this ratio (cheap to redesign early, expensive to redesign later)
  apply to software?''}
\item
  \emph{``The design phase involves thinking about edge cases ---
  unusual inputs. For a program that calculates how many buses are
  needed to transport N students, with each bus holding 40 students,
  what edge cases exist? (Consider: What if N=0? What if N=40? What if
  N=41? What if N is negative?)''}
\item
  \emph{``Some agile software development methodologies say `working
  software over comprehensive documentation.' Does this conflict with
  the value of design? How do you reconcile the two?''}
\item
  \emph{``Testing is described as part of the development process. But
  some teams test after development is complete. What is lost by testing
  at the end rather than continuously throughout? What is gained?''}
\end{enumerate}

\begin{center}\rule{0.5\linewidth}{0.5pt}\end{center}

\subsubsection{Lessons 2.2--2.5: Expressions, Input, Types,
Math}\label{lessons-2.22.5-expressions-input-types-math}

\paragraph{Instructional Cues}\label{instructional-cues-6}

\textbf{\texttt{input()} always returns a string --- make it tangible:}

\begin{Shaded}
\begin{Highlighting}[]
\NormalTok{user\_input }\OperatorTok{=} \BuiltInTok{input}\NormalTok{(}\StringTok{"Type the number 5: "}\NormalTok{)}
\BuiltInTok{print}\NormalTok{(}\BuiltInTok{type}\NormalTok{(user\_input))   }\CommentTok{\# \textless{}class \textquotesingle{}str\textquotesingle{}\textgreater{}}
\BuiltInTok{print}\NormalTok{(user\_input }\OperatorTok{+} \DecValTok{1}\NormalTok{)      }\CommentTok{\# TypeError}
\BuiltInTok{print}\NormalTok{(}\BuiltInTok{int}\NormalTok{(user\_input) }\OperatorTok{+} \DecValTok{1}\NormalTok{) }\CommentTok{\# 6}
\end{Highlighting}
\end{Shaded}

Run this sequence and let the student hear both the error and the fix.
The \texttt{TypeError} message should be read aloud and analyzed before
the fix is shown.

\textbf{Order of operations --- the hearing test:} Read
\texttt{10\ +\ 5\ *\ 2} aloud and ask for the answer before Python is
consulted. Most students say 30; Python says 20. This is a memorable
demonstration of why precise symbol reading matters. ``Your ear parsed
it as `add 10 and 5, then multiply by 2.' Python parses it as `10 plus
the result of 5 times 2.'\,''

\textbf{\texttt{math} module exploration via REPL:}

\begin{Shaded}
\begin{Highlighting}[]
\ImportTok{import}\NormalTok{ math}
\BuiltInTok{dir}\NormalTok{(math)  }\CommentTok{\# Lists all functions in the module}
\BuiltInTok{help}\NormalTok{(math.sqrt)  }\CommentTok{\# Reads the docstring aloud}
\end{Highlighting}
\end{Shaded}

\texttt{dir()} and \texttt{help()} are screen-reader-friendly tools for
exploring any module. Teach them as first-resort documentation tools.

\paragraph{Thought Experiments}\label{thought-experiments-8}

\begin{enumerate}
\def\labelenumi{\arabic{enumi}.}
\item
  \emph{``Type conversion functions like \texttt{int()} and
  \texttt{float()} can fail:
  \texttt{int(\textquotesingle{}hello\textquotesingle{})} raises a
  \texttt{ValueError}. Who is responsible --- the programmer who called
  \texttt{int()} on untrusted input, or the user who typed `hello'? How
  should a well-designed program handle this?''}
\item
  \emph{``\texttt{math.sqrt(25)} returns \texttt{5.0}, not \texttt{5}.
  Why does a function that logically produces a whole number return a
  float? What would break if it returned an integer instead?''}
\item
  \emph{``Floating-point numbers in computers are approximations ---
  \texttt{0.1\ +\ 0.2} gives \texttt{0.30000000000000004}. Describe a
  real-world application where this imprecision could cause serious
  problems. What strategies exist to avoid it?''}
\item
  \emph{``The \texttt{//} operator performs floor division:
  \texttt{7\ //\ 2} = \texttt{3}. The \texttt{\%} operator gives the
  remainder: \texttt{7\ \%\ 2} = \texttt{1}. Together they completely
  describe integer division. What real-world problem does each solve?
  Describe a program that would need both.''}
\end{enumerate}

\paragraph{Alternative Resources}\label{alternative-resources-1}

\begin{itemize}
\tightlist
\item
  \textbf{``Automate the Boring Stuff with Python'' Chapter 1--3 (free
  online):} Covers expressions, data types, and input with practical
  exercises. Plain HTML, accessible.
\item
  \textbf{Khan Academy ``Introduction to Programming'':} Video series
  with transcripts. Good for visual-thinker students who benefit from
  hearing concepts explained differently.
\item
  \textbf{SoloLearn Python course:} Mobile app with accessible
  exercises. Works well with iOS VoiceOver and Android TalkBack for
  students who prefer mobile.
\item
  \textbf{``Python Crash Course'' by Eric Matthes:} The primary textbook
  for this curriculum. Chapter 2 aligns directly with this unit.
\end{itemize}

\begin{center}\rule{0.5\linewidth}{0.5pt}\end{center}

\subsection{8. Unit 3 --- Control Flow, Decisions, and
Loops}\label{unit-3-control-flow-decisions-and-loops}

\subsubsection{Lessons 3.1--3.2:
Conditionals}\label{lessons-3.13.2-conditionals}

\paragraph{Instructional Cues}\label{instructional-cues-7}

\textbf{Verbal simulation protocol:} Before running any conditional
code, have the student perform the simulation verbally:

\begin{enumerate}
\def\labelenumi{\arabic{enumi}.}
\tightlist
\item
  ``What is the value of \texttt{score}?''
\item
  ``Evaluate the condition: is \texttt{85\ \textgreater{}=\ 90}? That's
  False.''
\item
  ``Because the condition is False, the if block is skipped.''
\item
  ``Evaluate the next condition: is \texttt{85\ \textgreater{}=\ 80}?
  That's True.''
\item
  ``The elif block runs. We print `Grade: B'.''
\end{enumerate}

This 5-step protocol should become automatic before the student runs any
code. It builds the mental simulation ability that is the core of
debugging.

\textbf{The assignment-vs-comparison trap:}

\begin{Shaded}
\begin{Highlighting}[]
\CommentTok{\# This is a bug that is SILENT in some languages}
\ControlFlowTok{if}\NormalTok{ score }\OperatorTok{=} \DecValTok{90}\NormalTok{:   }\CommentTok{\# SyntaxError in Python (good!)}
    \BuiltInTok{print}\NormalTok{(}\StringTok{"Perfect!"}\NormalTok{)}
\end{Highlighting}
\end{Shaded}

Python raises a \texttt{SyntaxError} here, which is actually a safety
feature. Other languages (C, JavaScript) allow this and it becomes a
silent logic bug. Appreciate Python's strictness here.

\paragraph{Thought Experiments}\label{thought-experiments-9}

\begin{enumerate}
\def\labelenumi{\arabic{enumi}.}
\item
  \emph{``If Python didn't have \texttt{elif}, only \texttt{if} and
  \texttt{else}, could you still write a grading program? Show how. What
  is the cost of doing it without \texttt{elif}?''}
\item
  \emph{``The \texttt{else} clause at the end of an if-elif chain is a
  catch-all --- it runs if nothing else matched. What is the programmer
  asserting when they write
  \texttt{else:\ print(\textquotesingle{}Invalid\ input\textquotesingle{})}?
  What assumptions about the universe of possible inputs does this make?
  When might those assumptions be wrong?''}
\item
  \emph{``Boolean logic (\texttt{and}, \texttt{or}, \texttt{not}) comes
  from mathematics (George Boole, 1847). Boolean logic also describes
  how circuits in a computer work. In what sense is every program you
  write built on top of millions of tiny \texttt{and}, \texttt{or}, and
  \texttt{not} operations happening in hardware?''}
\item
  \emph{``Short-circuit evaluation: in \texttt{if\ A\ and\ B}, Python
  doesn't evaluate \texttt{B} if \texttt{A} is False (because the result
  is already False). In \texttt{if\ A\ or\ B}, Python doesn't evaluate
  \texttt{B} if \texttt{A} is True. When is this behavior a useful
  feature rather than just an optimization?''}
\item
  \emph{``Consider a program that decides whether to grant access based
  on two conditions: correct password AND recognized device. Map this to
  Python \texttt{and} logic. Now consider: unlock if either condition
  passes (maybe for emergency access). Map this to \texttt{or}. What are
  the security implications of each design?''}
\end{enumerate}

\paragraph{Alternative Resources}\label{alternative-resources-2}

\begin{itemize}
\tightlist
\item
  \textbf{CS Unplugged ``Binary Numbers'' activity:} Physical,
  non-computer activity that illustrates Boolean thinking. Can be
  adapted for tactile use.
\item
  \textbf{``Think Python'' by Allen B. Downey (free online):} Chapter 5
  covers conditionals with clear explanation and exercises. HTML version
  is screen-reader-friendly.
\item
  \textbf{Parsons Problems (online):} Code ordering exercises --- the
  student arranges given lines into correct order. Does not require
  typing code from scratch. Available at parsons.problemsolving.org
  (accessibility varies; test first).
\end{itemize}

\begin{center}\rule{0.5\linewidth}{0.5pt}\end{center}

\subsubsection{Lessons 3.3--3.4: Lists and For
Loops}\label{lessons-3.33.4-lists-and-for-loops}

\paragraph{Instructional Cues}\label{instructional-cues-8}

\textbf{Mental models for lists --- choose one per student:}

\begin{itemize}
\tightlist
\item
  \emph{The train:} Each car has a number starting at 0. Car 0 is at the
  front. You can add cars to the end, remove from anywhere, and swap
  contents.
\item
  \emph{The playlist:} Songs in order, starting at position 0. You can
  shuffle, sort, add, or remove tracks.
\item
  \emph{The physical stack of cards:} A sequence of items that can be
  addressed by position.
\end{itemize}

\textbf{The for loop as a bookmark reader:} \emph{``Imagine Python has a
finger that moves along a list. It starts at position 0. It reads what's
there, runs the loop body, moves the finger to position 1, reads, runs
the body, and continues until the finger falls off the end.''}

\textbf{Tracing loops audibly --- required exercise:} For the first five
loops students write, require them to verbally trace every iteration
before running the code:

\begin{itemize}
\tightlist
\item
  ``Iteration 1: \texttt{fruit} gets `apple'. We print `I like
  apple.'\,''
\item
  ``Iteration 2: \texttt{fruit} gets `banana'. We print `I like
  banana.'\,''
\item
  etc.
\end{itemize}

\paragraph{Non-Obvious Accessibility
Strategies}\label{non-obvious-accessibility-strategies-1}

\begin{itemize}
\item
  \textbf{Printing list contents accessibly:} A bare
  \texttt{print(my\_list)} produces
  \texttt{{[}\textquotesingle{}apple\textquotesingle{},\ \textquotesingle{}banana\textquotesingle{},\ \textquotesingle{}cherry\textquotesingle{}{]}}
  with brackets and quotes. This is cluttered when read aloud. Teach the
  clean print pattern early:

\begin{Shaded}
\begin{Highlighting}[]
\ControlFlowTok{for}\NormalTok{ i, item }\KeywordTok{in} \BuiltInTok{enumerate}\NormalTok{(my\_list):}
    \BuiltInTok{print}\NormalTok{(}\SpecialStringTok{f"  }\SpecialCharTok{\{}\NormalTok{i}\SpecialCharTok{\}}\SpecialStringTok{: }\SpecialCharTok{\{}\NormalTok{item}\SpecialCharTok{\}}\SpecialStringTok{"}\NormalTok{)}
\end{Highlighting}
\end{Shaded}

  This gives numbered output, one item per line, that is easy to
  navigate and understand when heard.
\item
  \textbf{\texttt{len()} as a safety check before index access:} Teach
  students to always know the length before indexing into a list. This
  prevents \texttt{IndexError}. Make \texttt{len(my\_list)} the first
  thing printed when debugging a list problem.
\item
  \textbf{Index arithmetic verbalized:} When accessing
  \texttt{my\_list{[}-1{]}}, say ``the last element, which is at index
  \texttt{len(my\_list)\ minus\ 1}, which Python shortens to
  \texttt{-1}.'' Connecting the shorthand to the underlying meaning
  prevents confusion.
\end{itemize}

\paragraph{Thought Experiments}\label{thought-experiments-10}

\begin{enumerate}
\def\labelenumi{\arabic{enumi}.}
\item
  \emph{``Zero-based indexing means the first element is at position 0.
  This feels unnatural but is standard across almost all programming
  languages. Why? (Hint: think about memory addresses --- the address of
  the first element, and the address of any element as an offset from
  there.)''}
\item
  \emph{``A list allows duplicate values. A set (which you'll learn
  later) does not. When would duplicates in a list be meaningful data?
  When would they indicate an error in your program?''}
\item
  \emph{``A \texttt{for} loop iterates over every element in a list. Is
  there any computation a \texttt{for} loop can do that a \texttt{while}
  loop cannot? Is there any computation a \texttt{while} loop can do
  that a \texttt{for} loop cannot?''}
\item
  \emph{``What happens when you iterate over an empty list:
  \texttt{for\ item\ in\ {[}{]}:}? Does Python crash, print an error, or
  do nothing? Why is `do nothing' the correct behavior rather than an
  error?''}
\end{enumerate}

\begin{center}\rule{0.5\linewidth}{0.5pt}\end{center}

\subsubsection{Lesson 3.5: String
Operations}\label{lesson-3.5-string-operations}

\paragraph{Instructional Cues}\label{instructional-cues-9}

\textbf{Immutability demonstration:}

\begin{Shaded}
\begin{Highlighting}[]
\NormalTok{text }\OperatorTok{=} \StringTok{"hello"}
\NormalTok{text.upper()        }\CommentTok{\# Creates a new string {-}{-}{-} doesn\textquotesingle{}t modify text}
\BuiltInTok{print}\NormalTok{(text)         }\CommentTok{\# Still "hello"}
\NormalTok{text }\OperatorTok{=}\NormalTok{ text.upper() }\CommentTok{\# Now reassigned}
\BuiltInTok{print}\NormalTok{(text)         }\CommentTok{\# "HELLO"}
\end{Highlighting}
\end{Shaded}

This demonstration catches a persistent beginner misunderstanding. Ask
the student to predict the output of \texttt{print(text)} after
\texttt{text.upper()} before running it.

\textbf{Method chaining verbalized:}
\texttt{"\ \ hello\ world\ \ ".strip().title()} --- read this as:
``Start with the string. Strip whitespace from both ends. Then convert
the result to title case.''

\paragraph{Thought Experiments}\label{thought-experiments-11}

\begin{enumerate}
\def\labelenumi{\arabic{enumi}.}
\item
  \emph{``Strings are immutable --- they cannot be changed, only
  replaced. Lists are mutable --- they can be changed in-place. What
  architectural difference between strings and lists makes one mutable
  and the other not? (Think about how they're stored in memory.)''}
\item
  \emph{``\texttt{.split()} splits a string into a list of words.
  \texttt{.join()} combines a list of strings into one string. They are
  inverses. What is the inverse of \texttt{.replace()}? Does it have a
  direct inverse? If not, how would you reverse a replacement?''}
\item
  \emph{``String formatting with f-strings:
  \texttt{f\textquotesingle{}Value\ is\ \{x:.2f\}\textquotesingle{}}
  formats \texttt{x} as a float with 2 decimal places. Why does a
  programming language need a mini-language for formatting numbers? What
  would programs look like if we had to format every number manually
  with division and modulo?''}
\end{enumerate}

\paragraph{Alternative Resources}\label{alternative-resources-3}

\begin{itemize}
\tightlist
\item
  \textbf{Python String Methods documentation}
  (docs.python.org/3/library/stdtypes.html\#string-methods): Plain HTML
  table of all string methods. Screen-reader navigable. Essential
  reference.
\item
  \textbf{``Python for Data Analysis'' by Wes McKinney:} Chapter on
  string manipulation provides real-world examples beyond what this
  curriculum covers. For motivated students.
\end{itemize}

\begin{center}\rule{0.5\linewidth}{0.5pt}\end{center}

\subsection{9. Unit 4 --- Data Structures: Lists and
Dictionaries}\label{unit-4-data-structures-lists-and-dictionaries}

\emph{(See all lessons above for list material; this section adds
dictionary-specific guidance.)}

\subsubsection{Dictionary Lessons}\label{dictionary-lessons}

\paragraph{Instructional Cues}\label{instructional-cues-10}

\textbf{Multiple mental models for dictionaries:}

\begin{itemize}
\tightlist
\item
  \emph{Real dictionary:} You look up a word (key) to find its
  definition (value). Words are unique; definitions can be anything.
\item
  \emph{Contact list:} Each name (key) maps to a phone number (value).
  No two contacts have the same name in a well-maintained list.
\item
  \emph{JSON configuration file:} Settings names map to their values.
  This foreshadows Unit 8.
\end{itemize}

\textbf{The \texttt{.get()} method vs.~direct access:}

\begin{Shaded}
\begin{Highlighting}[]
\CommentTok{\# This crashes on missing key}
\BuiltInTok{print}\NormalTok{(my\_dict[}\StringTok{"nonexistent\_key"}\NormalTok{])   }\CommentTok{\# KeyError}

\CommentTok{\# This returns a default {-}{-}{-} never crashes}
\BuiltInTok{print}\NormalTok{(my\_dict.get(}\StringTok{"nonexistent\_key"}\NormalTok{, }\StringTok{"Not found"}\NormalTok{))  }\CommentTok{\# "Not found"}
\end{Highlighting}
\end{Shaded}

For a screen reader user who may not immediately catch a crash in
terminal output, teaching \texttt{.get()} as the default access pattern
reduces debugging overhead.

\paragraph{Thought Experiments}\label{thought-experiments-12}

\begin{enumerate}
\def\labelenumi{\arabic{enumi}.}
\item
  \emph{``Dictionaries use keys to retrieve values in O(1) time ---
  essentially instant regardless of dictionary size. Lists use indices,
  also O(1). But searching a list for a specific value (not by index)
  takes O(n) time --- proportional to list length. When does this
  difference matter in practice?''}
\item
  \emph{``Your dictionary cannot have duplicate keys --- adding a second
  value for an existing key replaces the first. What real-world data
  problems does this property solve? What problems does it create?''}
\item
  \emph{``JSON (JavaScript Object Notation) is essentially Python
  dictionaries written in a slightly different format. What does it mean
  that a data format invented for JavaScript has become the universal
  language of data exchange? What does this tell you about the value of
  common standards?''}
\item
  \emph{``A `nested dictionary' contains other dictionaries as values.
  Model a university course catalog as a nested dictionary. What are the
  outer keys? What are the inner keys? What would
  \texttt{catalog{[}\textquotesingle{}CS\textquotesingle{}{]}{[}\textquotesingle{}101\textquotesingle{}{]}{[}\textquotesingle{}instructor\textquotesingle{}{]}}
  represent?''}
\end{enumerate}

\begin{center}\rule{0.5\linewidth}{0.5pt}\end{center}

\subsection{10. Unit 5 --- User Input and While
Loops}\label{unit-5-user-input-and-while-loops}

\subsubsection{Thought Experiments}\label{thought-experiments-13}

\begin{enumerate}
\def\labelenumi{\arabic{enumi}.}
\item
  \emph{``A while loop with a flag variable and a while loop with
  \texttt{break} can often accomplish the same thing. Compare these two
  equivalent programs and explain when you would prefer each:}

\begin{Shaded}
\begin{Highlighting}[]
\CommentTok{\# Version A: flag}
\NormalTok{active }\OperatorTok{=} \VariableTok{True}
\ControlFlowTok{while}\NormalTok{ active:}
\NormalTok{    cmd }\OperatorTok{=} \BuiltInTok{input}\NormalTok{(}\StringTok{"\textgreater{} "}\NormalTok{)}
    \ControlFlowTok{if}\NormalTok{ cmd }\OperatorTok{==} \StringTok{"quit"}\NormalTok{: active }\OperatorTok{=} \VariableTok{False}
    \ControlFlowTok{else}\NormalTok{: }\BuiltInTok{print}\NormalTok{(}\StringTok{"You said:"}\NormalTok{, cmd)}

\CommentTok{\# Version B: break}
\ControlFlowTok{while} \VariableTok{True}\NormalTok{:}
\NormalTok{    cmd }\OperatorTok{=} \BuiltInTok{input}\NormalTok{(}\StringTok{"\textgreater{} "}\NormalTok{)}
    \ControlFlowTok{if}\NormalTok{ cmd }\OperatorTok{==} \StringTok{"quit"}\NormalTok{: }\ControlFlowTok{break}
    \BuiltInTok{print}\NormalTok{(}\StringTok{"You said:"}\NormalTok{, cmd)}
\end{Highlighting}
\end{Shaded}

  \emph{Which is clearer? Which is easier to extend if you add more exit
  conditions?{\kern0pt}``}
\item
  \emph{``An infinite loop that waits for user input is the core pattern
  of every interactive program --- email clients, video games, operating
  systems. Describe the `event loop' of your screen reader in terms of
  while loops and if statements. What is it waiting for? What does it do
  when it gets an event?''}
\item
  \emph{``The \texttt{continue} statement skips the rest of the current
  iteration. An alternative is to wrap the loop body in an \texttt{if}
  statement. When is \texttt{continue} cleaner? When does the
  \texttt{if} approach produce clearer code?''}
\end{enumerate}

\begin{center}\rule{0.5\linewidth}{0.5pt}\end{center}

\subsection{11. Unit 6 --- Functions}\label{unit-6-functions}

\subsubsection{Thought Experiments (in addition to existing
guide)}\label{thought-experiments-in-addition-to-existing-guide}

\begin{enumerate}
\def\labelenumi{\arabic{enumi}.}
\item
  \emph{``A function is described as a `black box' --- you give it
  inputs and get outputs without knowing the internal details. What are
  the benefits of this abstraction for the programmer who writes the
  function? For the programmer who uses it? What are the limits of this
  metaphor?''}
\item
  \emph{``Default parameter values:
  \texttt{def\ greet(name=\textquotesingle{}World\textquotesingle{})}.
  This allows the function to be called with or without an argument. How
  does this enable backward compatibility --- adding a new feature
  without breaking existing code that calls the function?''}
\item
  \emph{``Functions should `do one thing.' What does `one thing' mean?
  Is \texttt{calculate\_and\_print\_tax()} one thing or two? If you
  separate it into \texttt{calculate\_tax()} and \texttt{print\_tax()},
  what becomes possible that wasn't before?''}
\item
  \_``Recursive functions call themselves.
  \texttt{factorial(n)\ =\ n\ \_\ factorial(n-1)}. What is the base case
  and why is it required? What goes wrong if a recursive function has no
  base case?''*
\item
  \_``A function is a way to name a piece of code. Why does naming
  matter in programming? What is the cognitive difference between
  \texttt{x\ =\ a\ \_\ b\ +\ c\ -\ d}and\texttt{x\ =\ calculate\_net\_profit(revenue,\ cost,\ overhead,\ taxes)}?''*
\end{enumerate}

\paragraph{Alternative Resources}\label{alternative-resources-4}

\begin{itemize}
\tightlist
\item
  \textbf{``Structure and Interpretation of Computer Programs'' (SICP)
  --- MIT OpenCourseWare:} Conceptually deep treatment of functions as
  first-class objects. Lectures have transcripts. For advanced students.
\item
  \textbf{Exercism.io Python Track:} Mentor-guided exercises with
  community feedback. Web platform with reasonable accessibility.
  Students can complete exercises at their own pace.
\item
  \textbf{Python Tutor (pythontutor.com):} Step-by-step execution with
  variable state display. The text-based ``state panel'' below the code
  is screen-reader accessible and shows variable values at each step.
\end{itemize}

\begin{center}\rule{0.5\linewidth}{0.5pt}\end{center}

\subsection{12. Unit 7 --- Classes and Object-Oriented
Programming}\label{unit-7-classes-and-object-oriented-programming}

\subsubsection{Extended Instructional
Cues}\label{extended-instructional-cues}

\textbf{Staged introduction to OOP --- do not rush to syntax:}

Stage 1 (no code): Discuss the concept. What is an object? What does
``bundling data and behavior'' mean? Use tangible examples the student
knows.

Stage 2 (pseudocode): Write a class in plain English before Python
syntax:

\begin{verbatim}
Class: BankAccount
  Has: owner_name, balance, account_number
  Can do: deposit(amount), withdraw(amount), get_balance()
\end{verbatim}

Stage 3 (Python skeleton): Write the class structure with empty methods
and docstrings.

Stage 4 (full implementation): Fill in method bodies.

This staging prevents syntax confusion from obscuring the conceptual
learning.

\textbf{\texttt{\_\_str\_\_} is mandatory for accessibility:} Every
class a student creates must define \texttt{\_\_str\_\_}. Without it,
\texttt{print(my\_object)} produces
\texttt{\textless{}\_\_main\_\_.BankAccount\ object\ at\ 0x7f...\textgreater{}}
--- meaningless when heard. The requirement to define
\texttt{\_\_str\_\_} also reinforces thinking about what the object's
most important information is.

\begin{Shaded}
\begin{Highlighting}[]
\KeywordTok{def} \FunctionTok{\_\_str\_\_}\NormalTok{(}\VariableTok{self}\NormalTok{):}
    \ControlFlowTok{return} \SpecialStringTok{f"BankAccount(owner=}\SpecialCharTok{\{}\VariableTok{self}\SpecialCharTok{.}\NormalTok{owner\_name}\SpecialCharTok{\}}\SpecialStringTok{, balance=$}\SpecialCharTok{\{}\VariableTok{self}\SpecialCharTok{.}\NormalTok{balance}\SpecialCharTok{:.2f\}}\SpecialStringTok{)"}
\end{Highlighting}
\end{Shaded}

\paragraph{Thought Experiments}\label{thought-experiments-14}

\begin{enumerate}
\def\labelenumi{\arabic{enumi}.}
\item
  \emph{``Encapsulation means hiding internal details. A car's engine is
  encapsulated under the hood --- you don't need to know how fuel
  injection works to drive. What are the analogous `hidden internals' in
  a BankAccount class? What should be private? What should be public?
  Why?''}
\item
  \emph{``Inheritance creates `is-a' relationships. Composition creates
  `has-a' relationships. A \texttt{SavingsAccount} is-a
  \texttt{BankAccount}. A \texttt{Bank} has-a list of
  \texttt{BankAccount} objects. Design a simple banking system using
  both relationships. What would go wrong if you tried to use only
  inheritance?''}
\item
  \emph{``Method overriding allows a subclass to replace a parent's
  method. When is this the right design? When should you call
  \texttt{super().method()} inside the override vs.~replace the method
  entirely?''}
\item
  \emph{``Python allows multiple inheritance (inheriting from two parent
  classes). Most languages do not permit this. What problems does
  multiple inheritance create? What problems does it solve? Would you
  use it?''}
\item
  \emph{``Object-oriented programming was supposed to make code more
  reusable. Has it delivered on this promise? What kinds of code reuse
  are well-served by OOP? What kinds aren't?''}
\end{enumerate}

\paragraph{Alternative Resources}\label{alternative-resources-5}

\begin{itemize}
\tightlist
\item
  \textbf{``Object-Oriented Programming in Python'' by Michael
  Goldwasser and David Letscher:} Textbook specifically pairing Python
  and OOP concepts. More formal treatment for students who want depth.
\item
  \textbf{``Fluent Python'' by Luciano Ramalho:} Advanced but deeply
  accessible to screen readers. Covers Python's data model including
  \texttt{\_\_str\_\_}, \texttt{\_\_repr\_\_}, and special methods. For
  motivated students.
\item
  \textbf{Real Python ``Object-Oriented Programming in Python'' tutorial
  (realpython.com/python3-object-oriented-programming/):} Long but
  thorough. Plain HTML.
\end{itemize}

\begin{center}\rule{0.5\linewidth}{0.5pt}\end{center}

\subsection{13. Unit 8 --- Files and
Exceptions}\label{unit-8-files-and-exceptions}

\subsubsection{Extended Instructional
Cues}\label{extended-instructional-cues-1}

\textbf{Files as persistence --- make it tangible:} Run a program that
stores a number in a variable, quits, and starts again --- the variable
is gone. Then run the same program storing the number in a file --- it
persists. Ask: ``What changed? What is different about a file vs.~a
variable?'' Let the student reason through it before explaining.

\textbf{The \texttt{with} statement --- demonstrate the alternative:}

\begin{Shaded}
\begin{Highlighting}[]
\CommentTok{\# Without \textquotesingle{}with\textquotesingle{} {-}{-}{-} dangerous}
\NormalTok{f }\OperatorTok{=} \BuiltInTok{open}\NormalTok{(}\StringTok{"data.txt"}\NormalTok{, }\StringTok{"r"}\NormalTok{)}
\NormalTok{contents }\OperatorTok{=}\NormalTok{ f.read()}
\CommentTok{\# If an error occurs here, f.close() is never called}
\NormalTok{f.close()}

\CommentTok{\# With \textquotesingle{}with\textquotesingle{} {-}{-}{-} safe}
\ControlFlowTok{with} \BuiltInTok{open}\NormalTok{(}\StringTok{"data.txt"}\NormalTok{, }\StringTok{"r"}\NormalTok{) }\ImportTok{as}\NormalTok{ f:}
\NormalTok{    contents }\OperatorTok{=}\NormalTok{ f.read()}
\CommentTok{\# f is closed automatically even if an error occurs}
\end{Highlighting}
\end{Shaded}

\textbf{Exception types --- teach by demonstration, not memorization:}
Create a ``crash collection'' --- a series of small programs that each
produce a specific exception. Let the student run each one, read the
error, and discuss why it occurred before seeing the solution. The goal
is pattern recognition, not memorization.

\paragraph{Thought Experiments}\label{thought-experiments-15}

\begin{enumerate}
\def\labelenumi{\arabic{enumi}.}
\item
  \emph{``Files are persistent storage --- they survive when your
  program exits. Variables are temporary --- they're gone when the
  program ends. What determines which data should be stored in files
  vs.~kept only in memory? What are the tradeoffs?''}
\item
  \emph{``The \texttt{with} statement uses Python's `context manager'
  protocol. This protocol says: `when this block starts, do X; when it
  ends (even if it ends with an error), do Y.' What other situations in
  programming require guaranteed cleanup? What happens in real systems
  when cleanup is skipped?''}
\item
  \emph{``Exceptions propagate up the call stack --- if function A calls
  B which calls C, and C raises an exception, it moves up to B, then A,
  until something catches it. If nothing catches it, the program
  crashes. Design the exception handling strategy for a program that
  reads a config file, parses it, and then connects to a database. Where
  should exceptions be caught?''}
\item
  \emph{``JSON stores data as text. Binary formats (like \texttt{pickle}
  in Python) store data in a compact binary form. JSON is
  human-readable; binary is not. JSON is language-independent; pickle is
  Python-only. For each of these use cases, which format would you
  choose: (a) a config file the user might edit by hand, (b) a large
  dataset processed only by your own Python code, (c) data shared
  between a Python backend and a JavaScript frontend?''}
\end{enumerate}

\begin{center}\rule{0.5\linewidth}{0.5pt}\end{center}

\subsection{14. Unit 9 --- Testing and Data
Analysis}\label{unit-9-testing-and-data-analysis}

\subsubsection{Extended Instructional
Cues}\label{extended-instructional-cues-2}

\textbf{Tests as documentation --- make this concrete:} Show two
identical functions --- one with tests that describe the expected
behavior, one without. Ask: ``Which function is easier to understand
without reading the implementation?'' The tests are a specification in
executable form.

\textbf{Test-driven development (TDD) as a preview:} For motivated
students, introduce TDD: write the test before the implementation. This
inverts the usual workflow:

\begin{enumerate}
\def\labelenumi{\arabic{enumi}.}
\tightlist
\item
  Write a test that fails (because the function doesn't exist yet).
\item
  Write the minimum implementation to make the test pass.
\item
  Refactor.
\end{enumerate}

This discipline is genuinely useful for beginners because it forces them
to think about what a function should do before worrying about how.

\paragraph{Thought Experiments}\label{thought-experiments-16}

\begin{enumerate}
\def\labelenumi{\arabic{enumi}.}
\item
  \emph{``Unit testing tests one function in isolation. Integration
  testing tests multiple components working together. System testing
  tests the entire program. Why have three levels rather than just one?
  What does each level catch that the others might miss?''}
\item
  \emph{``A test can only prove that a program behaves correctly for the
  inputs tested. It cannot prove correctness for all possible inputs.
  What does this mean for the value of testing? How do you decide which
  inputs to test?''}
\item
  \emph{``Random numbers make testing difficult --- if the output
  changes every run, how do you write an assertion? Describe a strategy
  for testing a function that uses \texttt{random.randint()}. (Hint:
  think about seeding.)''}
\item
  \emph{``CSV data from the real world is often `dirty' --- missing
  values, extra commas, inconsistent formatting, encoding issues. What
  does a robust CSV parser need to handle? What assumptions would a
  naive parser make that would cause it to fail on real-world data?''}
\end{enumerate}

\begin{center}\rule{0.5\linewidth}{0.5pt}\end{center}

\subsection{15. Unit 10 --- Project-Based
Learning}\label{unit-10-project-based-learning}

\subsubsection{Extended Instructional
Cues}\label{extended-instructional-cues-3}

\textbf{Design document requirements:} Require a written design document
before any code is written. For a screen reader user, a plain text file
works best. The document should include:

\begin{verbatim}
Project Name: [title]
One-sentence purpose: [what problem does this solve?]
Target user: [who will use this?]

Features (ranked by priority):
1. [most important feature]
2. ...

Data I will store: [what information needs to persist?]
Data structures: [which Python structures will hold this?]

Functions I will write:
- [function name](inputs) -> output: [brief description]
- ...

Main loop pseudocode:
[plain English description of program flow]

Edge cases I have considered:
- [what happens if the user enters nothing?]
- [what happens if the file doesn't exist?]
- ...
\end{verbatim}

\textbf{Iterative development enforcement:} Require demonstrating a
working version at the end of every session, no matter how minimal.
``Working'' means runnable and producing correct output for at least one
input. A student who cannot demonstrate working code at session end
needs to reduce scope, not work harder overnight.

\paragraph{Thought Experiments}\label{thought-experiments-17}

\begin{enumerate}
\def\labelenumi{\arabic{enumi}.}
\item
  \emph{``A program designed for sighted users might use color to
  indicate status (red for error, green for success), icons to represent
  actions, and layout to group related information. How would you
  replicate each of these design patterns in a screen-reader-accessible
  terminal program?''}
\item
  \emph{``Your contact manager has a search function. How many different
  search behaviors could be useful? Search by exact name match, partial
  match, phone prefix, email domain, notes keywords? What should the
  program do when there are 0 results? 1 result? Multiple results?''}
\item
  \emph{``Error messages should be actionable --- they should tell the
  user not just what went wrong, but what to do next. Rewrite these
  unhelpful messages to be actionable: `Error: invalid input', `File not
  found', `Operation failed.'\,''}
\end{enumerate}

\begin{center}\rule{0.5\linewidth}{0.5pt}\end{center}

\subsection{16. Unit 11 --- Capstone
Project}\label{unit-11-capstone-project}

\subsubsection{Assessment Guidance}\label{assessment-guidance}

\paragraph{Process Checkpoints}\label{process-checkpoints}

{\def\LTcaptype{none} % do not increment counter
\begin{longtable}[]{@{}
  >{\raggedright\arraybackslash}p{(\linewidth - 4\tabcolsep) * \real{0.1917}}
  >{\raggedright\arraybackslash}p{(\linewidth - 4\tabcolsep) * \real{0.1250}}
  >{\raggedright\arraybackslash}p{(\linewidth - 4\tabcolsep) * \real{0.6833}}@{}}
\toprule\noalign{}
\begin{minipage}[b]{\linewidth}\raggedright
Checkpoint
\end{minipage} & \begin{minipage}[b]{\linewidth}\raggedright
Timing
\end{minipage} & \begin{minipage}[b]{\linewidth}\raggedright
What to Evaluate
\end{minipage} \\
\midrule\noalign{}
\endhead
\bottomrule\noalign{}
\endlastfoot
Written design document & Before any code & Is the scope realistic? Are
features prioritized? Are data structures appropriate? \\
Data model working & End of week 1 & Can items be created, stored in
memory, and printed? \\
Core CRUD functional & End of week 2 & Add, view, edit, delete all
work? \\
Persistence working & End of week 3 & Data survives program restart? \\
Error handling complete & End of week 4 & Program handles all invalid
inputs without crashing? \\
Accessibility review & End of week 4 & Is all output meaningful when
read aloud? \\
Final demo & End & Full demonstration with screen reader active \\
\end{longtable}
}

\paragraph{The Screen Reader Demo
Rubric}\label{the-screen-reader-demo-rubric}

The final demonstration should be performed entirely with the screen
reader active. Evaluate each:

{\def\LTcaptype{none} % do not increment counter
\begin{longtable}[]{@{}
  >{\raggedright\arraybackslash}p{(\linewidth - 6\tabcolsep) * \real{0.1316}}
  >{\raggedright\arraybackslash}p{(\linewidth - 6\tabcolsep) * \real{0.3474}}
  >{\raggedright\arraybackslash}p{(\linewidth - 6\tabcolsep) * \real{0.2737}}
  >{\raggedright\arraybackslash}p{(\linewidth - 6\tabcolsep) * \real{0.2474}}@{}}
\toprule\noalign{}
\begin{minipage}[b]{\linewidth}\raggedright
Criterion
\end{minipage} & \begin{minipage}[b]{\linewidth}\raggedright
Excellent
\end{minipage} & \begin{minipage}[b]{\linewidth}\raggedright
Satisfactory
\end{minipage} & \begin{minipage}[b]{\linewidth}\raggedright
Needs Work
\end{minipage} \\
\midrule\noalign{}
\endhead
\bottomrule\noalign{}
\endlastfoot
\textbf{Menu navigation} & Student navigates menu by listening; confirms
choice before acting & Student navigates but occasionally needs visual
cues & Student needs visual confirmation of menu state \\
\textbf{Input prompts} & All prompts clearly describe expected format
and type & Most prompts are descriptive & Some prompts are vague or
missing \\
\textbf{Confirmation messages} & Every action produces a clear spoken
confirmation & Most actions confirmed & Some actions succeed silently \\
\textbf{Error messages} & All error messages are specific and actionable
& Most errors give useful messages & Some errors produce generic or
missing messages \\
\textbf{Data persistence} & Student demonstrates save and load by
quitting and restarting & Save/load works but student doesn't
demonstrate live & Save or load not fully functional \\
\textbf{Code explanation} & Student can explain any line of their code
when asked & Student can explain most of their code & Student cannot
explain sections of their code \\
\end{longtable}
}

\begin{center}\rule{0.5\linewidth}{0.5pt}\end{center}

\subsection{17. Assessment Strategies}\label{assessment-strategies}

\subsubsection{Accessible Assessment
Design}\label{accessible-assessment-design}

Traditional programming assessments (written tests, timed coding) create
unnecessary barriers for blind students. Use alternatives that assess
the same competencies:

\textbf{Think-aloud problems:} Give the student a short piece of code
and ask them to read it aloud and explain what it does, step by step.
This assesses reading comprehension and debugging ability without
requiring written answers.

\textbf{Debugging exercises:} Provide a working program and a broken
version. Ask the student to identify the error using only their screen
reader and debugging tools. This directly tests real-world skill.

\textbf{Feature extension:} Give a working program and ask the student
to add a specific feature. This is a more realistic assessment than
writing from scratch and allows scaffolding.

\textbf{Verbal design review:} Present a programming problem and ask the
student to describe their approach --- data structures, functions,
algorithm --- without writing any code. Assess reasoning quality.

\textbf{Portfolio review:} Maintain a collection of all programs written
during the course. Assess growth, consistency, commenting quality, and
the student's ability to explain their own work.

\subsubsection{Grading Considerations}\label{grading-considerations}

\begin{itemize}
\tightlist
\item
  Do not penalize for missing visual formatting (e.g., a missing newline
  that only matters visually).
\item
  Do reward for accessibility features the student adds above
  requirements: clear prompts, confirmation messages, help text.
\item
  Recognize that a blind student may spend significantly more time
  debugging than a sighted peer on identical problems. Factor this into
  time expectations.
\item
  Consider allowing oral explanations as alternatives to written code
  comments for students who find comment-writing slow.
\end{itemize}

\begin{center}\rule{0.5\linewidth}{0.5pt}\end{center}

\subsection{18. Supporting Diverse Needs Within the Visually Impaired
Population}\label{supporting-diverse-needs-within-the-visually-impaired-population}

\subsubsection{For Students with Low
Vision}\label{for-students-with-low-vision}

\begin{itemize}
\tightlist
\item
  Use high-contrast VS Code themes: Dark High Contrast, Light High
  Contrast (available under Color Themes in settings).
\item
  Increase editor font size: \texttt{Ctrl+Shift+P} -\textgreater{}
  ``Preferences: Open Settings'' -\textgreater{}
  \texttt{editor.fontSize}.
\item
  Enable zoom magnification: \texttt{Ctrl++} to zoom the VS Code window.
\item
  Test whether syntax highlighting is visible with their specific visual
  condition. Some students benefit from; others are confused by color
  information they cannot fully perceive.
\end{itemize}

\subsubsection{For Students Who Are
Deaf-Blind}\label{for-students-who-are-deaf-blind}

\begin{itemize}
\tightlist
\item
  All communication with the teacher must be through the braille display
  or tactile signing --- no spoken instruction.
\item
  Screen reader audio verbosity settings must be managed by the student
  (since they cannot hear audio feedback during class).
\item
  Consider pairing with a trained interpreter who knows both tactile
  communication and computing concepts.
\item
  Review the Perkins School for the Blind's resources on deaf-blind
  education technology.
\end{itemize}

\subsubsection{For Students with Physical Disabilities in Addition to
Visual
Impairment}\label{for-students-with-physical-disabilities-in-addition-to-visual-impairment}

\begin{itemize}
\tightlist
\item
  Switch access: VS Code is fully keyboard-navigable and can be operated
  with adapted keyboards, switch interfaces, or eye-gaze systems paired
  with virtual keyboards.
\item
  Voice control: Dragon NaturallySpeaking and macOS Voice Control can
  both operate VS Code. Voice coding requires additional voice commands
  for symbols.
\item
  Sip-and-puff devices: Can control keyboard input. Students using these
  will have different typing rates --- extend all time-based tasks.
\end{itemize}

\subsubsection{For Students with Cognitive Differences in Addition to
Visual
Impairment}\label{for-students-with-cognitive-differences-in-addition-to-visual-impairment}

\begin{itemize}
\tightlist
\item
  Reduce the amount of new vocabulary per session.
\item
  Use consistent metaphors across the entire course --- do not introduce
  a new metaphor for the same concept.
\item
  Provide written summaries of every lesson that the student can review
  with their screen reader later.
\item
  Break complex concepts into smaller steps; complete each step before
  introducing the next.
\item
  Build in regular review sessions rather than only forward-moving
  instruction.
\end{itemize}

\begin{center}\rule{0.5\linewidth}{0.5pt}\end{center}

\subsection{19. Teaching AI Tools
Responsibly}\label{teaching-ai-tools-responsibly}

\subsubsection{Framing AI for the
Classroom}\label{framing-ai-for-the-classroom}

Establish a clear classroom policy before students begin using AI tools:

\textbf{Permitted:} Using AI to explain error messages, generate a first
draft for the student to critique and improve, explain a concept a
different way, suggest test cases to consider.

\textbf{Not permitted:} Submitting AI-generated code as completed work
without reading and understanding every line, using AI to complete
assessments that are meant to evaluate individual skill.

\textbf{Required when AI is used:} The student must be able to explain
any AI-generated code they submit as if they wrote it themselves.
Inability to explain code is evidence it was not understood.

\subsubsection{AI Tools and Blind Students --- Specific
Considerations}\label{ai-tools-and-blind-students-specific-considerations}

\textbf{Potential accessibility benefit:} AI tools like GitHub Copilot
are text-based --- suggestions appear as text that a screen reader can
read. Copilot's Accessible View integration in VS Code is improving.

\textbf{Potential accessibility risk:} AI may generate code that is
visually clear but auditorily confusing --- complex nested structures,
long lines, or dense comprehensions that are harder to navigate with a
screen reader than to read with eyes.

\textbf{Teaching point:} Ask students to evaluate AI-generated code for
screen-reader accessibility:

\begin{itemize}
\tightlist
\item
  Are variable names descriptive?
\item
  Are long lines split into readable pieces?
\item
  Are print statements and prompts clear and complete?
\item
  Does the code have docstrings?
\end{itemize}

If not, the AI-generated code needs to be modified before use ---
regardless of whether it runs correctly.

\subsubsection{Thought Experiments on
AI}\label{thought-experiments-on-ai}

\begin{enumerate}
\def\labelenumi{\arabic{enumi}.}
\item
  \emph{``An AI writes code that solves your problem but you don't
  understand it. You submit it. It passes tests. Have you learned
  anything? What is the difference between `having a solution' and
  `knowing how to solve problems'?''}
\item
  \emph{``AI code generators are trained on billions of lines of code
  written mostly by sighted programmers. What biases might this
  introduce in the code they generate? Think specifically about:
  variable names, output design, prompts, error messages,
  documentation.''}
\item
  \emph{``If AI can write most code, which skills become more valuable
  for future programmers? Consider: understanding requirements,
  debugging, testing, communicating about code, evaluating solutions,
  and accessibility design.''}
\end{enumerate}

\begin{center}\rule{0.5\linewidth}{0.5pt}\end{center}

\subsection{20. Alternative and Supplementary
Resources}\label{alternative-and-supplementary-resources}

\subsubsection{Textbooks and Reading}\label{textbooks-and-reading}

{\def\LTcaptype{none} % do not increment counter
\begin{longtable}[]{@{}
  >{\raggedright\arraybackslash}p{(\linewidth - 6\tabcolsep) * \real{0.3706}}
  >{\raggedright\arraybackslash}p{(\linewidth - 6\tabcolsep) * \real{0.1538}}
  >{\raggedright\arraybackslash}p{(\linewidth - 6\tabcolsep) * \real{0.1189}}
  >{\raggedright\arraybackslash}p{(\linewidth - 6\tabcolsep) * \real{0.3566}}@{}}
\toprule\noalign{}
\begin{minipage}[b]{\linewidth}\raggedright
Resource
\end{minipage} & \begin{minipage}[b]{\linewidth}\raggedright
Level
\end{minipage} & \begin{minipage}[b]{\linewidth}\raggedright
Format
\end{minipage} & \begin{minipage}[b]{\linewidth}\raggedright
Accessibility
\end{minipage} \\
\midrule\noalign{}
\endhead
\bottomrule\noalign{}
\endlastfoot
\textbf{Python Crash Course, 3rd Ed.} (Matthes) & Beginner & Print/ebook
& Good; primary textbook for this curriculum \\
\textbf{Python Programming: An Introduction to CS} (Zelle) & Beginner &
Print/ebook & Good; companion textbook \\
\textbf{Automate the Boring Stuff with Python} (Sweigart) &
Beginner--Intermediate & Free HTML / print & Excellent; HTML version
fully accessible \\
\textbf{Think Python} (Downey) & Beginner & Free HTML / print &
Excellent HTML version \\
\textbf{Learn Python 3 the Hard Way} (Shaw) & Beginner & Print/ebook &
Plain text exercises, highly accessible \\
\textbf{Python for Everybody} (Severance) & Beginner & Free HTML + video
& Video lectures have transcripts; good accessibility \\
\textbf{Fluent Python} (Ramalho) & Advanced & Print/ebook & Dense but
thorough; for motivated advanced students \\
\textbf{Python Docs} (python.org/3) & Reference & Free HTML & Excellent
screen-reader accessibility; definitive \\
\end{longtable}
}

\subsubsection{Online Courses and
Platforms}\label{online-courses-and-platforms}

{\def\LTcaptype{none} % do not increment counter
\begin{longtable}[]{@{}
  >{\raggedright\arraybackslash}p{(\linewidth - 6\tabcolsep) * \real{0.1524}}
  >{\raggedright\arraybackslash}p{(\linewidth - 6\tabcolsep) * \real{0.2378}}
  >{\raggedright\arraybackslash}p{(\linewidth - 6\tabcolsep) * \real{0.3293}}
  >{\raggedright\arraybackslash}p{(\linewidth - 6\tabcolsep) * \real{0.2805}}@{}}
\toprule\noalign{}
\begin{minipage}[b]{\linewidth}\raggedright
Platform
\end{minipage} & \begin{minipage}[b]{\linewidth}\raggedright
Course
\end{minipage} & \begin{minipage}[b]{\linewidth}\raggedright
Screen Reader Compatibility
\end{minipage} & \begin{minipage}[b]{\linewidth}\raggedright
Notes
\end{minipage} \\
\midrule\noalign{}
\endhead
\bottomrule\noalign{}
\endlastfoot
\textbf{edX / CS50P} (Harvard) & Introduction to Programming with Python
& High --- video transcripts, accessible exercises & Free audit. Very
high quality. \\
\textbf{Coursera} & Python for Everybody (Michigan) & Medium --- some
interactive elements inaccessible & Free audit available \\
\textbf{Codecademy} & Learn Python 3 & Partial --- editor accessible,
some UI not & Requires testing with specific screen reader \\
\textbf{SoloLearn} & Python Core & High --- mobile app works well with
VoiceOver/TalkBack & Good for mobile-first learners \\
\textbf{freeCodeCamp} & Python curriculum & Medium --- varies by module
& Free; large community \\
\textbf{Exercism} & Python Track & High --- terminal-based, mentor
feedback & Excellent for structured practice after basics \\
\textbf{CodingBat} & Python problems & High --- plain HTML,
browser-based & Good for logic practice; no complex UI \\
\textbf{replit.com} & Any Python project & Medium --- test with
student's setup & No-install option; accessibility variable \\
\end{longtable}
}

\subsubsection{Video Resources}\label{video-resources}

{\def\LTcaptype{none} % do not increment counter
\begin{longtable}[]{@{}
  >{\raggedright\arraybackslash}p{(\linewidth - 4\tabcolsep) * \real{0.2993}}
  >{\raggedright\arraybackslash}p{(\linewidth - 4\tabcolsep) * \real{0.3061}}
  >{\raggedright\arraybackslash}p{(\linewidth - 4\tabcolsep) * \real{0.3946}}@{}}
\toprule\noalign{}
\begin{minipage}[b]{\linewidth}\raggedright
Resource
\end{minipage} & \begin{minipage}[b]{\linewidth}\raggedright
Description
\end{minipage} & \begin{minipage}[b]{\linewidth}\raggedright
Accessibility
\end{minipage} \\
\midrule\noalign{}
\endhead
\bottomrule\noalign{}
\endlastfoot
\textbf{CS50P (Harvard)} & Full course lectures, free on YouTube and edX
& Closed captions, transcripts; very high production quality \\
\textbf{Corey Schafer Python tutorials} (YouTube) & In-depth concept
tutorials & Closed captions available \\
\textbf{Tech With Tim Python} (YouTube) & Project-based tutorials &
Closed captions available \\
\textbf{Sentdex Python tutorials} (YouTube) & Large series from beginner
to advanced & Captions available \\
\textbf{Python.org Getting Started} & Short official tutorials & Plain
text, no video \\
\end{longtable}
}

\subsubsection{Tools Specifically for Accessibility in Programming
Education}\label{tools-specifically-for-accessibility-in-programming-education}

{\def\LTcaptype{none} % do not increment counter
\begin{longtable}[]{@{}
  >{\raggedright\arraybackslash}p{(\linewidth - 4\tabcolsep) * \real{0.2329}}
  >{\raggedright\arraybackslash}p{(\linewidth - 4\tabcolsep) * \real{0.3151}}
  >{\raggedright\arraybackslash}p{(\linewidth - 4\tabcolsep) * \real{0.4521}}@{}}
\toprule\noalign{}
\begin{minipage}[b]{\linewidth}\raggedright
Tool
\end{minipage} & \begin{minipage}[b]{\linewidth}\raggedright
Purpose
\end{minipage} & \begin{minipage}[b]{\linewidth}\raggedright
Notes
\end{minipage} \\
\midrule\noalign{}
\endhead
\bottomrule\noalign{}
\endlastfoot
\textbf{Python Tutor} (pythontutor.com) & Step-by-step code execution
with state display & Text-based state panel is screen-reader
accessible \\
\textbf{Mu Editor} & Beginner-focused Python IDE & Simpler than VS Code;
reasonable screen reader support \\
\textbf{Thonny} & Beginner IDE with step debugger & Variable viewer is
navigable via keyboard \\
\textbf{NVDA + Remote} & Remote screen sharing for teacher support &
Student's screen reader on their machine; teacher can see output \\
\textbf{GitHub Codespaces} & Cloud VS Code environment & Same
accessibility as local VS Code; useful when local setup fails \\
\textbf{Braille Institute Resources} & General adaptive technology
guidance & brailleinstitute.org \\
\textbf{APH (American Printing House)} & Educational materials for blind
students & aph.org \\
\end{longtable}
}

\subsubsection{Communities and Professional
Organizations}\label{communities-and-professional-organizations}

{\def\LTcaptype{none} % do not increment counter
\begin{longtable}[]{@{}
  >{\raggedright\arraybackslash}p{(\linewidth - 2\tabcolsep) * \real{0.5168}}
  >{\raggedright\arraybackslash}p{(\linewidth - 2\tabcolsep) * \real{0.4832}}@{}}
\toprule\noalign{}
\begin{minipage}[b]{\linewidth}\raggedright
Organization
\end{minipage} & \begin{minipage}[b]{\linewidth}\raggedright
Relevance
\end{minipage} \\
\midrule\noalign{}
\endhead
\bottomrule\noalign{}
\endlastfoot
\textbf{National Federation of the Blind (NFB) --- Science \&
Engineering Division} & Blind scientists and engineers; mentorship
opportunities \\
\textbf{American Council of the Blind (ACB) --- Technology Interest
Group} & Accessible technology advocacy and resources \\
\textbf{CSUN Assistive Technology Conference} & Annual conference; many
papers on computing education for blind students \\
\textbf{SIGACCESS} (ACM SIG on Computers and Accessibility) & Academic
research on accessible computing \\
\textbf{Ability Central} & Technology resources for people with
disabilities \\
\textbf{Hadley Institute for the Blind} & Professional development for
educators of blind students \\
\textbf{Perkins School for the Blind} & Resources including STEM
education for blind students \\
\textbf{Carroll Center for the Blind} & Technology training for blind
adults; may have curriculum resources \\
\textbf{Stack Overflow Blind/VI Community} & Informal community; blind
developers helping each other \\
\end{longtable}
}

\begin{center}\rule{0.5\linewidth}{0.5pt}\end{center}

\subsection{21. Quick Reference Sheets}\label{quick-reference-sheets}

\subsubsection{Error Message Reading
Protocol}\label{error-message-reading-protocol}

When an error occurs, guide the student through this sequence:

\begin{verbatim}
1. Navigate to terminal output (Ctrl+` or click panel)
2. Read from the BOTTOM of the error output upward
3. Note the ERROR TYPE (last line): SyntaxError, NameError, TypeError, etc.
4. Note the ERROR MESSAGE: specific description of what went wrong
5. Note the LINE NUMBER in the traceback
6. Jump to that line with Ctrl+G -> [line number] -> Enter
7. Read the line and the lines immediately before and after it
8. Categorize the error (see table below)
9. Fix and re-run
\end{verbatim}

\subsubsection{Error Type Quick
Reference}\label{error-type-quick-reference}

{\def\LTcaptype{none} % do not increment counter
\begin{longtable}[]{@{}
  >{\raggedright\arraybackslash}p{(\linewidth - 4\tabcolsep) * \real{0.1387}}
  >{\raggedright\arraybackslash}p{(\linewidth - 4\tabcolsep) * \real{0.3650}}
  >{\raggedright\arraybackslash}p{(\linewidth - 4\tabcolsep) * \real{0.4964}}@{}}
\toprule\noalign{}
\begin{minipage}[b]{\linewidth}\raggedright
Error Type
\end{minipage} & \begin{minipage}[b]{\linewidth}\raggedright
Common Cause
\end{minipage} & \begin{minipage}[b]{\linewidth}\raggedright
First Thing to Check
\end{minipage} \\
\midrule\noalign{}
\endhead
\bottomrule\noalign{}
\endlastfoot
\texttt{SyntaxError} & Missing \texttt{:}, \texttt{)}, \texttt{"},
\texttt{{]}}, or \texttt{=} instead of \texttt{==} & The exact column
number in the error message \\
\texttt{IndentationError} & Inconsistent spaces/tabs & Line indentation;
use \texttt{cat\ -A} or \texttt{show\ whitespace} in VS Code \\
\texttt{NameError} & Variable not defined, or misspelled & Check
spelling; check if variable was defined before use \\
\texttt{TypeError} & Wrong data type & Use \texttt{type()} on each
variable involved in the operation \\
\texttt{ValueError} & Right type, wrong value & Check what you're
converting: \texttt{int(\textquotesingle{}hello\textquotesingle{})}
fails; \texttt{int(\textquotesingle{}5\textquotesingle{})} works \\
\texttt{IndexError} & List index out of range & Print \texttt{len(list)}
and check what index you're trying to use \\
\texttt{KeyError} & Dictionary key doesn't exist & Use \texttt{.get()}
or check with \texttt{if\ key\ in\ dict} first \\
\texttt{FileNotFoundError} & File doesn't exist at that path & Print
\texttt{os.getcwd()} and verify the file path \\
\texttt{AttributeError} & Called a method that doesn't exist & Check
spelling and whether the object is the type you think it is \\
\texttt{ZeroDivisionError} & Dividing by zero & Add a check before the
division: \texttt{if\ denominator\ !=\ 0} \\
\texttt{RecursionError} & Infinite recursion & Verify base case exists
and is reachable \\
\end{longtable}
}

\subsubsection{Accessible Output
Patterns}\label{accessible-output-patterns}

\begin{Shaded}
\begin{Highlighting}[]
\CommentTok{\# Instead of: print(my\_list)}
\CommentTok{\# Use:}
\BuiltInTok{print}\NormalTok{(}\StringTok{"Items in collection:"}\NormalTok{)}
\ControlFlowTok{for}\NormalTok{ i, item }\KeywordTok{in} \BuiltInTok{enumerate}\NormalTok{(my\_list, }\DecValTok{1}\NormalTok{):}
    \BuiltInTok{print}\NormalTok{(}\SpecialStringTok{f"  }\SpecialCharTok{\{}\NormalTok{i}\SpecialCharTok{\}}\SpecialStringTok{. }\SpecialCharTok{\{}\NormalTok{item}\SpecialCharTok{\}}\SpecialStringTok{"}\NormalTok{)}
\BuiltInTok{print}\NormalTok{(}\SpecialStringTok{f"Total: }\SpecialCharTok{\{}\BuiltInTok{len}\NormalTok{(my\_list)}\SpecialCharTok{\}}\SpecialStringTok{ items"}\NormalTok{)}

\CommentTok{\# Instead of: print(my\_dict)}
\CommentTok{\# Use:}
\BuiltInTok{print}\NormalTok{(}\StringTok{"Current settings:"}\NormalTok{)}
\ControlFlowTok{for}\NormalTok{ key, value }\KeywordTok{in}\NormalTok{ my\_dict.items():}
    \BuiltInTok{print}\NormalTok{(}\SpecialStringTok{f"  }\SpecialCharTok{\{}\NormalTok{key}\SpecialCharTok{\}}\SpecialStringTok{: }\SpecialCharTok{\{}\NormalTok{value}\SpecialCharTok{\}}\SpecialStringTok{"}\NormalTok{)}

\CommentTok{\# Menu display:}
\BuiltInTok{print}\NormalTok{(}\StringTok{"}\CharTok{\textbackslash{}n}\StringTok{{-}{-}{-} Main Menu {-}{-}{-}"}\NormalTok{)}
\BuiltInTok{print}\NormalTok{(}\StringTok{"1. Add item"}\NormalTok{)}
\BuiltInTok{print}\NormalTok{(}\StringTok{"2. View all items"}\NormalTok{)}
\BuiltInTok{print}\NormalTok{(}\StringTok{"3. Search"}\NormalTok{)}
\BuiltInTok{print}\NormalTok{(}\StringTok{"4. Delete item"}\NormalTok{)}
\BuiltInTok{print}\NormalTok{(}\StringTok{"5. Save and quit"}\NormalTok{)}
\BuiltInTok{print}\NormalTok{(}\StringTok{"{-}{-}{-}{-}{-}{-}{-}{-}{-}{-}{-}{-}{-}{-}{-}{-}{-}"}\NormalTok{)}
\NormalTok{choice }\OperatorTok{=} \BuiltInTok{input}\NormalTok{(}\StringTok{"Enter option number (1{-}5): "}\NormalTok{)}

\CommentTok{\# Confirmation messages:}
\BuiltInTok{print}\NormalTok{(}\SpecialStringTok{f"Contact \textquotesingle{}}\SpecialCharTok{\{}\NormalTok{name}\SpecialCharTok{\}}\SpecialStringTok{\textquotesingle{} added successfully."}\NormalTok{)}
\BuiltInTok{print}\NormalTok{(}\SpecialStringTok{f"Found }\SpecialCharTok{\{}\BuiltInTok{len}\NormalTok{(results)}\SpecialCharTok{\}}\SpecialStringTok{ matching records."}\NormalTok{)}
\BuiltInTok{print}\NormalTok{(}\StringTok{"Data saved to contacts.json."}\NormalTok{)}
\BuiltInTok{print}\NormalTok{(}\StringTok{"No matching records found."}\NormalTok{)}
\BuiltInTok{print}\NormalTok{(}\StringTok{"Error: Name field cannot be empty. Please try again."}\NormalTok{)}
\end{Highlighting}
\end{Shaded}

\subsubsection{VS Code Keyboard Shortcut Reference
Card}\label{vs-code-keyboard-shortcut-reference-card}

\begin{verbatim}
NAVIGATION
  Ctrl+G          Go to specific line number
  Ctrl+Home       Jump to start of file
  Ctrl+End        Jump to end of file
  Ctrl+Shift+O    Go to function/class (symbol picker)
  F12             Go to Definition
  Alt+F12         Peek Definition
  Shift+F12       Find all References
  Ctrl+P          Quick open file by name

ERRORS AND DIAGNOSTICS
  F8              Next error/warning
  Shift+F8        Previous error/warning
  Ctrl+Shift+M    Open Problems panel
  Shift+Alt+F2    Accessible View (for hover/tooltip content)

EDITING
  Shift+Alt+F     Format document (auto-indent)
  Ctrl+/          Toggle comment on selected lines
  Tab             Indent selected lines
  Shift+Tab       Outdent selected lines
  Ctrl+D          Select next occurrence of word
  Ctrl+H          Find and replace

TERMINAL
  Ctrl+`          Toggle integrated terminal
  Ctrl+Shift+`    New terminal
  Up/Down arrows  Navigate command history
  Tab             Auto-complete filenames
  Ctrl+C          Stop running program

SCREEN READER
  Ctrl+Shift+P    Command Palette
  Ctrl+Shift+P -> "Toggle Screen Reader Optimized"
  Shift+Alt+F1    Accessibility help
\end{verbatim}

\newpage

\section{Unit 0: Setting Up Your Accessible Python Development
Environment}\label{unit-0-setting-up-your-accessible-python-development-environment}

\subsection{Overview}\label{overview}

A detailed environment setup is not merely a technical prerequisite; it
empowers students and reduces dependency on external assistance. For
visually impaired students, relying on external help for every setup or
environment issue can be disempowering and slow down the learning
process. This unit provides extremely detailed, screen-reader-friendly
instructions for environment setup.

\textbf{Target Audience:} Students with no prior programming experience,
relying on screen readers for all computer interactions.

\begin{center}\rule{0.5\linewidth}{0.5pt}\end{center}

\subsection{Lesson 0.1: Installing Python and Visual Studio
Code}\label{lesson-0.1-installing-python-and-visual-studio-code}

\subsubsection{Reading References}\label{reading-references}

\begin{itemize}
\tightlist
\item
  Matthes: Chapter 1, pages 3--14
\item
  Matthes: Appendix A, pages 463--468
\end{itemize}

\subsubsection{Installing Python 3}\label{installing-python-3}

Python is not always installed on a computer by default. Download the
latest version from \url{https://www.python.org/downloads/}.

\paragraph{Windows}\label{windows}

\begin{enumerate}
\def\labelenumi{\arabic{enumi}.}
\tightlist
\item
  Locate the downloaded \texttt{.exe} file (usually in your Downloads
  folder). Activate it.
\item
  \textbf{CRUCIAL STEP:} When the installer window appears, check the
  \textbf{``Add Python to PATH''} checkbox. Use Tab to navigate to this
  checkbox and press Spacebar to check it.
\item
  Navigate to ``Install Now'' and activate it.
\item
  Follow any further prompts. Activate ``Close'' when finished.
\end{enumerate}

\paragraph{macOS}\label{macos}

\begin{enumerate}
\def\labelenumi{\arabic{enumi}.}
\tightlist
\item
  Locate the downloaded \texttt{.pkg} file. Activate it.
\item
  Activate ``Continue'' through the steps.
\item
  Agree to the license terms. Choose the installation location (default
  is fine).
\item
  Enter your administrator password when prompted.
\item
  Activate ``Install'' and then ``Close'' when finished.
\end{enumerate}

\paragraph{Linux}\label{linux}

Python 3 is often pre-installed. Verify in a terminal:

\begin{Shaded}
\begin{Highlighting}[]
\ExtensionTok{python3} \AttributeTok{{-}{-}version}
\end{Highlighting}
\end{Shaded}

If not installed (Debian/Ubuntu):

\begin{Shaded}
\begin{Highlighting}[]
\FunctionTok{sudo}\NormalTok{ apt update}
\FunctionTok{sudo}\NormalTok{ apt install python3}
\FunctionTok{sudo}\NormalTok{ apt install python3{-}pip}
\end{Highlighting}
\end{Shaded}

\subsubsection{Verifying Python
Installation}\label{verifying-python-installation}

Open a terminal and type:

\begin{Shaded}
\begin{Highlighting}[]
\ExtensionTok{python3} \AttributeTok{{-}{-}version}
\end{Highlighting}
\end{Shaded}

You should hear the version number (e.g., ``Python 3.10.12''). Also
check pip:

\begin{Shaded}
\begin{Highlighting}[]
\ExtensionTok{pip3} \AttributeTok{{-}{-}version}
\end{Highlighting}
\end{Shaded}

\subsubsection{Installing Visual Studio
Code}\label{installing-visual-studio-code}

\begin{enumerate}
\def\labelenumi{\arabic{enumi}.}
\tightlist
\item
  Navigate to \url{https://code.visualstudio.com/download}.
\item
  Download the installer for your OS.
\end{enumerate}

\paragraph{Windows}\label{windows-1}

\begin{itemize}
\tightlist
\item
  \textbf{CRUCIAL STEP:} On the ``Select Additional Tasks'' screen,
  check \textbf{``Add to PATH''}.
\item
  Activate ``Next'' and then ``Install''.
\end{itemize}

\paragraph{macOS}\label{macos-1}

\begin{itemize}
\tightlist
\item
  Drag the ``Visual Studio Code'' application to your ``Applications''
  folder.
\end{itemize}

\paragraph{Linux (Debian/Ubuntu)}\label{linux-debianubuntu}

\begin{Shaded}
\begin{Highlighting}[]
\BuiltInTok{cd}\NormalTok{ Downloads}
\FunctionTok{sudo}\NormalTok{ dpkg }\AttributeTok{{-}i}\NormalTok{ code\_}\PreprocessorTok{*}\NormalTok{.deb}
\end{Highlighting}
\end{Shaded}

\subsubsection{Verify VS Code}\label{verify-vs-code}

\begin{Shaded}
\begin{Highlighting}[]
\ExtensionTok{code} \AttributeTok{{-}{-}version}
\end{Highlighting}
\end{Shaded}

\subsubsection{Skills Checklist}\label{skills-checklist}

\begin{itemize}
\tightlist
\item[$\square$]
  I can successfully download the Python 3 installer for my operating
  system.
\item[$\square$]
  I can run the Python installer and correctly select the ``Add Python
  to PATH'' option.
\item[$\square$]
  I can successfully download the VS Code installer for my operating
  system.
\item[$\square$]
  I can run the VS Code installer and correctly select the ``Add to
  PATH'' option.
\item[$\square$]
  I can open my computer's system terminal.
\item[$\square$]
  I can verify Python's installation by running
  \texttt{python3\ -\/-version} in the terminal.
\item[$\square$]
  I can verify VS Code's installation by running
  \texttt{code\ -\/-version} in the terminal.
\end{itemize}

\begin{center}\rule{0.5\linewidth}{0.5pt}\end{center}

\subsection{Lesson 0.2: Configuring Your Screen Reader for VS
Code}\label{lesson-0.2-configuring-your-screen-reader-for-vs-code-1}

\subsubsection{Reading References}\label{reading-references-1}

\begin{itemize}
\tightlist
\item
  Matthes: Appendix B, pages 469--476
\item
  Matthes: Appendix C, pages 477--482
\end{itemize}

\subsubsection{General Screen Reader Tips for VS
Code}\label{general-screen-reader-tips-for-vs-code}

\textbf{1. Use Screen Reader Mode}

\begin{itemize}
\tightlist
\item
  Press \texttt{Ctrl\ +\ Shift\ +\ P} to open the Command Palette.
\item
  Type ``Toggle Screen Reader Optimized'' and press Enter.
\end{itemize}

\textbf{2. Accessible View}

\begin{itemize}
\tightlist
\item
  Press \texttt{Shift\ +\ Alt\ +\ F2} (Windows/Linux) or
  \texttt{Shift\ +\ Option\ +\ F2} (macOS).
\item
  Opens a flat, text-based representation of complex UI elements.
\item
  Press \texttt{Esc} to close.
\end{itemize}

\textbf{3. Keyboard-First Navigation}

\begin{itemize}
\tightlist
\item
  \texttt{Tab} and \texttt{Shift+Tab}: Move focus between UI elements.
\item
  Arrow Keys: Navigate within active panes.
\end{itemize}

\textbf{4. Status Bar} The status bar at the bottom provides current
line/column number, file type, active Python environment, and
error/warning counts.

\subsubsection{Screen Reader Indentation
Configuration}\label{screen-reader-indentation-configuration}

Python relies heavily on indentation. Configure your screen reader to
announce indentation changes.

\paragraph{NVDA (Windows)}\label{nvda-windows}

\begin{enumerate}
\def\labelenumi{\arabic{enumi}.}
\tightlist
\item
  Press \texttt{NVDA\ Key\ +\ N} -\textgreater{} Preferences
  -\textgreater{} Settings.
\item
  Find ``Document Formatting''.
\item
  Locate ``Indentation'' and set it to \textbf{``Say all''} or
  \textbf{``Beeps''}.
\item
  Activate ``OK''.
\end{enumerate}

\paragraph{JAWS (Windows)}\label{jaws-windows}

\begin{enumerate}
\def\labelenumi{\arabic{enumi}.}
\tightlist
\item
  Press \texttt{JAWS\ Key\ +\ J} -\textgreater{} Alt+O for Options
  -\textgreater{} B for Basics.
\item
  Find ``Text Processing'' settings.
\item
  Configure indentation announcement.
\end{enumerate}

\paragraph{VoiceOver (macOS)}\label{voiceover-macos}

VoiceOver generally handles indentation well by default, announcing ``4
spaces'' or ``tab'' as you navigate.

\subsubsection{Skills Checklist}\label{skills-checklist-1}

\begin{itemize}
\tightlist
\item[$\square$]
  I can identify and toggle VS Code's ``Screen Reader Optimized Mode''.
\item[$\square$]
  I can open and close the ``Accessible View''.
\item[$\square$]
  I can navigate between VS Code panes using Tab and Shift+Tab.
\item[$\square$]
  I can navigate line by line and character by character within the
  editor.
\item[$\square$]
  I can locate and listen to information in the VS Code Status Bar.
\item[$\square$]
  I can configure my screen reader to announce indentation changes.
\end{itemize}

\begin{center}\rule{0.5\linewidth}{0.5pt}\end{center}

\subsection{Lesson 0.3: Setting Up Your VS Code
Workspace}\label{lesson-0.3-setting-up-your-vs-code-workspace-1}

\subsubsection{Reading References}\label{reading-references-2}

\begin{itemize}
\tightlist
\item
  Matthes: Appendix B, pages 469--476
\item
  Matthes: Chapter 8, pages 149--152
\end{itemize}

\subsubsection{Installing VS Code
Extensions}\label{installing-vs-code-extensions}

\begin{enumerate}
\def\labelenumi{\arabic{enumi}.}
\tightlist
\item
  Open the Extensions View: \texttt{Ctrl\ +\ Shift\ +\ X}
  (Windows/Linux) or \texttt{Command\ +\ Shift\ +\ X} (macOS).
\item
  Search and install \textbf{``Python''} by Microsoft --- the official
  and most important Python extension.
\item
  Search and install \textbf{``Pylance''} --- provides excellent
  IntelliSense and type checking.
\item
  Optional: \textbf{``Code Runner''} --- allows running code files with
  \texttt{Ctrl+Alt+N}.
\end{enumerate}

\subsubsection{Creating Your First Workspace
Folder}\label{creating-your-first-workspace-folder}

\textbf{Windows:} Open File Explorer (\texttt{Windows\ Key\ +\ E}).
Navigate to Documents. Press \texttt{Ctrl\ +\ Shift\ +\ N} to create a
new folder named \texttt{PythonProjects}.

\textbf{macOS:} Open Finder. Navigate to Documents. Press
\texttt{Command\ +\ Shift\ +\ N} to create \texttt{PythonProjects}.

\subsubsection{Opening the Folder in VS
Code}\label{opening-the-folder-in-vs-code}

\textbf{Method 1 (Recommended):} Open terminal, navigate to the folder,
then:

\begin{Shaded}
\begin{Highlighting}[]
\ExtensionTok{code}\NormalTok{ .}
\end{Highlighting}
\end{Shaded}

\textbf{Method 2:} In VS Code, press \texttt{Ctrl\ +\ K,\ Ctrl\ +\ O} to
``Open Folder''.

\subsubsection{VS Code Navigation
Shortcuts}\label{vs-code-navigation-shortcuts}

{\def\LTcaptype{none} % do not increment counter
\begin{longtable}[]{@{}ll@{}}
\toprule\noalign{}
Area & Shortcut \\
\midrule\noalign{}
\endhead
\bottomrule\noalign{}
\endlastfoot
Explorer & \texttt{Ctrl\ +\ Shift\ +\ E} \\
Search & \texttt{Ctrl\ +\ Shift\ +\ F} \\
Extensions & \texttt{Ctrl\ +\ Shift\ +\ X} \\
Toggle Sidebar & \texttt{Ctrl\ +\ B} \\
Command Palette & \texttt{Ctrl\ +\ Shift\ +\ P} \\
Problems View & \texttt{Ctrl\ +\ Shift\ +\ M} \\
Next Error & \texttt{F8} \\
Previous Error & \texttt{Shift\ +\ F8} \\
Integrated Terminal & `Ctrl + `` \\
\end{longtable}
}

\subsubsection{Skills Checklist}\label{skills-checklist-2}

\begin{itemize}
\tightlist
\item[$\square$]
  I can open the Extensions View in VS Code.
\item[$\square$]
  I can search for and install the ``Python'' extension.
\item[$\square$]
  I can search for and install the ``Pylance'' extension.
\item[$\square$]
  I can create a new folder for Python projects.
\item[$\square$]
  I can open a specific folder as a workspace in VS Code.
\item[$\square$]
  I can navigate the Activity Bar, Explorer, Editor, and Panel using
  keyboard shortcuts.
\item[$\square$]
  I can effectively use the Command Palette
  (\texttt{Ctrl\ +\ Shift\ +\ P}).
\end{itemize}

\begin{center}\rule{0.5\linewidth}{0.5pt}\end{center}

\subsection{Lesson 0.4: Python Environments and Terminal
Usage}\label{lesson-0.4-python-environments-and-terminal-usage-1}

\subsubsection{Reading References}\label{reading-references-3}

\begin{itemize}
\tightlist
\item
  Matthes: Chapter 1, pages 3--14
\item
  Matthes: Chapter 8, pages 149--152
\item
  Matthes: Chapter 11, pages 210--211
\end{itemize}

\subsubsection{Creating Your First Python
File}\label{creating-your-first-python-file}

\begin{enumerate}
\def\labelenumi{\arabic{enumi}.}
\tightlist
\item
  Ensure Explorer is focused: \texttt{Ctrl\ +\ Shift\ +\ E}.
\item
  Create new file: \texttt{Ctrl\ +\ N}.
\item
  Save: \texttt{Ctrl\ +\ S}. Name it \texttt{hello.py}.
\item
  Type:
\end{enumerate}

\begin{Shaded}
\begin{Highlighting}[]
\BuiltInTok{print}\NormalTok{(}\StringTok{"Hello from VS Code!"}\NormalTok{)}
\end{Highlighting}
\end{Shaded}

\subsubsection{Running Python Scripts}\label{running-python-scripts}

\textbf{From System Terminal:}

\begin{Shaded}
\begin{Highlighting}[]
\BuiltInTok{cd}\NormalTok{ \textasciitilde{}/Documents/PythonProjects}
\ExtensionTok{python3}\NormalTok{ hello.py}
\end{Highlighting}
\end{Shaded}

\textbf{From VS Code Integrated Terminal:}

\begin{enumerate}
\def\labelenumi{\arabic{enumi}.}
\tightlist
\item
  Open terminal: `Ctrl + ``
\item
  The terminal opens in your project folder automatically.
\item
  Run: \texttt{python3\ hello.py}
\end{enumerate}

\subsubsection{Python Virtual
Environments}\label{python-virtual-environments}

Virtual environments keep project dependencies separate, preventing
conflicts.

\textbf{Create a virtual environment:}

\begin{Shaded}
\begin{Highlighting}[]
\ExtensionTok{python3} \AttributeTok{{-}m}\NormalTok{ venv .venv}
\end{Highlighting}
\end{Shaded}

\textbf{Activate:}

\begin{itemize}
\tightlist
\item
  Windows (Command Prompt):
  \texttt{.venv\textbackslash{}Scripts\textbackslash{}activate.bat}
\item
  Windows (PowerShell):
  \texttt{.venv\textbackslash{}Scripts\textbackslash{}Activate.ps1}
\item
  macOS/Linux: \texttt{source\ ./.venv/bin/activate}
\end{itemize}

You will hear \texttt{(.venv)} at the beginning of your prompt when
active.

\textbf{Deactivate:}

\begin{Shaded}
\begin{Highlighting}[]
\ExtensionTok{deactivate}
\end{Highlighting}
\end{Shaded}

\textbf{Select Interpreter in VS Code:}

\begin{enumerate}
\def\labelenumi{\arabic{enumi}.}
\tightlist
\item
  Press \texttt{Ctrl\ +\ Shift\ +\ P}.
\item
  Type ``Python: Select Interpreter'' and press Enter.
\item
  Select the one that says ``Python 3.x.x (.venv)''.
\end{enumerate}

\textbf{Install packages:}

\begin{Shaded}
\begin{Highlighting}[]
\ExtensionTok{pip}\NormalTok{ install requests}
\ExtensionTok{pip}\NormalTok{ list}
\end{Highlighting}
\end{Shaded}

\subsubsection{Skills Checklist}\label{skills-checklist-3}

\begin{itemize}
\tightlist
\item[$\square$]
  I can create a new Python file (.py) in VS Code.
\item[$\square$]
  I can write a simple \texttt{print()} statement.
\item[$\square$]
  I can run a Python script from the system terminal.
\item[$\square$]
  I can run a Python script from the integrated terminal in VS Code.
\item[$\square$]
  I understand why virtual environments are useful.
\item[$\square$]
  I can create, activate, and deactivate a virtual environment.
\item[$\square$]
  I can select the correct Python interpreter in VS Code.
\item[$\square$]
  I can install Python packages using pip.
\end{itemize}

\begin{center}\rule{0.5\linewidth}{0.5pt}\end{center}

\subsection{Lesson 0.5: General Accessibility Tools and
Strategies}\label{lesson-0.5-general-accessibility-tools-and-strategies}

\subsubsection{VS Code Accessibility
Features}\label{vs-code-accessibility-features}

\begin{itemize}
\tightlist
\item
  \textbf{Accessible View} (\texttt{Shift+Alt+F2}): Reveals hover
  information and error details.
\item
  \textbf{Focus on Errors} (\texttt{F8} / \texttt{Shift+F8}): Jump
  between errors without visual scanning.
\item
  \textbf{Go to Definition} (\texttt{F12}): Navigate to where a function
  or variable is defined.
\item
  \textbf{Go to Reference} (\texttt{Shift+F12}): Find all usages of a
  symbol.
\item
  \textbf{Go to Line/Column} (\texttt{Ctrl+G}): Jump to a specific line
  number.
\end{itemize}

\subsubsection{Python Documentation}\label{python-documentation}

The official Python documentation is well-structured and accessible:
\url{https://docs.python.org/3/}

\subsubsection{Unit 0 Completion
Checklist}\label{unit-0-completion-checklist}

\begin{itemize}
\tightlist
\item[$\square$]
  Python 3 is installed and \texttt{python3\ -\/-version} works in your
  system terminal.
\item[$\square$]
  VS Code is installed and \texttt{code\ -\/-version} works in your
  system terminal.
\item[$\square$]
  Your screen reader is configured to announce indentation changes in VS
  Code.
\item[$\square$]
  You can open your PythonProjects folder in VS Code.
\item[$\square$]
  You have installed the official ``Python'' VS Code extension.
\item[$\square$]
  You can create a new .py file, write \texttt{print("Hello!")}, and
  save it.
\item[$\square$]
  You can run your hello.py file from both terminals.
\item[$\square$]
  You can create and activate a virtual environment (.venv).
\item[$\square$]
  You can select the virtual environment interpreter in VS Code.
\item[$\square$]
  You can install a package using pip.
\item[$\square$]
  You know how to use \texttt{Ctrl\ +\ Shift\ +\ P} (Command Palette)
  and \texttt{F8} (Go to next error).
\end{itemize}

\textbf{Congratulations! You now have a powerful and accessible
development environment ready for your Python journey.}

\newpage

\section{Unit 1: Foundations of Computing and Getting Started with
Python}\label{unit-1-foundations-of-computing-and-getting-started-with-python}

\subsection{Overview}\label{overview-1}

This unit covers basic computer science concepts, setting up the Python
environment, and writing and executing simple Python programs.

\textbf{Objectives:}

\begin{itemize}
\tightlist
\item
  Understand what computer science is and how hardware relates to
  software.
\item
  Write and run your first Python program.
\item
  Use variables to store and manipulate simple data types.
\end{itemize}

\begin{center}\rule{0.5\linewidth}{0.5pt}\end{center}

\subsection{Lesson 1.1: Welcome to Computer
Science}\label{lesson-1.1-welcome-to-computer-science-1}

\subsubsection{Reading References}\label{reading-references-4}

\begin{itemize}
\tightlist
\item
  Zelle: Chapter 1, pages 1--7
\item
  Matthes: Chapter 1, pages 3--14
\end{itemize}

\subsubsection{What Is Computer
Science?}\label{what-is-computer-science}

A computer is a \textbf{universal machine} --- it can be programmed to
solve any problem that can be broken down into a series of precise,
logical steps. The true power of a computer doesn't come from its
hardware but from the software we create.

\textbf{Computer science} is the study of problem-solving, not just
programming. It's an interdisciplinary field involving breaking down
complex problems, designing algorithms, and implementing efficient
solutions.

\subsubsection{Hardware vs.~Software}\label{hardware-vs.-software}

{\def\LTcaptype{none} % do not increment counter
\begin{longtable}[]{@{}
  >{\raggedright\arraybackslash}p{(\linewidth - 2\tabcolsep) * \real{0.1860}}
  >{\raggedright\arraybackslash}p{(\linewidth - 2\tabcolsep) * \real{0.8140}}@{}}
\toprule\noalign{}
\begin{minipage}[b]{\linewidth}\raggedright
Component
\end{minipage} & \begin{minipage}[b]{\linewidth}\raggedright
Description
\end{minipage} \\
\midrule\noalign{}
\endhead
\bottomrule\noalign{}
\endlastfoot
\textbf{Hardware} & Physical components you can touch: keyboard, screen,
processor, memory \\
\textbf{Software} & Instructions that direct hardware: programs,
operating systems, apps \\
\textbf{CPU} & The processor that executes instructions \\
\textbf{Memory (RAM)} & Temporary storage for active programs and
data \\
\textbf{Storage} & Permanent storage for files \\
\end{longtable}
}

\textbf{Analogy:} Think of a kitchen. The stove, refrigerator, and pots
are \textbf{hardware}. A cookie recipe is \textbf{software} --- it
provides precise steps for the ``hardware'' (the cook) to follow.

\subsubsection{Why Python?}\label{why-python}

Python is an ideal language for beginners because of its:

\begin{itemize}
\tightlist
\item
  \textbf{Readability} --- English-like syntax makes code easy to
  understand.
\item
  \textbf{Simplicity} --- Less ``boilerplate'' code than other
  languages.
\item
  \textbf{Versatility} --- Used in web development, data science, AI,
  and more.
\end{itemize}

\subsubsection{The Problem-Solving
Approach}\label{the-problem-solving-approach}

\begin{enumerate}
\def\labelenumi{\arabic{enumi}.}
\tightlist
\item
  Understand the problem.
\item
  Break it down into smaller pieces.
\item
  Design a step-by-step solution (an algorithm).
\item
  Implement the solution in Python.
\item
  Test and refine your code.
\end{enumerate}

\subsubsection{Activities}\label{activities}

\textbf{Activity 1.1.1: ``Human Computer'' Role-Play}

\begin{itemize}
\tightlist
\item
  In pairs, one student is the ``computer'' and the other is the
  ``programmer.''
\item
  The programmer provides simple, precise instructions for a common task
  (e.g., ``make a peanut butter and jelly sandwich'').
\item
  The ``computer'' follows only what is literally said.
\item
  Purpose: Illustrate the need for precise, unambiguous instructions.
\end{itemize}

\textbf{Activity 1.1.2: Identify Hardware and Software}

\begin{itemize}
\tightlist
\item
  List five items you use daily.
\item
  For each, identify whether it is hardware or software and explain your
  choice.
\end{itemize}

\subsubsection{Skills Checklist}\label{skills-checklist-4}

\begin{itemize}
\tightlist
\item[$\square$]
  I can define what computer science is.
\item[$\square$]
  I can distinguish between computer hardware and software.
\item[$\square$]
  I can explain what a computer program is.
\item[$\square$]
  I can identify Python as a programming language.
\end{itemize}

\begin{center}\rule{0.5\linewidth}{0.5pt}\end{center}

\subsection{Lesson 1.2: Your First Python Program --- ``Hello
World!''}\label{lesson-1.2-your-first-python-program-hello-world}

\subsubsection{Reading References}\label{reading-references-5}

\begin{itemize}
\tightlist
\item
  Zelle: Chapter 1, pages 7--13
\item
  Matthes: Chapter 1, pages 3--14
\end{itemize}

\subsubsection{\texorpdfstring{The \texttt{print()}
Function}{The print() Function}}\label{the-print-function}

The \texttt{print()} function is Python's primary method for displaying
output to the user.

\begin{Shaded}
\begin{Highlighting}[]
\BuiltInTok{print}\NormalTok{(}\StringTok{"Hello, World!"}\NormalTok{)}
\end{Highlighting}
\end{Shaded}

\textbf{Expected output:}

\begin{verbatim}
Hello, World!
\end{verbatim}

\subsubsection{Key Concepts}\label{key-concepts}

\begin{itemize}
\tightlist
\item
  \textbf{Strings} are text data enclosed in quotation marks (single or
  double).
\item
  \textbf{Comments} start with \texttt{\#} and are ignored by Python.
\item
  Python executes code \textbf{line by line}, from top to bottom.
\end{itemize}

\subsubsection{Examples}\label{examples}

\textbf{Example 1.2.1: Simple Greeting}

\begin{Shaded}
\begin{Highlighting}[]
\CommentTok{\# This program prints a classic greeting.}
\BuiltInTok{print}\NormalTok{(}\StringTok{"Hello, World!"}\NormalTok{)}
\end{Highlighting}
\end{Shaded}

\textbf{Example 1.2.2: Personalized Message}

\begin{Shaded}
\begin{Highlighting}[]
\CommentTok{\# This program prints a message with your name.}
\BuiltInTok{print}\NormalTok{(}\StringTok{"Hello, Alice!"}\NormalTok{)}
\BuiltInTok{print}\NormalTok{(}\StringTok{"Welcome to Python programming!"}\NormalTok{)}
\end{Highlighting}
\end{Shaded}

Expected output:

\begin{verbatim}
Hello, Alice!
Welcome to Python programming!
\end{verbatim}

\subsubsection{Common Mistakes}\label{common-mistakes}

\begin{Shaded}
\begin{Highlighting}[]
\BuiltInTok{print}\NormalTok{(}\StringTok{"Hello, World"}        \CommentTok{\# Missing closing parenthesis {-}{-}{-} SyntaxError}
\BuiltInTok{print} \StringTok{"Hello"}               \CommentTok{\# Python 2 style {-}{-}{-} SyntaxError in Python 3}
\NormalTok{Print(}\StringTok{"Hello"}\NormalTok{)              }\CommentTok{\# Capital P {-}{-}{-} NameError}
\end{Highlighting}
\end{Shaded}

\subsubsection{Activities}\label{activities-1}

\textbf{Activity 1.2.1: My First Message}

\begin{enumerate}
\def\labelenumi{\arabic{enumi}.}
\item
  Open VS Code and confirm the PythonProjects folder is open.
\item
  Create a new file as \texttt{my\_message.py}.
\item
  Type:

\begin{Shaded}
\begin{Highlighting}[]
\BuiltInTok{print}\NormalTok{(}\StringTok{"My favorite color is blue."}\NormalTok{)}
\BuiltInTok{print}\NormalTok{(}\StringTok{"My favorite animal is a dog."}\NormalTok{)}
\end{Highlighting}
\end{Shaded}
\item
  Save and run the file.
\end{enumerate}

\textbf{Activity 1.2.2: Exploring Print Statements}

\begin{enumerate}
\def\labelenumi{\arabic{enumi}.}
\tightlist
\item
  Experiment with:

  \begin{itemize}
  \tightlist
  \item
    Printing numbers: \texttt{print(123)}
  \item
    Printing sentences with punctuation.
  \item
    Omitting quotes: \texttt{print(Hello)} --- expect a
    \texttt{NameError}.
  \end{itemize}
\end{enumerate}

\subsubsection{Skills Checklist}\label{skills-checklist-5}

\begin{itemize}
\tightlist
\item[$\square$]
  I can write a \texttt{print()} statement in Python.
\item[$\square$]
  I can save a Python file with the \texttt{.py} extension.
\item[$\square$]
  I can execute a Python script from the terminal.
\item[$\square$]
  I understand that the output of \texttt{print()} appears in the
  terminal.
\end{itemize}

\begin{center}\rule{0.5\linewidth}{0.5pt}\end{center}

\subsection{Lesson 1.3: Variables and Simple Data
Types}\label{lesson-1.3-variables-and-simple-data-types-1}

\subsubsection{Reading References}\label{reading-references-6}

\begin{itemize}
\tightlist
\item
  Zelle: Chapter 2, pages 24--26
\item
  Matthes: Chapter 2, pages 15--32
\end{itemize}

\subsubsection{What Is a Variable?}\label{what-is-a-variable}

A \textbf{variable} is a name that refers to a value stored in the
computer's memory. The \textbf{assignment operator} \texttt{=} creates a
connection between a name and a value.

\begin{Shaded}
\begin{Highlighting}[]
\NormalTok{age }\OperatorTok{=} \DecValTok{25}
\NormalTok{name }\OperatorTok{=} \StringTok{"Alice"}
\NormalTok{height }\OperatorTok{=} \FloatTok{5.6}
\end{Highlighting}
\end{Shaded}

\textbf{Rules for variable names:}

\begin{itemize}
\tightlist
\item
  Must start with a letter or underscore
\item
  Can only contain letters, digits, and underscores
\item
  Are case-sensitive (\texttt{age} is different from \texttt{Age})
\item
  Cannot be a Python keyword (\texttt{if}, \texttt{for}, \texttt{while},
  etc.)
\end{itemize}

\textbf{Best practices:}

\begin{itemize}
\tightlist
\item
  Use descriptive names: \texttt{student\_count}, not \texttt{sc}
\item
  Use lowercase with underscores: \texttt{first\_name}
\end{itemize}

\subsubsection{Essential Data Types}\label{essential-data-types}

{\def\LTcaptype{none} % do not increment counter
\begin{longtable}[]{@{}lll@{}}
\toprule\noalign{}
Type & Description & Examples \\
\midrule\noalign{}
\endhead
\bottomrule\noalign{}
\endlastfoot
\texttt{str} (String) & Text data & \texttt{"hello"},
\texttt{"Python"} \\
\texttt{int} (Integer) & Whole numbers & \texttt{10}, \texttt{-5},
\texttt{0} \\
\texttt{float} (Float) & Numbers with decimals & \texttt{3.14},
\texttt{9.99} \\
\texttt{bool} (Boolean) & True or False & \texttt{True},
\texttt{False} \\
\end{longtable}
}

\subsubsection{Examples}\label{examples-1}

\textbf{Example 1.3.1: Storing and Printing a String}

\begin{Shaded}
\begin{Highlighting}[]
\NormalTok{student\_name }\OperatorTok{=} \StringTok{"Charlie"}
\BuiltInTok{print}\NormalTok{(}\StringTok{"Hello, "} \OperatorTok{+}\NormalTok{ student\_name }\OperatorTok{+} \StringTok{"!"}\NormalTok{)}
\end{Highlighting}
\end{Shaded}

Output:

\begin{verbatim}
Hello, Charlie!
\end{verbatim}

\textbf{Example 1.3.2: Storing and Printing Numbers}

\begin{Shaded}
\begin{Highlighting}[]
\NormalTok{apple\_count }\OperatorTok{=} \DecValTok{5}
\NormalTok{orange\_count }\OperatorTok{=} \DecValTok{3}
\NormalTok{total\_fruit }\OperatorTok{=}\NormalTok{ apple\_count }\OperatorTok{+}\NormalTok{ orange\_count}
\BuiltInTok{print}\NormalTok{(}\StringTok{"You have"}\NormalTok{, total\_fruit, }\StringTok{"fruits in total."}\NormalTok{)}

\NormalTok{price\_per\_item }\OperatorTok{=} \FloatTok{2.50}
\NormalTok{quantity }\OperatorTok{=} \DecValTok{4}
\NormalTok{total\_cost }\OperatorTok{=}\NormalTok{ price\_per\_item }\OperatorTok{*}\NormalTok{ quantity}
\BuiltInTok{print}\NormalTok{(}\StringTok{"The total cost is:"}\NormalTok{, total\_cost)}
\end{Highlighting}
\end{Shaded}

Output:

\begin{verbatim}
You have 8 fruits in total.
The total cost is: 10.0
\end{verbatim}

\subsubsection{Common Errors}\label{common-errors}

\begin{Shaded}
\begin{Highlighting}[]
\DecValTok{2}\ErrorTok{nd\_number} \OperatorTok{=} \DecValTok{10}         \CommentTok{\# Invalid variable name {-}{-}{-} starts with digit}
\BuiltInTok{print}\NormalTok{(}\StringTok{"Result: "} \OperatorTok{+} \DecValTok{42}\NormalTok{)  }\CommentTok{\# TypeError {-}{-}{-} can\textquotesingle{}t concatenate str and int}
\end{Highlighting}
\end{Shaded}

\subsubsection{Activities}\label{activities-2}

\textbf{Activity 1.3.1: Personal Information Variables}

\begin{Shaded}
\begin{Highlighting}[]
\NormalTok{my\_name }\OperatorTok{=} \StringTok{"Jamie"}
\NormalTok{favorite\_number }\OperatorTok{=} \DecValTok{7}
\NormalTok{height\_in\_meters }\OperatorTok{=} \FloatTok{1.75}
\NormalTok{favorite\_food }\OperatorTok{=} \StringTok{"pizza"}

\BuiltInTok{print}\NormalTok{(}\StringTok{"My name is:"}\NormalTok{, my\_name)}
\BuiltInTok{print}\NormalTok{(}\StringTok{"My favorite number is:"}\NormalTok{, favorite\_number)}
\BuiltInTok{print}\NormalTok{(}\StringTok{"My height is:"}\NormalTok{, height\_in\_meters, }\StringTok{"meters"}\NormalTok{)}
\BuiltInTok{print}\NormalTok{(}\StringTok{"My favorite food is:"}\NormalTok{, favorite\_food)}
\end{Highlighting}
\end{Shaded}

\textbf{Activity 1.3.2: Simple Math Challenge}

\begin{Shaded}
\begin{Highlighting}[]
\NormalTok{num1 }\OperatorTok{=} \DecValTok{15}
\NormalTok{num2 }\OperatorTok{=} \DecValTok{7}
\NormalTok{sum\_result }\OperatorTok{=}\NormalTok{ num1 }\OperatorTok{+}\NormalTok{ num2}
\NormalTok{difference }\OperatorTok{=}\NormalTok{ num1 }\OperatorTok{{-}}\NormalTok{ num2}
\NormalTok{product }\OperatorTok{=}\NormalTok{ num1 }\OperatorTok{*}\NormalTok{ num2}
\NormalTok{quotient }\OperatorTok{=}\NormalTok{ num1 }\OperatorTok{/}\NormalTok{ num2}

\BuiltInTok{print}\NormalTok{(}\StringTok{"Sum:"}\NormalTok{, sum\_result)}
\BuiltInTok{print}\NormalTok{(}\StringTok{"Difference:"}\NormalTok{, difference)}
\BuiltInTok{print}\NormalTok{(}\StringTok{"Product:"}\NormalTok{, product)}
\BuiltInTok{print}\NormalTok{(}\StringTok{"Quotient:"}\NormalTok{, quotient)}
\end{Highlighting}
\end{Shaded}

\subsubsection{Key Table: Python Data Types and
Operations}\label{key-table-python-data-types-and-operations}

{\def\LTcaptype{none} % do not increment counter
\begin{longtable}[]{@{}
  >{\raggedright\arraybackslash}p{(\linewidth - 6\tabcolsep) * \real{0.0978}}
  >{\raggedright\arraybackslash}p{(\linewidth - 6\tabcolsep) * \real{0.2391}}
  >{\raggedright\arraybackslash}p{(\linewidth - 6\tabcolsep) * \real{0.2283}}
  >{\raggedright\arraybackslash}p{(\linewidth - 6\tabcolsep) * \real{0.4348}}@{}}
\toprule\noalign{}
\begin{minipage}[b]{\linewidth}\raggedright
Data Type
\end{minipage} & \begin{minipage}[b]{\linewidth}\raggedright
Description
\end{minipage} & \begin{minipage}[b]{\linewidth}\raggedright
Example Value
\end{minipage} & \begin{minipage}[b]{\linewidth}\raggedright
Common Operations
\end{minipage} \\
\midrule\noalign{}
\endhead
\bottomrule\noalign{}
\endlastfoot
\texttt{str} & Text, words, sentences & \texttt{"Hello"},
\texttt{"Python"} & \texttt{+} (concatenation), \texttt{len()},
\texttt{.upper()} \\
\texttt{int} & Whole numbers & \texttt{10}, \texttt{-5}, \texttt{0} &
\texttt{+}, \texttt{-}, \texttt{*}, \texttt{/}, \texttt{\%} \\
\texttt{float} & Numbers with decimals & \texttt{3.14}, \texttt{9.99},
\texttt{0.0} & \texttt{+}, \texttt{-}, \texttt{*}, \texttt{/} \\
\end{longtable}
}

\subsubsection{Skills Checklist}\label{skills-checklist-6}

\begin{itemize}
\tightlist
\item[$\square$]
  I can define what a variable is in programming.
\item[$\square$]
  I can create a valid variable name in Python.
\item[$\square$]
  I can assign string, integer, and float values to variables.
\item[$\square$]
  I can use \texttt{print()} to display the value of a variable.
\item[$\square$]
  I can perform basic arithmetic operations with numbers in Python.
\end{itemize}

\newpage

\section{Unit 2: Writing Simple Programs and Numeric
Computations}\label{unit-2-writing-simple-programs-and-numeric-computations}

\subsection{Overview}\label{overview-2}

This unit explores the software development process, more complex
numeric operations, and how to obtain and use input from users.

\textbf{Objectives:}

\begin{itemize}
\tightlist
\item
  Understand the phases of software development.
\item
  Work with different numeric data types (integers, floats).
\item
  Perform type conversions.
\item
  Use mathematical operators and the math library.
\item
  Get user input with \texttt{input()}.
\end{itemize}

\begin{center}\rule{0.5\linewidth}{0.5pt}\end{center}

\subsection{Lesson 2.1: The Software Development
Process}\label{lesson-2.1-the-software-development-process-1}

\subsubsection{Reading References}\label{reading-references-7}

\begin{itemize}
\tightlist
\item
  Zelle: Chapter 2, pages 21--22 and pages 35--36
\end{itemize}

\subsubsection{The Three-Phase Development
Process}\label{the-three-phase-development-process}

\textbf{1. Design (Plan first)} Before writing any code, fully
understand the problem:

\begin{itemize}
\tightlist
\item
  Identify \textbf{inputs} (information the program needs) and
  \textbf{outputs} (results it must produce).
\item
  Write an \textbf{algorithm} --- a step-by-step plan in plain English
  (\textbf{pseudocode}).
\item
  Consider \textbf{edge cases} --- unusual inputs that might break your
  program.
\end{itemize}

\textbf{2. Implement (Write the code)}

\begin{itemize}
\tightlist
\item
  Translate your design into Python syntax.
\item
  Code \textbf{incrementally} --- build in small pieces and test as you
  go.
\item
  Follow coding standards and use descriptive variable names.
\end{itemize}

\textbf{3. Test and Debug (Verify it works)}

\begin{itemize}
\tightlist
\item
  Run with a variety of inputs, including normal values and edge cases.
\item
  Read error messages carefully --- they tell you the type of error and
  the line number.
\item
  Use \texttt{print()} statements to check variable values during
  debugging.
\end{itemize}

\subsubsection{Why This Process Matters}\label{why-this-process-matters}

\begin{itemize}
\tightlist
\item
  \textbf{Reduces frustration} --- Having a plan prevents getting stuck.
\item
  \textbf{Improves code quality} --- A thoughtful design leads to
  cleaner code.
\item
  \textbf{Simplifies debugging} --- Test each component in isolation.
\item
  \textbf{Saves time} --- Good design prevents time-consuming rewrites.
\end{itemize}

\subsubsection{Activities}\label{activities-3}

\textbf{Activity 2.1.1: Planning a Tip Calculator} Write pseudocode
before coding:

\begin{verbatim}
1. GET bill amount from user
2. CALCULATE tip as 15% of bill amount
3. CALCULATE total amount (bill + tip)
4. PRINT tip amount
5. PRINT total amount
\end{verbatim}

\textbf{Activity 2.1.2: Tracing a Simple Program} Without running the
code, trace the execution and predict output:

\begin{Shaded}
\begin{Highlighting}[]
\NormalTok{a }\OperatorTok{=} \DecValTok{10}
\NormalTok{b }\OperatorTok{=} \DecValTok{5}
\NormalTok{c }\OperatorTok{=}\NormalTok{ a }\OperatorTok{*}\NormalTok{ b}
\NormalTok{a }\OperatorTok{=}\NormalTok{ a }\OperatorTok{+} \DecValTok{2}
\BuiltInTok{print}\NormalTok{(c)}
\BuiltInTok{print}\NormalTok{(a)}
\end{Highlighting}
\end{Shaded}

\subsubsection{Skills Checklist}\label{skills-checklist-7}

\begin{itemize}
\tightlist
\item[$\square$]
  I can describe the purpose of the software development process.
\item[$\square$]
  I can list and briefly explain the main phases of software
  development.
\item[$\square$]
  I understand that problem-solving and design happen before writing
  code.
\end{itemize}

\begin{center}\rule{0.5\linewidth}{0.5pt}\end{center}

\subsection{Lesson 2.2: Elements of Programs --- Expressions and
Assignment}\label{lesson-2.2-elements-of-programs-expressions-and-assignment}

\subsubsection{Reading References}\label{reading-references-8}

\begin{itemize}
\tightlist
\item
  Zelle: Chapter 2, pages 24--26
\end{itemize}

\subsubsection{Expressions and
Assignment}\label{expressions-and-assignment}

An \textbf{expression} is any combination of values, variables, and
operators that evaluates to a single result. Python follows the standard
mathematical order of operations (PEMDAS).

\begin{Shaded}
\begin{Highlighting}[]
\NormalTok{result1 }\OperatorTok{=} \DecValTok{10} \OperatorTok{/} \DecValTok{2} \OperatorTok{+} \DecValTok{3}  \CommentTok{\# 10/2 is 5, then 5+3 = 8.0}
\NormalTok{message }\OperatorTok{=} \StringTok{"Python"} \OperatorTok{+} \StringTok{" "} \OperatorTok{+} \StringTok{"is fun!"}
\end{Highlighting}
\end{Shaded}

\textbf{Assignment:}

\begin{Shaded}
\begin{Highlighting}[]
\NormalTok{score }\OperatorTok{=} \DecValTok{50}
\NormalTok{score }\OperatorTok{=}\NormalTok{ score }\OperatorTok{+} \DecValTok{20}  \CommentTok{\# score is now 70}
\end{Highlighting}
\end{Shaded}

\subsubsection{Examples}\label{examples-2}

\textbf{Example 2.2.1: Expression Evaluation}

\begin{Shaded}
\begin{Highlighting}[]
\NormalTok{result1 }\OperatorTok{=} \DecValTok{10} \OperatorTok{/} \DecValTok{2} \OperatorTok{+} \DecValTok{3}
\BuiltInTok{print}\NormalTok{(}\StringTok{"Result 1:"}\NormalTok{, result1)   }\CommentTok{\# Output: Result 1: 8.0}

\NormalTok{message }\OperatorTok{=} \StringTok{"Python"} \OperatorTok{+} \StringTok{" "} \OperatorTok{+} \StringTok{"is fun!"}
\BuiltInTok{print}\NormalTok{(}\StringTok{"Message:"}\NormalTok{, message)    }\CommentTok{\# Output: Message: Python is fun!}
\end{Highlighting}
\end{Shaded}

\textbf{Example 2.2.2: Reassigning Variables}

\begin{Shaded}
\begin{Highlighting}[]
\NormalTok{score }\OperatorTok{=} \DecValTok{50}
\BuiltInTok{print}\NormalTok{(}\StringTok{"Initial score:"}\NormalTok{, score)}
\NormalTok{score }\OperatorTok{=}\NormalTok{ score }\OperatorTok{+} \DecValTok{20}
\BuiltInTok{print}\NormalTok{(}\StringTok{"New score:"}\NormalTok{, score)     }\CommentTok{\# Output: New score: 70}
\end{Highlighting}
\end{Shaded}

\subsubsection{Skills Checklist}\label{skills-checklist-8}

\begin{itemize}
\tightlist
\item[$\square$]
  I can identify an expression in Python code.
\item[$\square$]
  I can predict the result of simple Python expressions.
\item[$\square$]
  I can explain how the assignment operator \texttt{=} works.
\item[$\square$]
  I can demonstrate how to update a variable's value.
\end{itemize}

\begin{center}\rule{0.5\linewidth}{0.5pt}\end{center}

\subsection{Lesson 2.3: Input and Output
Revisited}\label{lesson-2.3-input-and-output-revisited}

\subsubsection{Reading References}\label{reading-references-9}

\begin{itemize}
\tightlist
\item
  Zelle: Chapter 2, pages 27--33
\item
  Matthes: Chapter 7, pages 114--116
\end{itemize}

\subsubsection{\texorpdfstring{The \texttt{input()}
Function}{The input() Function}}\label{the-input-function}

The \texttt{input()} function pauses program execution and waits for the
user to type something and press Enter. \textbf{It always returns a
string.}

\begin{Shaded}
\begin{Highlighting}[]
\NormalTok{user\_name }\OperatorTok{=} \BuiltInTok{input}\NormalTok{(}\StringTok{"What is your name? "}\NormalTok{)}
\BuiltInTok{print}\NormalTok{(}\SpecialStringTok{f"Hello, }\SpecialCharTok{\{}\NormalTok{user\_name}\SpecialCharTok{\}}\SpecialStringTok{! Nice to meet you."}\NormalTok{)}
\end{Highlighting}
\end{Shaded}

\subsubsection{F-Strings (Formatted String
Literals)}\label{f-strings-formatted-string-literals}

F-strings provide a clean way to embed variable values directly into a
string:

\begin{Shaded}
\begin{Highlighting}[]
\NormalTok{name }\OperatorTok{=} \StringTok{"Alice"}
\NormalTok{age }\OperatorTok{=} \DecValTok{30}
\BuiltInTok{print}\NormalTok{(}\SpecialStringTok{f"Hello, }\SpecialCharTok{\{}\NormalTok{name}\SpecialCharTok{\}}\SpecialStringTok{! You are }\SpecialCharTok{\{}\NormalTok{age}\SpecialCharTok{\}}\SpecialStringTok{ years old."}\NormalTok{)}
\end{Highlighting}
\end{Shaded}

\subsubsection{Converting Input for
Calculations}\label{converting-input-for-calculations}

Since \texttt{input()} always returns a string, you must convert it for
math:

\begin{Shaded}
\begin{Highlighting}[]
\NormalTok{age\_str }\OperatorTok{=} \BuiltInTok{input}\NormalTok{(}\StringTok{"How old are you? "}\NormalTok{)}
\NormalTok{age\_int }\OperatorTok{=} \BuiltInTok{int}\NormalTok{(age\_str)}
\BuiltInTok{print}\NormalTok{(}\SpecialStringTok{f"Next year you will be }\SpecialCharTok{\{}\NormalTok{age\_int }\OperatorTok{+} \DecValTok{1}\SpecialCharTok{\}}\SpecialStringTok{ years old."}\NormalTok{)}
\end{Highlighting}
\end{Shaded}

\subsubsection{Examples}\label{examples-3}

\textbf{Example 2.3.1: Basic Input}

\begin{Shaded}
\begin{Highlighting}[]
\NormalTok{user\_name }\OperatorTok{=} \BuiltInTok{input}\NormalTok{(}\StringTok{"What is your name? "}\NormalTok{)}
\BuiltInTok{print}\NormalTok{(}\SpecialStringTok{f"Hello, }\SpecialCharTok{\{}\NormalTok{user\_name}\SpecialCharTok{\}}\SpecialStringTok{! Nice to meet you."}\NormalTok{)}
\end{Highlighting}
\end{Shaded}

\textbf{Example 2.3.2: Numeric Input and Calculation}

\begin{Shaded}
\begin{Highlighting}[]
\NormalTok{num1\_str }\OperatorTok{=} \BuiltInTok{input}\NormalTok{(}\StringTok{"Enter the first number: "}\NormalTok{)}
\NormalTok{num1 }\OperatorTok{=} \BuiltInTok{int}\NormalTok{(num1\_str)}

\NormalTok{num2\_str }\OperatorTok{=} \BuiltInTok{input}\NormalTok{(}\StringTok{"Enter the second number: "}\NormalTok{)}
\NormalTok{num2 }\OperatorTok{=} \BuiltInTok{int}\NormalTok{(num2\_str)}

\NormalTok{sum\_result }\OperatorTok{=}\NormalTok{ num1 }\OperatorTok{+}\NormalTok{ num2}
\BuiltInTok{print}\NormalTok{(}\SpecialStringTok{f"The sum of }\SpecialCharTok{\{}\NormalTok{num1}\SpecialCharTok{\}}\SpecialStringTok{ and }\SpecialCharTok{\{}\NormalTok{num2}\SpecialCharTok{\}}\SpecialStringTok{ is }\SpecialCharTok{\{}\NormalTok{sum\_result}\SpecialCharTok{\}}\SpecialStringTok{."}\NormalTok{)}
\end{Highlighting}
\end{Shaded}

\subsubsection{Activities}\label{activities-4}

\textbf{Activity 2.3.1: Temperature Converter}

\begin{Shaded}
\begin{Highlighting}[]
\NormalTok{fahrenheit\_str }\OperatorTok{=} \BuiltInTok{input}\NormalTok{(}\StringTok{"Enter temperature in Fahrenheit: "}\NormalTok{)}
\NormalTok{fahrenheit }\OperatorTok{=} \BuiltInTok{float}\NormalTok{(fahrenheit\_str)}
\NormalTok{celsius }\OperatorTok{=}\NormalTok{ (fahrenheit }\OperatorTok{{-}} \DecValTok{32}\NormalTok{) }\OperatorTok{*} \DecValTok{5} \OperatorTok{/} \DecValTok{9}
\BuiltInTok{print}\NormalTok{(}\SpecialStringTok{f"}\SpecialCharTok{\{}\NormalTok{fahrenheit}\SpecialCharTok{\}}\SpecialStringTok{°F is }\SpecialCharTok{\{}\NormalTok{celsius}\SpecialCharTok{:.2f\}}\SpecialStringTok{°C"}\NormalTok{)}
\end{Highlighting}
\end{Shaded}

\subsubsection{Skills Checklist}\label{skills-checklist-9}

\begin{itemize}
\tightlist
\item[$\square$]
  I can use the \texttt{input()} function to ask the user for
  information.
\item[$\square$]
  I can provide a clear prompt message when using \texttt{input()}.
\item[$\square$]
  I remember that \texttt{input()} returns a string, even for numbers.
\item[$\square$]
  I can combine text and variable values in a \texttt{print()} statement
  using f-strings.
\end{itemize}

\begin{center}\rule{0.5\linewidth}{0.5pt}\end{center}

\subsection{Lesson 2.4: Computing with Numbers and Type
Conversions}\label{lesson-2.4-computing-with-numbers-and-type-conversions}

\subsubsection{Reading References}\label{reading-references-10}

\begin{itemize}
\tightlist
\item
  Zelle: Chapter 3, pages 41--48
\item
  Matthes: Chapter 2, pages 26--28
\end{itemize}

\subsubsection{Integers and Floats}\label{integers-and-floats}

{\def\LTcaptype{none} % do not increment counter
\begin{longtable}[]{@{}lll@{}}
\toprule\noalign{}
Type & Description & Example \\
\midrule\noalign{}
\endhead
\bottomrule\noalign{}
\endlastfoot
\texttt{int} & Whole numbers (exact) & \texttt{5}, \texttt{-20},
\texttt{1000000} \\
\texttt{float} & Numbers with decimals & \texttt{3.14}, \texttt{-2.5},
\texttt{6.0} \\
\end{longtable}
}

\textbf{Key behaviors:}

\begin{itemize}
\tightlist
\item
  Division \texttt{/} always returns a float: \texttt{10\ /\ 2}
  -\textgreater{} \texttt{5.0}
\item
  Floor division \texttt{//} returns the integer part:
  \texttt{10\ //\ 3} -\textgreater{} \texttt{3}
\item
  Mixed operations promote to float: \texttt{10\ +\ 3.5} -\textgreater{}
  \texttt{13.5}
\end{itemize}

\subsubsection{Type Conversion
Functions}\label{type-conversion-functions}

\begin{Shaded}
\begin{Highlighting}[]
\BuiltInTok{int}\NormalTok{(}\StringTok{"123"}\NormalTok{)      }\CommentTok{\# {-}\textgreater{} 123  (truncates: int(3.9) {-}\textgreater{} 3)}
\BuiltInTok{float}\NormalTok{(}\StringTok{"3.14"}\NormalTok{)   }\CommentTok{\# {-}\textgreater{} 3.14}
\BuiltInTok{str}\NormalTok{(}\DecValTok{42}\NormalTok{)         }\CommentTok{\# {-}\textgreater{} "42"}
\end{Highlighting}
\end{Shaded}

\subsubsection{Examples}\label{examples-4}

\textbf{Example 2.4.1: Type Conversion}

\begin{Shaded}
\begin{Highlighting}[]
\NormalTok{num\_str }\OperatorTok{=} \StringTok{"123"}
\NormalTok{num\_int }\OperatorTok{=} \BuiltInTok{int}\NormalTok{(num\_str)}
\BuiltInTok{print}\NormalTok{(}\SpecialStringTok{f"String \textquotesingle{}}\SpecialCharTok{\{}\NormalTok{num\_str}\SpecialCharTok{\}}\SpecialStringTok{\textquotesingle{} as integer: }\SpecialCharTok{\{}\NormalTok{num\_int}\SpecialCharTok{\}}\SpecialStringTok{, type: }\SpecialCharTok{\{}\BuiltInTok{type}\NormalTok{(num\_int)}\SpecialCharTok{\}}\SpecialStringTok{"}\NormalTok{)}

\NormalTok{num\_float }\OperatorTok{=} \BuiltInTok{float}\NormalTok{(num\_str)}
\BuiltInTok{print}\NormalTok{(}\SpecialStringTok{f"String \textquotesingle{}}\SpecialCharTok{\{}\NormalTok{num\_str}\SpecialCharTok{\}}\SpecialStringTok{\textquotesingle{} as float: }\SpecialCharTok{\{}\NormalTok{num\_float}\SpecialCharTok{\}}\SpecialStringTok{, type: }\SpecialCharTok{\{}\BuiltInTok{type}\NormalTok{(num\_float)}\SpecialCharTok{\}}\SpecialStringTok{"}\NormalTok{)}

\NormalTok{my\_number }\OperatorTok{=} \DecValTok{42}
\NormalTok{my\_text }\OperatorTok{=} \BuiltInTok{str}\NormalTok{(my\_number)}
\BuiltInTok{print}\NormalTok{(}\SpecialStringTok{f"Integer }\SpecialCharTok{\{}\NormalTok{my\_number}\SpecialCharTok{\}}\SpecialStringTok{ as string: \textquotesingle{}}\SpecialCharTok{\{}\NormalTok{my\_text}\SpecialCharTok{\}}\SpecialStringTok{\textquotesingle{}, type: }\SpecialCharTok{\{}\BuiltInTok{type}\NormalTok{(my\_text)}\SpecialCharTok{\}}\SpecialStringTok{"}\NormalTok{)}
\end{Highlighting}
\end{Shaded}

\textbf{Common Type Error:}

\begin{Shaded}
\begin{Highlighting}[]
\CommentTok{\# This will cause a TypeError:}
\CommentTok{\# print("My age is " + 14)}
\CommentTok{\# Fix: print("My age is " + str(14))}
\CommentTok{\# Or: print(f"My age is \{14\}")}
\end{Highlighting}
\end{Shaded}

\subsubsection{Activities}\label{activities-5}

\textbf{Activity 2.4.1: Mixed Bag Calculator}

\begin{Shaded}
\begin{Highlighting}[]
\NormalTok{age\_str }\OperatorTok{=} \BuiltInTok{input}\NormalTok{(}\StringTok{"Enter your age: "}\NormalTok{)}
\NormalTok{age }\OperatorTok{=} \BuiltInTok{int}\NormalTok{(age\_str)}

\NormalTok{weight\_str }\OperatorTok{=} \BuiltInTok{input}\NormalTok{(}\StringTok{"Enter your weight in kg: "}\NormalTok{)}
\NormalTok{weight }\OperatorTok{=} \BuiltInTok{float}\NormalTok{(weight\_str)}

\NormalTok{phrase }\OperatorTok{=} \BuiltInTok{input}\NormalTok{(}\StringTok{"Enter a short phrase: "}\NormalTok{)}

\BuiltInTok{print}\NormalTok{(}\SpecialStringTok{f"Age doubled: }\SpecialCharTok{\{}\NormalTok{age }\OperatorTok{*} \DecValTok{2}\SpecialCharTok{\}}\SpecialStringTok{"}\NormalTok{)}
\BuiltInTok{print}\NormalTok{(}\SpecialStringTok{f"Weight + age: }\SpecialCharTok{\{}\NormalTok{weight }\OperatorTok{+}\NormalTok{ age}\SpecialCharTok{\}}\SpecialStringTok{"}\NormalTok{)}
\BuiltInTok{print}\NormalTok{(}\SpecialStringTok{f"Phrase length: }\SpecialCharTok{\{}\BuiltInTok{len}\NormalTok{(phrase)}\SpecialCharTok{\}}\SpecialStringTok{ characters"}\NormalTok{)}
\end{Highlighting}
\end{Shaded}

\subsubsection{Skills Checklist}\label{skills-checklist-10}

\begin{itemize}
\tightlist
\item[$\square$]
  I can identify when to use an integer versus a float.
\item[$\square$]
  I can convert a string to an integer using \texttt{int()}.
\item[$\square$]
  I can convert a string to a float using \texttt{float()}.
\item[$\square$]
  I can convert a number to a string using \texttt{str()}.
\item[$\square$]
  I can apply type conversions to perform calculations on user input.
\item[$\square$]
  I understand and can apply the order of operations.
\end{itemize}

\begin{center}\rule{0.5\linewidth}{0.5pt}\end{center}

\subsection{Lesson 2.5: Using the Math
Library}\label{lesson-2.5-using-the-math-library}

\subsubsection{Reading References}\label{reading-references-11}

\begin{itemize}
\tightlist
\item
  Zelle: Chapter 3, pages 49--51
\end{itemize}

\subsubsection{Importing the Math
Module}\label{importing-the-math-module}

\begin{Shaded}
\begin{Highlighting}[]
\ImportTok{import}\NormalTok{ math}
\end{Highlighting}
\end{Shaded}

\subsubsection{Key Math Functions and
Constants}\label{key-math-functions-and-constants}

{\def\LTcaptype{none} % do not increment counter
\begin{longtable}[]{@{}
  >{\raggedright\arraybackslash}p{(\linewidth - 4\tabcolsep) * \real{0.2533}}
  >{\raggedright\arraybackslash}p{(\linewidth - 4\tabcolsep) * \real{0.3467}}
  >{\raggedright\arraybackslash}p{(\linewidth - 4\tabcolsep) * \real{0.4000}}@{}}
\toprule\noalign{}
\begin{minipage}[b]{\linewidth}\raggedright
Function/Constant
\end{minipage} & \begin{minipage}[b]{\linewidth}\raggedright
Description
\end{minipage} & \begin{minipage}[b]{\linewidth}\raggedright
Example
\end{minipage} \\
\midrule\noalign{}
\endhead
\bottomrule\noalign{}
\endlastfoot
\texttt{math.pi} & π ≈ 3.14159 & \texttt{math.pi} \\
\texttt{math.e} & e ≈ 2.71828 & \texttt{math.e} \\
\texttt{math.sqrt(x)} & Square root of x & \texttt{math.sqrt(25)}
-\textgreater{} \texttt{5.0} \\
\texttt{math.pow(x,\ y)} & x raised to power y &
\texttt{math.pow(2,\ 4)} -\textgreater{} \texttt{16.0} \\
\texttt{math.ceil(x)} & Round up & \texttt{math.ceil(4.2)}
-\textgreater{} \texttt{5} \\
\texttt{math.floor(x)} & Round down & \texttt{math.floor(4.8)}
-\textgreater{} \texttt{4} \\
\texttt{math.log(x,\ base)} & Logarithm & \texttt{math.log(100,\ 10)}
-\textgreater{} \texttt{2.0} \\
\texttt{math.sin(x)} & Sine (in radians) & \texttt{math.sin(math.pi/2)}
-\textgreater{} \texttt{1.0} \\
\texttt{math.radians(deg)} & Convert degrees to radians &
\texttt{math.radians(90)} \\
\end{longtable}
}

\subsubsection{Examples}\label{examples-5}

\textbf{Example 2.5.1: Circle Calculations}

\begin{Shaded}
\begin{Highlighting}[]
\ImportTok{import}\NormalTok{ math}

\NormalTok{radius }\OperatorTok{=} \DecValTok{5}
\NormalTok{area }\OperatorTok{=}\NormalTok{ math.pi }\OperatorTok{*}\NormalTok{ radius }\OperatorTok{**} \DecValTok{2}
\NormalTok{circumference }\OperatorTok{=} \DecValTok{2} \OperatorTok{*}\NormalTok{ math.pi }\OperatorTok{*}\NormalTok{ radius}

\BuiltInTok{print}\NormalTok{(}\SpecialStringTok{f"Area of circle with radius }\SpecialCharTok{\{}\NormalTok{radius}\SpecialCharTok{\}}\SpecialStringTok{: }\SpecialCharTok{\{}\NormalTok{area}\SpecialCharTok{:.2f\}}\SpecialStringTok{"}\NormalTok{)}
\BuiltInTok{print}\NormalTok{(}\SpecialStringTok{f"Circumference: }\SpecialCharTok{\{}\NormalTok{circumference}\SpecialCharTok{:.2f\}}\SpecialStringTok{"}\NormalTok{)}

\NormalTok{number }\OperatorTok{=} \DecValTok{81}
\NormalTok{square\_root }\OperatorTok{=}\NormalTok{ math.sqrt(number)}
\BuiltInTok{print}\NormalTok{(}\SpecialStringTok{f"Square root of }\SpecialCharTok{\{}\NormalTok{number}\SpecialCharTok{\}}\SpecialStringTok{: }\SpecialCharTok{\{}\NormalTok{square\_root}\SpecialCharTok{\}}\SpecialStringTok{"}\NormalTok{)}
\end{Highlighting}
\end{Shaded}

\subsubsection{Activities}\label{activities-6}

\textbf{Activity 2.5.1: Circle Calculator}

\begin{Shaded}
\begin{Highlighting}[]
\ImportTok{import}\NormalTok{ math}

\NormalTok{radius\_str }\OperatorTok{=} \BuiltInTok{input}\NormalTok{(}\StringTok{"Enter the radius: "}\NormalTok{)}
\NormalTok{radius }\OperatorTok{=} \BuiltInTok{float}\NormalTok{(radius\_str)}

\NormalTok{area }\OperatorTok{=}\NormalTok{ math.pi }\OperatorTok{*}\NormalTok{ radius }\OperatorTok{**} \DecValTok{2}
\NormalTok{circumference }\OperatorTok{=} \DecValTok{2} \OperatorTok{*}\NormalTok{ math.pi }\OperatorTok{*}\NormalTok{ radius}

\BuiltInTok{print}\NormalTok{(}\SpecialStringTok{f"Area: }\SpecialCharTok{\{}\NormalTok{area}\SpecialCharTok{:.2f\}}\SpecialStringTok{"}\NormalTok{)}
\BuiltInTok{print}\NormalTok{(}\SpecialStringTok{f"Circumference: }\SpecialCharTok{\{}\NormalTok{circumference}\SpecialCharTok{:.2f\}}\SpecialStringTok{"}\NormalTok{)}
\end{Highlighting}
\end{Shaded}

\textbf{Activity 2.5.2: Distance Finder}

\begin{Shaded}
\begin{Highlighting}[]
\ImportTok{import}\NormalTok{ math}

\NormalTok{x }\OperatorTok{=} \BuiltInTok{float}\NormalTok{(}\BuiltInTok{input}\NormalTok{(}\StringTok{"Enter x coordinate: "}\NormalTok{))}
\NormalTok{y }\OperatorTok{=} \BuiltInTok{float}\NormalTok{(}\BuiltInTok{input}\NormalTok{(}\StringTok{"Enter y coordinate: "}\NormalTok{))}
\NormalTok{distance }\OperatorTok{=}\NormalTok{ math.sqrt(x}\OperatorTok{**}\DecValTok{2} \OperatorTok{+}\NormalTok{ y}\OperatorTok{**}\DecValTok{2}\NormalTok{)}
\BuiltInTok{print}\NormalTok{(}\SpecialStringTok{f"Distance from origin: }\SpecialCharTok{\{}\NormalTok{distance}\SpecialCharTok{:.2f\}}\SpecialStringTok{"}\NormalTok{)}
\end{Highlighting}
\end{Shaded}

\subsubsection{Skills Checklist}\label{skills-checklist-11}

\begin{itemize}
\tightlist
\item[$\square$]
  I can explain what a Python module is.
\item[$\square$]
  I can import the \texttt{math} module into my program.
\item[$\square$]
  I can use \texttt{math.pi} and \texttt{math.sqrt()}.
\item[$\square$]
  I can use \texttt{math.ceil()} to round up and \texttt{math.floor()}
  to round down.
\end{itemize}

\newpage

\section{Unit 3: Control Flow --- Decisions and
Loops}\label{unit-3-control-flow-decisions-and-loops-1}

\subsection{Overview}\label{overview-3}

This unit introduces how programs make decisions using \texttt{if}
statements and perform repetitive tasks with \texttt{for} loops, along
with fundamental string manipulation techniques.

\textbf{Objectives:}

\begin{itemize}
\tightlist
\item
  Control program flow with \texttt{if}, \texttt{elif}, and
  \texttt{else}.
\item
  Iterate over sequences using \texttt{for} loops.
\item
  Understand and manipulate strings.
\end{itemize}

\begin{center}\rule{0.5\linewidth}{0.5pt}\end{center}

\subsection{\texorpdfstring{Lesson 3.1: Boolean Expressions and
\texttt{if}
Statements}{Lesson 3.1: Boolean Expressions and if Statements}}\label{lesson-3.1-boolean-expressions-and-if-statements}

\subsubsection{Reading References}\label{reading-references-12}

\begin{itemize}
\tightlist
\item
  Zelle: Chapter 7, pages 165--173
\item
  Matthes: Chapter 5, pages 71--78
\end{itemize}

\subsubsection{Comparison and Boolean
Operators}\label{comparison-and-boolean-operators}

{\def\LTcaptype{none} % do not increment counter
\begin{longtable}[]{@{}llll@{}}
\toprule\noalign{}
Operator & Meaning & Example & Result \\
\midrule\noalign{}
\endhead
\bottomrule\noalign{}
\endlastfoot
\texttt{==} & Equal to & \texttt{5\ ==\ 5} & \texttt{True} \\
\texttt{!=} & Not equal to & \texttt{5\ !=\ 10} & \texttt{True} \\
\texttt{\textless{}} & Less than & \texttt{5\ \textless{}\ 10} &
\texttt{True} \\
\texttt{\textgreater{}} & Greater than & \texttt{10\ \textgreater{}\ 5}
& \texttt{True} \\
\texttt{\textless{}=} & Less than or equal to &
\texttt{5\ \textless{}=\ 5} & \texttt{True} \\
\texttt{\textgreater{}=} & Greater than or equal &
\texttt{10\ \textgreater{}=\ 5} & \texttt{True} \\
\end{longtable}
}

{\def\LTcaptype{none} % do not increment counter
\begin{longtable}[]{@{}
  >{\raggedright\arraybackslash}p{(\linewidth - 6\tabcolsep) * \real{0.2388}}
  >{\raggedright\arraybackslash}p{(\linewidth - 6\tabcolsep) * \real{0.4179}}
  >{\raggedright\arraybackslash}p{(\linewidth - 6\tabcolsep) * \real{0.2388}}
  >{\raggedright\arraybackslash}p{(\linewidth - 6\tabcolsep) * \real{0.1045}}@{}}
\toprule\noalign{}
\begin{minipage}[b]{\linewidth}\raggedright
Boolean Operator
\end{minipage} & \begin{minipage}[b]{\linewidth}\raggedright
Description
\end{minipage} & \begin{minipage}[b]{\linewidth}\raggedright
Example
\end{minipage} & \begin{minipage}[b]{\linewidth}\raggedright
Result
\end{minipage} \\
\midrule\noalign{}
\endhead
\bottomrule\noalign{}
\endlastfoot
\texttt{and} & True if both are true & \texttt{True\ and\ False} &
\texttt{False} \\
\texttt{or} & True if at least one is true & \texttt{True\ or\ False} &
\texttt{True} \\
\texttt{not} & Inverts the truth value & \texttt{not\ True} &
\texttt{False} \\
\end{longtable}
}

\subsubsection{\texorpdfstring{The \texttt{if}
Statement}{The if Statement}}\label{the-if-statement}

\begin{Shaded}
\begin{Highlighting}[]
\NormalTok{age }\OperatorTok{=} \DecValTok{15}
\ControlFlowTok{if}\NormalTok{ age }\OperatorTok{\textgreater{}=} \DecValTok{18}\NormalTok{:}
    \BuiltInTok{print}\NormalTok{(}\StringTok{"You are old enough to vote."}\NormalTok{)}
\BuiltInTok{print}\NormalTok{(}\StringTok{"Program finished."}\NormalTok{)  }\CommentTok{\# This runs regardless}
\end{Highlighting}
\end{Shaded}

\textbf{Indentation defines which lines belong to the \texttt{if}
block.}

\subsubsection{Examples}\label{examples-6}

\textbf{Example 3.1.1: Simple Age Check}

\begin{Shaded}
\begin{Highlighting}[]
\NormalTok{age }\OperatorTok{=} \DecValTok{17}
\ControlFlowTok{if}\NormalTok{ age }\OperatorTok{\textgreater{}=} \DecValTok{18}\NormalTok{:}
    \BuiltInTok{print}\NormalTok{(}\StringTok{"You are an adult."}\NormalTok{)}
\ControlFlowTok{else}\NormalTok{:}
    \BuiltInTok{print}\NormalTok{(}\StringTok{"You are a minor."}\NormalTok{)}
\end{Highlighting}
\end{Shaded}

\textbf{Example 3.1.2: Password Length Check}

\begin{Shaded}
\begin{Highlighting}[]
\NormalTok{password }\OperatorTok{=} \StringTok{"mysecret"}
\ControlFlowTok{if} \BuiltInTok{len}\NormalTok{(password) }\OperatorTok{\textless{}} \DecValTok{8}\NormalTok{:}
    \BuiltInTok{print}\NormalTok{(}\StringTok{"Weak password: Must be at least 8 characters."}\NormalTok{)}
\end{Highlighting}
\end{Shaded}

\subsubsection{Activities}\label{activities-7}

\textbf{Activity 3.1.1: Number Comparison Game}

\begin{Shaded}
\begin{Highlighting}[]
\NormalTok{num1 }\OperatorTok{=} \BuiltInTok{int}\NormalTok{(}\BuiltInTok{input}\NormalTok{(}\StringTok{"Enter the first number: "}\NormalTok{))}
\NormalTok{num2 }\OperatorTok{=} \BuiltInTok{int}\NormalTok{(}\BuiltInTok{input}\NormalTok{(}\StringTok{"Enter the second number: "}\NormalTok{))}

\ControlFlowTok{if}\NormalTok{ num1 }\OperatorTok{\textgreater{}}\NormalTok{ num2:}
    \BuiltInTok{print}\NormalTok{(}\SpecialStringTok{f"}\SpecialCharTok{\{}\NormalTok{num1}\SpecialCharTok{\}}\SpecialStringTok{ is greater than }\SpecialCharTok{\{}\NormalTok{num2}\SpecialCharTok{\}}\SpecialStringTok{."}\NormalTok{)}
\ControlFlowTok{if}\NormalTok{ num1 }\OperatorTok{==}\NormalTok{ num2:}
    \BuiltInTok{print}\NormalTok{(}\StringTok{"The numbers are equal."}\NormalTok{)}
\ControlFlowTok{if}\NormalTok{ num1 }\OperatorTok{\textless{}}\NormalTok{ num2:}
    \BuiltInTok{print}\NormalTok{(}\SpecialStringTok{f"}\SpecialCharTok{\{}\NormalTok{num1}\SpecialCharTok{\}}\SpecialStringTok{ is less than }\SpecialCharTok{\{}\NormalTok{num2}\SpecialCharTok{\}}\SpecialStringTok{."}\NormalTok{)}
\end{Highlighting}
\end{Shaded}

\subsubsection{Skills Checklist}\label{skills-checklist-12}

\begin{itemize}
\tightlist
\item[$\square$]
  I can identify \texttt{True} and \texttt{False} as Boolean values.
\item[$\square$]
  I can use comparison operators to create Boolean expressions.
\item[$\square$]
  I can write a basic \texttt{if} statement.
\item[$\square$]
  I can correctly indent the code block belonging to an \texttt{if}
  statement.
\end{itemize}

\begin{center}\rule{0.5\linewidth}{0.5pt}\end{center}

\subsection{\texorpdfstring{Lesson 3.2: Chained Conditionals
(\texttt{elif},
\texttt{else})}{Lesson 3.2: Chained Conditionals (elif, else)}}\label{lesson-3.2-chained-conditionals-elif-else}

\subsubsection{Reading References}\label{reading-references-13}

\begin{itemize}
\tightlist
\item
  Zelle: Chapter 7, pages 173--176
\item
  Matthes: Chapter 5, pages 79--84
\end{itemize}

\subsubsection{\texorpdfstring{\texttt{if-elif-else}
Structure}{if-elif-else Structure}}\label{if-elif-else-structure}

Python checks each condition in \textbf{sequential order}. As soon as a
condition is \texttt{True}, it runs that block and skips the rest.

\begin{Shaded}
\begin{Highlighting}[]
\NormalTok{score }\OperatorTok{=} \DecValTok{85}

\ControlFlowTok{if}\NormalTok{ score }\OperatorTok{\textgreater{}=} \DecValTok{90}\NormalTok{:}
    \BuiltInTok{print}\NormalTok{(}\StringTok{"Grade: A"}\NormalTok{)}
\ControlFlowTok{elif}\NormalTok{ score }\OperatorTok{\textgreater{}=} \DecValTok{80}\NormalTok{:}
    \BuiltInTok{print}\NormalTok{(}\StringTok{"Grade: B"}\NormalTok{)}
\ControlFlowTok{elif}\NormalTok{ score }\OperatorTok{\textgreater{}=} \DecValTok{70}\NormalTok{:}
    \BuiltInTok{print}\NormalTok{(}\StringTok{"Grade: C"}\NormalTok{)}
\ControlFlowTok{else}\NormalTok{:}
    \BuiltInTok{print}\NormalTok{(}\StringTok{"Grade: F"}\NormalTok{)}
\end{Highlighting}
\end{Shaded}

Output: \texttt{Grade:\ B}

\subsubsection{Examples}\label{examples-7}

\textbf{Example 3.2.1: Stages of Life}

\begin{Shaded}
\begin{Highlighting}[]
\NormalTok{age }\OperatorTok{=} \DecValTok{25}
\ControlFlowTok{if}\NormalTok{ age }\OperatorTok{\textless{}} \DecValTok{2}\NormalTok{:}
    \BuiltInTok{print}\NormalTok{(}\StringTok{"You are a baby."}\NormalTok{)}
\ControlFlowTok{elif}\NormalTok{ age }\OperatorTok{\textless{}} \DecValTok{4}\NormalTok{:}
    \BuiltInTok{print}\NormalTok{(}\StringTok{"You are a toddler."}\NormalTok{)}
\ControlFlowTok{elif}\NormalTok{ age }\OperatorTok{\textless{}} \DecValTok{13}\NormalTok{:}
    \BuiltInTok{print}\NormalTok{(}\StringTok{"You are a kid."}\NormalTok{)}
\ControlFlowTok{elif}\NormalTok{ age }\OperatorTok{\textless{}} \DecValTok{20}\NormalTok{:}
    \BuiltInTok{print}\NormalTok{(}\StringTok{"You are a teenager."}\NormalTok{)}
\ControlFlowTok{else}\NormalTok{:}
    \BuiltInTok{print}\NormalTok{(}\StringTok{"You are an adult."}\NormalTok{)}
\end{Highlighting}
\end{Shaded}

\subsubsection{Activities}\label{activities-8}

\textbf{Activity 3.2.1: Movie Ticket Pricing}

\begin{Shaded}
\begin{Highlighting}[]
\NormalTok{age }\OperatorTok{=} \BuiltInTok{int}\NormalTok{(}\BuiltInTok{input}\NormalTok{(}\StringTok{"Enter your age: "}\NormalTok{))}

\ControlFlowTok{if}\NormalTok{ age }\OperatorTok{\textless{}=} \DecValTok{5}\NormalTok{:}
    \BuiltInTok{print}\NormalTok{(}\StringTok{"Ticket price: Free"}\NormalTok{)}
\ControlFlowTok{elif}\NormalTok{ age }\OperatorTok{\textless{}=} \DecValTok{12}\NormalTok{:}
    \BuiltInTok{print}\NormalTok{(}\StringTok{"Ticket price: $7"}\NormalTok{)}
\ControlFlowTok{elif}\NormalTok{ age }\OperatorTok{\textless{}=} \DecValTok{64}\NormalTok{:}
    \BuiltInTok{print}\NormalTok{(}\StringTok{"Ticket price: $12"}\NormalTok{)}
\ControlFlowTok{else}\NormalTok{:}
    \BuiltInTok{print}\NormalTok{(}\StringTok{"Ticket price: $5"}\NormalTok{)}
\end{Highlighting}
\end{Shaded}

\subsubsection{Skills Checklist}\label{skills-checklist-13}

\begin{itemize}
\tightlist
\item[$\square$]
  I can use an \texttt{else} statement to provide an alternative path.
\item[$\square$]
  I can use \texttt{elif} to check multiple conditions in order.
\item[$\square$]
  I can use \texttt{and}, \texttt{or}, and \texttt{not} operators to
  combine conditions.
\end{itemize}

\begin{center}\rule{0.5\linewidth}{0.5pt}\end{center}

\subsection{Lesson 3.3: Introduction to
Lists}\label{lesson-3.3-introduction-to-lists}

\subsubsection{Reading References}\label{reading-references-14}

\begin{itemize}
\tightlist
\item
  Zelle: Chapter 5, pages 105--106
\item
  Matthes: Chapter 3, pages 33--47
\end{itemize}

\subsubsection{Creating and Accessing
Lists}\label{creating-and-accessing-lists}

\begin{Shaded}
\begin{Highlighting}[]
\NormalTok{colors }\OperatorTok{=}\NormalTok{ [}\StringTok{"red"}\NormalTok{, }\StringTok{"blue"}\NormalTok{, }\StringTok{"green"}\NormalTok{, }\StringTok{"yellow"}\NormalTok{]}

\BuiltInTok{print}\NormalTok{(colors[}\DecValTok{0}\NormalTok{])   }\CommentTok{\# "red"   (first item)}
\BuiltInTok{print}\NormalTok{(colors[}\OperatorTok{{-}}\DecValTok{1}\NormalTok{])  }\CommentTok{\# "yellow" (last item)}
\BuiltInTok{print}\NormalTok{(}\BuiltInTok{len}\NormalTok{(colors)) }\CommentTok{\# 4}
\end{Highlighting}
\end{Shaded}

\textbf{Zero-based indexing:} The first item is at index \texttt{0}.

\subsubsection{Modifying Lists}\label{modifying-lists}

\begin{Shaded}
\begin{Highlighting}[]
\NormalTok{fruits }\OperatorTok{=}\NormalTok{ [}\StringTok{"apple"}\NormalTok{, }\StringTok{"banana"}\NormalTok{, }\StringTok{"cherry"}\NormalTok{]}
\NormalTok{fruits[}\DecValTok{1}\NormalTok{] }\OperatorTok{=} \StringTok{"grape"}  \CommentTok{\# Change "banana" to "grape"}
\BuiltInTok{print}\NormalTok{(fruits)         }\CommentTok{\# ["apple", "grape", "cherry"]}
\end{Highlighting}
\end{Shaded}

\subsubsection{Activities}\label{activities-9}

\textbf{Activity 3.3.1: My Top 3 Hobbies}

\begin{Shaded}
\begin{Highlighting}[]
\NormalTok{hobbies }\OperatorTok{=}\NormalTok{ [}\StringTok{"reading"}\NormalTok{, }\StringTok{"coding"}\NormalTok{, }\StringTok{"hiking"}\NormalTok{]}
\BuiltInTok{print}\NormalTok{(}\StringTok{"My hobbies:"}\NormalTok{, hobbies)}
\BuiltInTok{print}\NormalTok{(}\StringTok{"My second hobby:"}\NormalTok{, hobbies[}\DecValTok{1}\NormalTok{])}
\NormalTok{hobbies[}\DecValTok{2}\NormalTok{] }\OperatorTok{=} \StringTok{"swimming"}
\BuiltInTok{print}\NormalTok{(}\StringTok{"Updated hobbies:"}\NormalTok{, hobbies)}
\end{Highlighting}
\end{Shaded}

\subsubsection{Skills Checklist}\label{skills-checklist-14}

\begin{itemize}
\tightlist
\item[$\square$]
  I can create a list in Python.
\item[$\square$]
  I can access individual items in a list using their index.
\item[$\square$]
  I can access list items from the end using negative indices.
\item[$\square$]
  I can change the value of an item at a specific index.
\end{itemize}

\begin{center}\rule{0.5\linewidth}{0.5pt}\end{center}

\subsection{\texorpdfstring{Lesson 3.4: Iterating with \texttt{for}
Loops}{Lesson 3.4: Iterating with for Loops}}\label{lesson-3.4-iterating-with-for-loops}

\subsubsection{Reading References}\label{reading-references-15}

\begin{itemize}
\tightlist
\item
  Zelle: Chapter 2, pages 34--36
\item
  Matthes: Chapter 4, pages 49--60
\end{itemize}

\subsubsection{\texorpdfstring{The \texttt{for} Loop
Structure}{The for Loop Structure}}\label{the-for-loop-structure}

\begin{Shaded}
\begin{Highlighting}[]
\ControlFlowTok{for}\NormalTok{ item }\KeywordTok{in}\NormalTok{ sequence:}
    \CommentTok{\# code that runs for each item}
\end{Highlighting}
\end{Shaded}

\subsubsection{Examples}\label{examples-8}

\textbf{Example 3.4.1: Looping Through a List}

\begin{Shaded}
\begin{Highlighting}[]
\NormalTok{fruits }\OperatorTok{=}\NormalTok{ [}\StringTok{"apple"}\NormalTok{, }\StringTok{"banana"}\NormalTok{, }\StringTok{"cherry"}\NormalTok{]}
\ControlFlowTok{for}\NormalTok{ fruit }\KeywordTok{in}\NormalTok{ fruits:}
    \BuiltInTok{print}\NormalTok{(}\SpecialStringTok{f"I like }\SpecialCharTok{\{}\NormalTok{fruit}\SpecialCharTok{\}}\SpecialStringTok{."}\NormalTok{)}
\BuiltInTok{print}\NormalTok{(}\StringTok{"Finished listing fruits."}\NormalTok{)}
\end{Highlighting}
\end{Shaded}

\textbf{Example 3.4.2: Looping Through a String}

\begin{Shaded}
\begin{Highlighting}[]
\NormalTok{word }\OperatorTok{=} \StringTok{"Python"}
\ControlFlowTok{for}\NormalTok{ char }\KeywordTok{in}\NormalTok{ word:}
    \BuiltInTok{print}\NormalTok{(char)}
\end{Highlighting}
\end{Shaded}

\textbf{Example 3.4.3: Using \texttt{range()}}

\begin{Shaded}
\begin{Highlighting}[]
\ControlFlowTok{for}\NormalTok{ i }\KeywordTok{in} \BuiltInTok{range}\NormalTok{(}\DecValTok{5}\NormalTok{):      }\CommentTok{\# 0, 1, 2, 3, 4}
    \BuiltInTok{print}\NormalTok{(i)}

\ControlFlowTok{for}\NormalTok{ i }\KeywordTok{in} \BuiltInTok{range}\NormalTok{(}\DecValTok{2}\NormalTok{, }\DecValTok{8}\NormalTok{):   }\CommentTok{\# 2, 3, 4, 5, 6, 7}
    \BuiltInTok{print}\NormalTok{(i)}

\ControlFlowTok{for}\NormalTok{ i }\KeywordTok{in} \BuiltInTok{range}\NormalTok{(}\DecValTok{0}\NormalTok{, }\DecValTok{10}\NormalTok{, }\DecValTok{2}\NormalTok{):  }\CommentTok{\# 0, 2, 4, 6, 8}
    \BuiltInTok{print}\NormalTok{(i)}
\end{Highlighting}
\end{Shaded}

\subsubsection{Activities}\label{activities-10}

\textbf{Activity 3.4.1: My Favorite Things}

\begin{Shaded}
\begin{Highlighting}[]
\NormalTok{favorite\_things }\OperatorTok{=}\NormalTok{ [}\StringTok{"sunshine"}\NormalTok{, }\StringTok{"coffee"}\NormalTok{, }\StringTok{"books"}\NormalTok{, }\StringTok{"music"}\NormalTok{, }\StringTok{"mountains"}\NormalTok{]}
\ControlFlowTok{for}\NormalTok{ thing }\KeywordTok{in}\NormalTok{ favorite\_things:}
    \BuiltInTok{print}\NormalTok{(}\SpecialStringTok{f"I love }\SpecialCharTok{\{}\NormalTok{thing}\SpecialCharTok{\}}\SpecialStringTok{!"}\NormalTok{)}
\end{Highlighting}
\end{Shaded}

\textbf{Activity 3.4.2: Countdown Timer}

\begin{Shaded}
\begin{Highlighting}[]
\NormalTok{start }\OperatorTok{=} \BuiltInTok{int}\NormalTok{(}\BuiltInTok{input}\NormalTok{(}\StringTok{"Enter a starting number: "}\NormalTok{))}
\ControlFlowTok{for}\NormalTok{ i }\KeywordTok{in} \BuiltInTok{range}\NormalTok{(start, }\DecValTok{0}\NormalTok{, }\OperatorTok{{-}}\DecValTok{1}\NormalTok{):}
    \BuiltInTok{print}\NormalTok{(i)}
\BuiltInTok{print}\NormalTok{(}\StringTok{"Blast off!"}\NormalTok{)}
\end{Highlighting}
\end{Shaded}

\subsubsection{Skills Checklist}\label{skills-checklist-15}

\begin{itemize}
\tightlist
\item[$\square$]
  I can explain the purpose of a \texttt{for} loop.
\item[$\square$]
  I can write a \texttt{for} loop to iterate over items in a list.
\item[$\square$]
  I can write a \texttt{for} loop to iterate over characters in a
  string.
\item[$\square$]
  I can use the \texttt{range()} function to control how many times a
  loop runs.
\end{itemize}

\begin{center}\rule{0.5\linewidth}{0.5pt}\end{center}

\subsection{Lesson 3.5: Basic String
Operations}\label{lesson-3.5-basic-string-operations}

\subsubsection{Reading References}\label{reading-references-16}

\begin{itemize}
\tightlist
\item
  Zelle: Chapter 5, pages 99--115
\item
  Matthes: Chapter 2, pages 19--25
\end{itemize}

\subsubsection{Common String Methods}\label{common-string-methods}

\begin{Shaded}
\begin{Highlighting}[]
\NormalTok{message }\OperatorTok{=} \StringTok{"  Hello Python!  "}

\NormalTok{message.strip()          }\CommentTok{\# "Hello Python!"   {-}{-}{-} removes whitespace}
\NormalTok{message.upper()          }\CommentTok{\# "  HELLO PYTHON!  "}
\NormalTok{message.lower()          }\CommentTok{\# "  hello python!  "}
\NormalTok{message.replace(}\StringTok{"Hello"}\NormalTok{, }\StringTok{"Hi"}\NormalTok{)  }\CommentTok{\# "  Hi Python!  "}
\BuiltInTok{len}\NormalTok{(message)             }\CommentTok{\# 17}
\end{Highlighting}
\end{Shaded}

\textbf{Important:} String methods return a \textbf{new string} --- they
do not modify the original.

\subsubsection{Examples}\label{examples-9}

\textbf{Example 3.5.1: String Manipulation}

\begin{Shaded}
\begin{Highlighting}[]
\NormalTok{message }\OperatorTok{=} \StringTok{"  Hello Python  "}
\BuiltInTok{print}\NormalTok{(}\SpecialStringTok{f"Original: \textquotesingle{}}\SpecialCharTok{\{}\NormalTok{message}\SpecialCharTok{\}}\SpecialStringTok{\textquotesingle{}"}\NormalTok{)}
\BuiltInTok{print}\NormalTok{(}\SpecialStringTok{f"Length: }\SpecialCharTok{\{}\BuiltInTok{len}\NormalTok{(message)}\SpecialCharTok{\}}\SpecialStringTok{"}\NormalTok{)}
\BuiltInTok{print}\NormalTok{(}\SpecialStringTok{f"Uppercase: \textquotesingle{}}\SpecialCharTok{\{}\NormalTok{message}\SpecialCharTok{.}\NormalTok{upper()}\SpecialCharTok{\}}\SpecialStringTok{\textquotesingle{}"}\NormalTok{)}
\BuiltInTok{print}\NormalTok{(}\SpecialStringTok{f"Stripped: \textquotesingle{}}\SpecialCharTok{\{}\NormalTok{message}\SpecialCharTok{.}\NormalTok{strip()}\SpecialCharTok{\}}\SpecialStringTok{\textquotesingle{}"}\NormalTok{)}
\end{Highlighting}
\end{Shaded}

\textbf{Example 3.5.2: Combining Strings}

\begin{Shaded}
\begin{Highlighting}[]
\NormalTok{first\_name }\OperatorTok{=} \StringTok{"Alice"}
\NormalTok{last\_name }\OperatorTok{=} \StringTok{"Smith"}
\NormalTok{full\_name }\OperatorTok{=}\NormalTok{ first\_name }\OperatorTok{+} \StringTok{" "} \OperatorTok{+}\NormalTok{ last\_name}
\BuiltInTok{print}\NormalTok{(}\SpecialStringTok{f"Full name: }\SpecialCharTok{\{}\NormalTok{full\_name}\SpecialCharTok{\}}\SpecialStringTok{"}\NormalTok{)}
\end{Highlighting}
\end{Shaded}

\subsubsection{Activities}\label{activities-11}

\textbf{Activity 3.5.1: Name Formatter}

\begin{Shaded}
\begin{Highlighting}[]
\NormalTok{first\_name }\OperatorTok{=} \BuiltInTok{input}\NormalTok{(}\StringTok{"Enter your first name: "}\NormalTok{)}
\NormalTok{last\_name }\OperatorTok{=} \BuiltInTok{input}\NormalTok{(}\StringTok{"Enter your last name: "}\NormalTok{)}
\NormalTok{full\_name }\OperatorTok{=}\NormalTok{ first\_name }\OperatorTok{+} \StringTok{" "} \OperatorTok{+}\NormalTok{ last\_name}
\BuiltInTok{print}\NormalTok{(}\SpecialStringTok{f"Full name: }\SpecialCharTok{\{}\NormalTok{full\_name}\SpecialCharTok{\}}\SpecialStringTok{"}\NormalTok{)}
\BuiltInTok{print}\NormalTok{(}\SpecialStringTok{f"Uppercase: }\SpecialCharTok{\{}\NormalTok{full\_name}\SpecialCharTok{.}\NormalTok{upper()}\SpecialCharTok{\}}\SpecialStringTok{"}\NormalTok{)}
\BuiltInTok{print}\NormalTok{(}\SpecialStringTok{f"Total characters: }\SpecialCharTok{\{}\BuiltInTok{len}\NormalTok{(full\_name)}\SpecialCharTok{\}}\SpecialStringTok{"}\NormalTok{)}
\end{Highlighting}
\end{Shaded}

\subsubsection{Skills Checklist}\label{skills-checklist-16}

\begin{itemize}
\tightlist
\item[$\square$]
  I can combine strings using concatenation.
\item[$\square$]
  I can find the length of a string using \texttt{len()}.
\item[$\square$]
  I can convert a string to uppercase/lowercase.
\item[$\square$]
  I can remove extra spaces using \texttt{.strip()}.
\item[$\square$]
  I can replace parts of a string using \texttt{.replace()}.
\end{itemize}

\newpage

\section{Unit 4: Data Structures --- Lists and
Dictionaries}\label{unit-4-data-structures-lists-and-dictionaries-1}

\subsection{Overview}\label{overview-4}

This unit deepens understanding of lists and introduces dictionaries as
powerful mechanisms for organizing data.

\textbf{Objectives:}

\begin{itemize}
\tightlist
\item
  Add, remove, and organize list elements.
\item
  Work with tuples (immutable lists).
\item
  Use dictionaries for key-value pair storage.
\item
  Loop through dictionaries.
\end{itemize}

\begin{center}\rule{0.5\linewidth}{0.5pt}\end{center}

\subsection{Lesson 4.1: More on Lists}\label{lesson-4.1-more-on-lists}

\subsubsection{Adding and Removing
Elements}\label{adding-and-removing-elements}

\begin{Shaded}
\begin{Highlighting}[]
\NormalTok{my\_list }\OperatorTok{=}\NormalTok{ [}\DecValTok{10}\NormalTok{, }\DecValTok{20}\NormalTok{, }\DecValTok{30}\NormalTok{]}

\NormalTok{my\_list.append(}\DecValTok{40}\NormalTok{)          }\CommentTok{\# Add to end {-}\textgreater{} [10, 20, 30, 40]}
\NormalTok{my\_list.insert(}\DecValTok{1}\NormalTok{, }\DecValTok{15}\NormalTok{)       }\CommentTok{\# Insert at index 1 {-}\textgreater{} [10, 15, 20, 30, 40]}
\KeywordTok{del}\NormalTok{ my\_list[}\DecValTok{0}\NormalTok{]              }\CommentTok{\# Delete by index {-}\textgreater{} [15, 20, 30, 40]}
\NormalTok{removed }\OperatorTok{=}\NormalTok{ my\_list.pop()     }\CommentTok{\# Remove and return last {-}\textgreater{} removed=40, list=[15, 20, 30]}
\NormalTok{my\_list.remove(}\DecValTok{20}\NormalTok{)          }\CommentTok{\# Remove by value {-}\textgreater{} [15, 30]}
\end{Highlighting}
\end{Shaded}

\subsubsection{Example}\label{example}

\begin{Shaded}
\begin{Highlighting}[]
\NormalTok{todo\_list }\OperatorTok{=}\NormalTok{ [}\StringTok{"Buy groceries"}\NormalTok{, }\StringTok{"Walk the dog"}\NormalTok{]}

\NormalTok{todo\_list.append(}\StringTok{"Call grandma"}\NormalTok{)}
\NormalTok{todo\_list.insert(}\DecValTok{0}\NormalTok{, }\StringTok{"Wake up early"}\NormalTok{)}
\KeywordTok{del}\NormalTok{ todo\_list[}\DecValTok{1}\NormalTok{]}
\NormalTok{completed\_task }\OperatorTok{=}\NormalTok{ todo\_list.pop()}
\NormalTok{todo\_list.remove(}\StringTok{"Walk the dog"}\NormalTok{)}
\BuiltInTok{print}\NormalTok{(todo\_list)}
\end{Highlighting}
\end{Shaded}

\subsubsection{Skills Checklist}\label{skills-checklist-17}

\begin{itemize}
\tightlist
\item[$\square$]
  I can add an item to the end of a list using \texttt{.append()}.
\item[$\square$]
  I can insert an item at a specific position using \texttt{.insert()}.
\item[$\square$]
  I can remove an item by index using \texttt{del}.
\item[$\square$]
  I can remove an item and retrieve its value using \texttt{.pop()}.
\item[$\square$]
  I can remove the first occurrence of a value using \texttt{.remove()}.
\end{itemize}

\begin{center}\rule{0.5\linewidth}{0.5pt}\end{center}

\subsection{Lesson 4.2: Organizing
Lists}\label{lesson-4.2-organizing-lists}

\subsubsection{Sorting and Reversing}\label{sorting-and-reversing}

\begin{Shaded}
\begin{Highlighting}[]
\NormalTok{numbers }\OperatorTok{=}\NormalTok{ [}\DecValTok{5}\NormalTok{, }\DecValTok{2}\NormalTok{, }\DecValTok{9}\NormalTok{, }\DecValTok{1}\NormalTok{]}
\NormalTok{numbers.sort()                          }\CommentTok{\# Permanent sort in{-}place}
\BuiltInTok{print}\NormalTok{(numbers)                          }\CommentTok{\# [1, 2, 5, 9]}

\NormalTok{fruits }\OperatorTok{=}\NormalTok{ [}\StringTok{"orange"}\NormalTok{, }\StringTok{"apple"}\NormalTok{, }\StringTok{"banana"}\NormalTok{]}
\NormalTok{sorted\_fruits }\OperatorTok{=} \BuiltInTok{sorted}\NormalTok{(fruits)          }\CommentTok{\# Temporary sort {-}{-}{-} original unchanged}
\BuiltInTok{print}\NormalTok{(fruits)                           }\CommentTok{\# [\textquotesingle{}orange\textquotesingle{}, \textquotesingle{}apple\textquotesingle{}, \textquotesingle{}banana\textquotesingle{}]}
\BuiltInTok{print}\NormalTok{(sorted\_fruits)                    }\CommentTok{\# [\textquotesingle{}apple\textquotesingle{}, \textquotesingle{}banana\textquotesingle{}, \textquotesingle{}orange\textquotesingle{}]}

\NormalTok{numbers.reverse()                       }\CommentTok{\# Reverse in{-}place}
\NormalTok{numbers.sort(reverse}\OperatorTok{=}\VariableTok{True}\NormalTok{)              }\CommentTok{\# Sort descending}
\BuiltInTok{print}\NormalTok{(}\BuiltInTok{len}\NormalTok{(numbers))                     }\CommentTok{\# Length}
\end{Highlighting}
\end{Shaded}

\subsubsection{Skills Checklist}\label{skills-checklist-18}

\begin{itemize}
\tightlist
\item[$\square$]
  I can sort a list permanently using \texttt{.sort()}.
\item[$\square$]
  I can get a temporarily sorted version using \texttt{sorted()}.
\item[$\square$]
  I can reverse a list using \texttt{.reverse()}.
\item[$\square$]
  I can find how many items are in a list using \texttt{len()}.
\end{itemize}

\begin{center}\rule{0.5\linewidth}{0.5pt}\end{center}

\subsection{Lesson 4.3: Tuples --- Immutable
Lists}\label{lesson-4.3-tuples-immutable-lists}

\subsubsection{Tuples}\label{tuples}

Tuples are defined with parentheses and \textbf{cannot be changed} after
creation.

\begin{Shaded}
\begin{Highlighting}[]
\NormalTok{coordinates }\OperatorTok{=}\NormalTok{ (}\DecValTok{10}\NormalTok{, }\DecValTok{20}\NormalTok{, }\DecValTok{30}\NormalTok{)}
\BuiltInTok{print}\NormalTok{(coordinates[}\DecValTok{0}\NormalTok{])     }\CommentTok{\# 10}
\BuiltInTok{print}\NormalTok{(coordinates[}\OperatorTok{{-}}\DecValTok{1}\NormalTok{])    }\CommentTok{\# 30}

\CommentTok{\# This would cause a TypeError:}
\CommentTok{\# coordinates[1] = 25}
\end{Highlighting}
\end{Shaded}

\textbf{When to use tuples:} For data that should remain constant
(coordinates, RGB values, fixed settings).

\subsubsection{Skills Checklist}\label{skills-checklist-19}

\begin{itemize}
\tightlist
\item[$\square$]
  I can create a tuple in Python.
\item[$\square$]
  I can explain the main difference between a list and a tuple
  (mutability).
\item[$\square$]
  I can access items in a tuple by their index.
\end{itemize}

\begin{center}\rule{0.5\linewidth}{0.5pt}\end{center}

\subsection{Lesson 4.4: Introducing
Dictionaries}\label{lesson-4.4-introducing-dictionaries}

\subsubsection{Reading References}\label{reading-references-17}

\begin{itemize}
\tightlist
\item
  Zelle: Chapter 11, pages 307--309
\item
  Matthes: Chapter 6, pages 91--98
\end{itemize}

\subsubsection{Dictionary Basics}\label{dictionary-basics}

A dictionary stores \textbf{key-value pairs}. Keys must be unique and
immutable (strings, numbers, tuples). Values can be anything.

\begin{Shaded}
\begin{Highlighting}[]
\NormalTok{student\_profile }\OperatorTok{=}\NormalTok{ \{}
    \StringTok{"name"}\NormalTok{: }\StringTok{"Emily"}\NormalTok{,}
    \StringTok{"age"}\NormalTok{: }\DecValTok{13}\NormalTok{,}
    \StringTok{"grade"}\NormalTok{: }\StringTok{"7th"}\NormalTok{,}
    \StringTok{"favorite\_subject"}\NormalTok{: }\StringTok{"Science"}
\NormalTok{\}}

\BuiltInTok{print}\NormalTok{(student\_profile[}\StringTok{"name"}\NormalTok{])      }\CommentTok{\# "Emily"}
\NormalTok{student\_profile[}\StringTok{"city"}\NormalTok{] }\OperatorTok{=} \StringTok{"Springfield"}   \CommentTok{\# Add new key}
\NormalTok{student\_profile[}\StringTok{"age"}\NormalTok{] }\OperatorTok{=} \DecValTok{14}              \CommentTok{\# Modify existing key}
\KeywordTok{del}\NormalTok{ student\_profile[}\StringTok{"grade"}\NormalTok{]             }\CommentTok{\# Remove key}
\end{Highlighting}
\end{Shaded}

\subsubsection{Examples}\label{examples-10}

\textbf{Example 4.4.1: Book Information Dictionary}

\begin{Shaded}
\begin{Highlighting}[]
\NormalTok{my\_book }\OperatorTok{=}\NormalTok{ \{}
    \StringTok{"title"}\NormalTok{: }\StringTok{"The Hitchhiker\textquotesingle{}s Guide to the Galaxy"}\NormalTok{,}
    \StringTok{"author"}\NormalTok{: }\StringTok{"Douglas Adams"}\NormalTok{,}
    \StringTok{"publication\_year"}\NormalTok{: }\DecValTok{1979}\NormalTok{,}
    \StringTok{"genre"}\NormalTok{: }\StringTok{"Science Fiction"}
\NormalTok{\}}

\BuiltInTok{print}\NormalTok{(}\SpecialStringTok{f"Title: }\SpecialCharTok{\{}\NormalTok{my\_book[}\StringTok{\textquotesingle{}title\textquotesingle{}}\NormalTok{]}\SpecialCharTok{\}}\SpecialStringTok{"}\NormalTok{)}
\BuiltInTok{print}\NormalTok{(}\SpecialStringTok{f"Author: }\SpecialCharTok{\{}\NormalTok{my\_book[}\StringTok{\textquotesingle{}author\textquotesingle{}}\NormalTok{]}\SpecialCharTok{\}}\SpecialStringTok{"}\NormalTok{)}

\NormalTok{my\_book[}\StringTok{"number\_of\_pages"}\NormalTok{] }\OperatorTok{=} \DecValTok{193}
\NormalTok{my\_book[}\StringTok{"genre"}\NormalTok{] }\OperatorTok{=} \StringTok{"Comedy Science Fiction"}
\BuiltInTok{print}\NormalTok{(my\_book)}
\end{Highlighting}
\end{Shaded}

\subsubsection{Skills Checklist}\label{skills-checklist-20}

\begin{itemize}
\tightlist
\item[$\square$]
  I can create a dictionary in Python.
\item[$\square$]
  I can explain what a key-value pair is.
\item[$\square$]
  I can access a value in a dictionary using its key.
\item[$\square$]
  I can add, modify, and remove key-value pairs.
\end{itemize}

\begin{center}\rule{0.5\linewidth}{0.5pt}\end{center}

\subsection{Lesson 4.5: Looping Through
Dictionaries}\label{lesson-4.5-looping-through-dictionaries}

\subsubsection{Dictionary Iteration
Methods}\label{dictionary-iteration-methods}

\begin{Shaded}
\begin{Highlighting}[]
\NormalTok{inventory }\OperatorTok{=}\NormalTok{ \{}\StringTok{"apples"}\NormalTok{: }\DecValTok{50}\NormalTok{, }\StringTok{"bananas"}\NormalTok{: }\DecValTok{30}\NormalTok{, }\StringTok{"oranges"}\NormalTok{: }\DecValTok{40}\NormalTok{\}}

\CommentTok{\# Loop through keys (default)}
\ControlFlowTok{for}\NormalTok{ item\_name }\KeywordTok{in}\NormalTok{ inventory:}
    \BuiltInTok{print}\NormalTok{(item\_name)}

\CommentTok{\# Loop through values only}
\ControlFlowTok{for}\NormalTok{ count }\KeywordTok{in}\NormalTok{ inventory.values():}
    \BuiltInTok{print}\NormalTok{(count)}

\CommentTok{\# Loop through key{-}value pairs}
\ControlFlowTok{for}\NormalTok{ item\_name, count }\KeywordTok{in}\NormalTok{ inventory.items():}
    \BuiltInTok{print}\NormalTok{(}\SpecialStringTok{f"}\SpecialCharTok{\{}\NormalTok{item\_name}\SpecialCharTok{\}}\SpecialStringTok{: }\SpecialCharTok{\{}\NormalTok{count}\SpecialCharTok{\}}\SpecialStringTok{ in stock"}\NormalTok{)}
\end{Highlighting}
\end{Shaded}

\subsubsection{Skills Checklist}\label{skills-checklist-21}

\begin{itemize}
\tightlist
\item[$\square$]
  I can loop through all the keys in a dictionary.
\item[$\square$]
  I can loop through all the values in a dictionary.
\item[$\square$]
  I can loop through both keys and values at the same time.
\end{itemize}

\begin{center}\rule{0.5\linewidth}{0.5pt}\end{center}

\section{\texorpdfstring{Unit 5: User Input and \texttt{while}
Loops}{Unit 5: User Input and while Loops}}\label{unit-5-user-input-and-while-loops-1}

\subsection{Overview}\label{overview-5}

This unit focuses on obtaining diverse input from users and controlling
program execution through indefinite loops.

\textbf{Objectives:}

\begin{itemize}
\tightlist
\item
  Implement indefinite loops using \texttt{while} statements.
\item
  Control \texttt{while} loop execution with \texttt{break} and
  \texttt{continue}.
\item
  Avoid infinite loops.
\end{itemize}

\begin{center}\rule{0.5\linewidth}{0.5pt}\end{center}

\subsection{\texorpdfstring{Lesson 5.1: How the \texttt{input()}
Function
Works}{Lesson 5.1: How the input() Function Works}}\label{lesson-5.1-how-the-input-function-works}

\subsubsection{Key Reminder}\label{key-reminder}

\texttt{input()} \textbf{always returns a string}. Convert to numeric
types as needed.

\begin{Shaded}
\begin{Highlighting}[]
\NormalTok{age\_input }\OperatorTok{=} \BuiltInTok{input}\NormalTok{(}\StringTok{"How old are you? "}\NormalTok{)}
\BuiltInTok{print}\NormalTok{(}\SpecialStringTok{f"Type: }\SpecialCharTok{\{}\BuiltInTok{type}\NormalTok{(age\_input)}\SpecialCharTok{\}}\SpecialStringTok{"}\NormalTok{)    }\CommentTok{\# \textless{}class \textquotesingle{}str\textquotesingle{}\textgreater{}}

\NormalTok{age\_int }\OperatorTok{=} \BuiltInTok{int}\NormalTok{(age\_input)}
\BuiltInTok{print}\NormalTok{(}\SpecialStringTok{f"In 5 years: }\SpecialCharTok{\{}\NormalTok{age\_int }\OperatorTok{+} \DecValTok{5}\SpecialCharTok{\}}\SpecialStringTok{"}\NormalTok{)}
\end{Highlighting}
\end{Shaded}

\subsubsection{Skills Checklist}\label{skills-checklist-22}

\begin{itemize}
\tightlist
\item[$\square$]
  I can effectively use \texttt{input()} to get user data.
\item[$\square$]
  I can confirm that \texttt{input()}'s return value is always a string.
\item[$\square$]
  I can convert user input to a numeric type when necessary.
\end{itemize}

\begin{center}\rule{0.5\linewidth}{0.5pt}\end{center}

\subsection{\texorpdfstring{Lesson 5.2: The \texttt{while} Loop in
Action}{Lesson 5.2: The while Loop in Action}}\label{lesson-5.2-the-while-loop-in-action}

\subsubsection{\texorpdfstring{\texttt{for} vs.~\texttt{while} Loop
Comparison}{for vs.~while Loop Comparison}}\label{for-vs.-while-loop-comparison}

{\def\LTcaptype{none} % do not increment counter
\begin{longtable}[]{@{}
  >{\raggedright\arraybackslash}p{(\linewidth - 4\tabcolsep) * \real{0.2128}}
  >{\raggedright\arraybackslash}p{(\linewidth - 4\tabcolsep) * \real{0.3404}}
  >{\raggedright\arraybackslash}p{(\linewidth - 4\tabcolsep) * \real{0.4468}}@{}}
\toprule\noalign{}
\begin{minipage}[b]{\linewidth}\raggedright
Feature
\end{minipage} & \begin{minipage}[b]{\linewidth}\raggedright
\texttt{for} Loop
\end{minipage} & \begin{minipage}[b]{\linewidth}\raggedright
\texttt{while} Loop
\end{minipage} \\
\midrule\noalign{}
\endhead
\bottomrule\noalign{}
\endlastfoot
\textbf{Purpose} & Definite iteration (known count) & Indefinite
iteration (condition-based) \\
\textbf{When to Use} & Iterating over lists, ranges & Repeating until a
condition is met \\
\textbf{Control Variable} & Automatically updated & Must be manually
updated \\
\textbf{Risk} & Less prone to infinite loops & High risk if condition
never becomes false \\
\end{longtable}
}

\subsubsection{Examples}\label{examples-11}

\textbf{Example 5.2.1: Simple Counter}

\begin{Shaded}
\begin{Highlighting}[]
\NormalTok{count }\OperatorTok{=} \DecValTok{1}
\ControlFlowTok{while}\NormalTok{ count }\OperatorTok{\textless{}=} \DecValTok{5}\NormalTok{:}
    \BuiltInTok{print}\NormalTok{(count)}
\NormalTok{    count }\OperatorTok{+=} \DecValTok{1}
\BuiltInTok{print}\NormalTok{(}\StringTok{"Counting finished!"}\NormalTok{)}
\end{Highlighting}
\end{Shaded}

\textbf{Example 5.2.2: User Input Loop (Sentinel Value)}

\begin{Shaded}
\begin{Highlighting}[]
\NormalTok{message }\OperatorTok{=} \StringTok{""}
\ControlFlowTok{while}\NormalTok{ message }\OperatorTok{!=} \StringTok{\textquotesingle{}quit\textquotesingle{}}\NormalTok{:}
\NormalTok{    message }\OperatorTok{=} \BuiltInTok{input}\NormalTok{(}\StringTok{"Enter a message, or type \textquotesingle{}quit\textquotesingle{} to exit: "}\NormalTok{)}
    \ControlFlowTok{if}\NormalTok{ message }\OperatorTok{!=} \StringTok{\textquotesingle{}quit\textquotesingle{}}\NormalTok{:}
        \BuiltInTok{print}\NormalTok{(}\SpecialStringTok{f"You said: }\SpecialCharTok{\{}\NormalTok{message}\SpecialCharTok{\}}\SpecialStringTok{"}\NormalTok{)}
\BuiltInTok{print}\NormalTok{(}\StringTok{"Goodbye!"}\NormalTok{)}
\end{Highlighting}
\end{Shaded}

\subsubsection{Skills Checklist}\label{skills-checklist-23}

\begin{itemize}
\tightlist
\item[$\square$]
  I can explain the difference between a \texttt{for} loop and a
  \texttt{while} loop.
\item[$\square$]
  I can write a basic \texttt{while} loop.
\item[$\square$]
  I can make sure my \texttt{while} loop has a way to stop.
\end{itemize}

\begin{center}\rule{0.5\linewidth}{0.5pt}\end{center}

\subsection{\texorpdfstring{Lesson 5.3: Using Flags and \texttt{break}
to Exit
Loops}{Lesson 5.3: Using Flags and break to Exit Loops}}\label{lesson-5.3-using-flags-and-break-to-exit-loops}

\subsubsection{Flag Variable}\label{flag-variable}

\begin{Shaded}
\begin{Highlighting}[]
\NormalTok{active }\OperatorTok{=} \VariableTok{True}
\ControlFlowTok{while}\NormalTok{ active:}
\NormalTok{    message }\OperatorTok{=} \BuiltInTok{input}\NormalTok{(}\StringTok{"Say something (or \textquotesingle{}quit\textquotesingle{} to stop): "}\NormalTok{)}
    \ControlFlowTok{if}\NormalTok{ message }\OperatorTok{==} \StringTok{\textquotesingle{}quit\textquotesingle{}}\NormalTok{:}
\NormalTok{        active }\OperatorTok{=} \VariableTok{False}
    \ControlFlowTok{else}\NormalTok{:}
        \BuiltInTok{print}\NormalTok{(}\SpecialStringTok{f"You said: }\SpecialCharTok{\{}\NormalTok{message}\SpecialCharTok{\}}\SpecialStringTok{"}\NormalTok{)}
\end{Highlighting}
\end{Shaded}

\subsubsection{\texorpdfstring{\texttt{break}
Statement}{break Statement}}\label{break-statement}

\begin{Shaded}
\begin{Highlighting}[]
\ControlFlowTok{while} \VariableTok{True}\NormalTok{:}
\NormalTok{    city }\OperatorTok{=} \BuiltInTok{input}\NormalTok{(}\StringTok{"Enter a city (or \textquotesingle{}done\textquotesingle{} to finish): "}\NormalTok{)}
    \ControlFlowTok{if}\NormalTok{ city }\OperatorTok{==} \StringTok{\textquotesingle{}done\textquotesingle{}}\NormalTok{:}
        \ControlFlowTok{break}
    \BuiltInTok{print}\NormalTok{(}\SpecialStringTok{f"I\textquotesingle{}ve heard }\SpecialCharTok{\{}\NormalTok{city}\SpecialCharTok{\}}\SpecialStringTok{ is a great place!"}\NormalTok{)}
\end{Highlighting}
\end{Shaded}

\subsubsection{Skills Checklist}\label{skills-checklist-24}

\begin{itemize}
\tightlist
\item[$\square$]
  I can use a Boolean ``flag'' variable to control a \texttt{while}
  loop.
\item[$\square$]
  I can use the \texttt{break} statement to exit a loop early.
\end{itemize}

\begin{center}\rule{0.5\linewidth}{0.5pt}\end{center}

\subsection{\texorpdfstring{Lesson 5.4: Using \texttt{continue} in
Loops}{Lesson 5.4: Using continue in Loops}}\label{lesson-5.4-using-continue-in-loops}

\subsubsection{\texorpdfstring{\texttt{continue} --- Skip the Current
Iteration}{continue --- Skip the Current Iteration}}\label{continue-skip-the-current-iteration}

\begin{Shaded}
\begin{Highlighting}[]
\NormalTok{number }\OperatorTok{=} \DecValTok{0}
\ControlFlowTok{while}\NormalTok{ number }\OperatorTok{\textless{}} \DecValTok{10}\NormalTok{:}
\NormalTok{    number }\OperatorTok{+=} \DecValTok{1}
    \ControlFlowTok{if}\NormalTok{ number }\OperatorTok{\%} \DecValTok{2} \OperatorTok{==} \DecValTok{0}\NormalTok{:}
        \ControlFlowTok{continue}        \CommentTok{\# Skip even numbers}
    \BuiltInTok{print}\NormalTok{(number)       }\CommentTok{\# Prints 1, 3, 5, 7, 9}
\end{Highlighting}
\end{Shaded}

\subsubsection{\texorpdfstring{\texttt{continue}
vs.~\texttt{break}}{continue vs.~break}}\label{continue-vs.-break}

\begin{itemize}
\tightlist
\item
  \texttt{break} --- \textbf{exits the loop entirely}
\item
  \texttt{continue} --- \textbf{skips only the current iteration}, loop
  continues
\end{itemize}

\subsubsection{Skills Checklist}\label{skills-checklist-25}

\begin{itemize}
\tightlist
\item[$\square$]
  I can use \texttt{continue} to skip the rest of the current loop
  iteration.
\item[$\square$]
  I can identify scenarios where \texttt{continue} would be useful.
\end{itemize}

\begin{center}\rule{0.5\linewidth}{0.5pt}\end{center}

\subsection{Lesson 5.5: Avoiding Infinite
Loops}\label{lesson-5.5-avoiding-infinite-loops}

An \textbf{infinite loop} never ends because its condition always
remains \texttt{True}.

\textbf{Common causes:}

\begin{enumerate}
\def\labelenumi{\arabic{enumi}.}
\tightlist
\item
  Forgetting to update the loop variable.
\item
  Updating the variable in the wrong direction.
\item
  Using an incorrect logical condition.
\end{enumerate}

\begin{Shaded}
\begin{Highlighting}[]
\CommentTok{\# INFINITE LOOP {-}{-}{-} Don\textquotesingle{}t run!}
\NormalTok{count }\OperatorTok{=} \DecValTok{1}
\ControlFlowTok{while}\NormalTok{ count }\OperatorTok{\textless{}=} \DecValTok{5}\NormalTok{:}
    \BuiltInTok{print}\NormalTok{(count)}
    \CommentTok{\# Missing: count += 1}

\CommentTok{\# CORRECT}
\NormalTok{count }\OperatorTok{=} \DecValTok{1}
\ControlFlowTok{while}\NormalTok{ count }\OperatorTok{\textless{}=} \DecValTok{5}\NormalTok{:}
    \BuiltInTok{print}\NormalTok{(count)}
\NormalTok{    count }\OperatorTok{+=} \DecValTok{1}
\end{Highlighting}
\end{Shaded}

\textbf{To stop a running infinite loop:} Press \texttt{Ctrl\ +\ C} in
the terminal.

\subsubsection{Skills Checklist}\label{skills-checklist-26}

\begin{itemize}
\tightlist
\item[$\square$]
  I can define what an infinite loop is.
\item[$\square$]
  I can identify common causes of infinite loops.
\item[$\square$]
  I know how to stop a running program in an infinite loop.
\end{itemize}

\begin{center}\rule{0.5\linewidth}{0.5pt}\end{center}

\section{Unit 6: Functions}\label{unit-6-functions-1}

\subsection{Overview}\label{overview-6}

This unit introduces functions, enabling organization of code into
reusable blocks, passing arguments, and returning values.

\textbf{Objectives:}

\begin{itemize}
\tightlist
\item
  Define and call functions.
\item
  Pass arguments to functions.
\item
  Return values from functions.
\item
  Understand the concept of scope.
\end{itemize}

\begin{center}\rule{0.5\linewidth}{0.5pt}\end{center}

\subsection{Lesson 6.1: Defining
Functions}\label{lesson-6.1-defining-functions}

\subsubsection{Reading References}\label{reading-references-18}

\begin{itemize}
\tightlist
\item
  Zelle: Chapter 6, pages 137--142
\item
  Matthes: Chapter 8, pages 129--130
\end{itemize}

\subsubsection{Function Syntax}\label{function-syntax}

\begin{Shaded}
\begin{Highlighting}[]
\KeywordTok{def}\NormalTok{ function\_name(parameter1, parameter2):}
    \CommentTok{"""Docstring describing what the function does."""}
    \CommentTok{\# Function body}
    \BuiltInTok{print}\NormalTok{(}\SpecialStringTok{f"Hello, }\SpecialCharTok{\{}\NormalTok{parameter1}\SpecialCharTok{\}}\SpecialStringTok{!"}\NormalTok{)}
\end{Highlighting}
\end{Shaded}

\subsubsection{Examples}\label{examples-12}

\textbf{Example 6.1.1: Simple Greeting Function}

\begin{Shaded}
\begin{Highlighting}[]
\KeywordTok{def}\NormalTok{ greet\_user():}
    \CommentTok{"""Prints a simple greeting message."""}
    \BuiltInTok{print}\NormalTok{(}\StringTok{"Hello, Python programmer!"}\NormalTok{)}

\NormalTok{greet\_user()   }\CommentTok{\# Call the function}
\NormalTok{greet\_user()   }\CommentTok{\# Call it again}
\end{Highlighting}
\end{Shaded}

\textbf{Example 6.1.2: Function with a Parameter}

\begin{Shaded}
\begin{Highlighting}[]
\KeywordTok{def}\NormalTok{ greet\_name(name):}
    \CommentTok{"""Prints a personalized greeting."""}
    \BuiltInTok{print}\NormalTok{(}\SpecialStringTok{f"Hello, }\SpecialCharTok{\{}\NormalTok{name}\SpecialCharTok{\}}\SpecialStringTok{!"}\NormalTok{)}

\NormalTok{greet\_name(}\StringTok{"Alice"}\NormalTok{)}
\NormalTok{greet\_name(}\StringTok{"Bob"}\NormalTok{)}
\end{Highlighting}
\end{Shaded}

\subsubsection{Skills Checklist}\label{skills-checklist-27}

\begin{itemize}
\tightlist
\item[$\square$]
  I can explain why functions are useful.
\item[$\square$]
  I can define a simple function in Python.
\item[$\square$]
  I can use parameters to specify function inputs.
\item[$\square$]
  I can call a function by its name.
\end{itemize}

\begin{center}\rule{0.5\linewidth}{0.5pt}\end{center}

\subsection{Lesson 6.2: Passing
Arguments}\label{lesson-6.2-passing-arguments}

\subsubsection{Positional vs.~Keyword
Arguments}\label{positional-vs.-keyword-arguments}

\begin{Shaded}
\begin{Highlighting}[]
\KeywordTok{def}\NormalTok{ describe\_animal(animal\_type, animal\_name):}
    \BuiltInTok{print}\NormalTok{(}\SpecialStringTok{f"I have a }\SpecialCharTok{\{}\NormalTok{animal\_type}\SpecialCharTok{\}}\SpecialStringTok{ named }\SpecialCharTok{\{}\NormalTok{animal\_name}\SpecialCharTok{\}}\SpecialStringTok{."}\NormalTok{)}

\CommentTok{\# Positional arguments (order matters)}
\NormalTok{describe\_animal(}\StringTok{"dog"}\NormalTok{, }\StringTok{"Buddy"}\NormalTok{)}

\CommentTok{\# Keyword arguments (order doesn\textquotesingle{}t matter)}
\NormalTok{describe\_animal(animal\_name}\OperatorTok{=}\StringTok{"Max"}\NormalTok{, animal\_type}\OperatorTok{=}\StringTok{"cat"}\NormalTok{)}
\end{Highlighting}
\end{Shaded}

\subsubsection{Default Values}\label{default-values}

\begin{Shaded}
\begin{Highlighting}[]
\KeywordTok{def}\NormalTok{ make\_shirt(size, message}\OperatorTok{=}\StringTok{"I love Python"}\NormalTok{):}
    \BuiltInTok{print}\NormalTok{(}\SpecialStringTok{f"Making a }\SpecialCharTok{\{}\NormalTok{size}\SpecialCharTok{\}}\SpecialStringTok{{-}sized shirt: \textquotesingle{}}\SpecialCharTok{\{}\NormalTok{message}\SpecialCharTok{\}}\SpecialStringTok{\textquotesingle{}"}\NormalTok{)}

\NormalTok{make\_shirt(}\StringTok{"large"}\NormalTok{)                       }\CommentTok{\# Uses default message}
\NormalTok{make\_shirt(}\StringTok{"medium"}\NormalTok{, }\StringTok{"Hello World"}\NormalTok{)       }\CommentTok{\# Custom message}
\end{Highlighting}
\end{Shaded}

\subsubsection{Skills Checklist}\label{skills-checklist-28}

\begin{itemize}
\tightlist
\item[$\square$]
  I can differentiate between a parameter and an argument.
\item[$\square$]
  I can pass arguments using positional order.
\item[$\square$]
  I can pass arguments using keyword arguments.
\end{itemize}

\begin{center}\rule{0.5\linewidth}{0.5pt}\end{center}

\subsection{Lesson 6.3: Return Values}\label{lesson-6.3-return-values}

\subsubsection{\texorpdfstring{The \texttt{return}
Statement}{The return Statement}}\label{the-return-statement}

\begin{Shaded}
\begin{Highlighting}[]
\KeywordTok{def}\NormalTok{ add\_numbers(num1, num2):}
    \CommentTok{"""Adds two numbers and returns the result."""}
\NormalTok{    total }\OperatorTok{=}\NormalTok{ num1 }\OperatorTok{+}\NormalTok{ num2}
    \ControlFlowTok{return}\NormalTok{ total}

\NormalTok{result }\OperatorTok{=}\NormalTok{ add\_numbers(}\DecValTok{5}\NormalTok{, }\DecValTok{3}\NormalTok{)}
\BuiltInTok{print}\NormalTok{(}\SpecialStringTok{f"The sum is: }\SpecialCharTok{\{}\NormalTok{result}\SpecialCharTok{\}}\SpecialStringTok{"}\NormalTok{)   }\CommentTok{\# 8}
\end{Highlighting}
\end{Shaded}

\textbf{Key difference:} \texttt{print()} displays a value;
\texttt{return} sends a value back to the caller for further use.

\subsubsection{Skills Checklist}\label{skills-checklist-29}

\begin{itemize}
\tightlist
\item[$\square$]
  I can use \texttt{return} to send a value back from a function.
\item[$\square$]
  I can explain the difference between \texttt{print()} and
  \texttt{return}.
\item[$\square$]
  I can store the returned value of a function in a variable.
\end{itemize}

\begin{center}\rule{0.5\linewidth}{0.5pt}\end{center}

\subsection{Lesson 6.4: Functions with Lists and
Dictionaries}\label{lesson-6.4-functions-with-lists-and-dictionaries}

Since lists are \textbf{mutable}, modifying them inside a function
changes the original:

\begin{Shaded}
\begin{Highlighting}[]
\KeywordTok{def}\NormalTok{ add\_new\_user(user\_list, new\_user):}
\NormalTok{    user\_list.append(new\_user)}

\NormalTok{users }\OperatorTok{=}\NormalTok{ [}\StringTok{\textquotesingle{}Alice\textquotesingle{}}\NormalTok{, }\StringTok{\textquotesingle{}Bob\textquotesingle{}}\NormalTok{]}
\NormalTok{add\_new\_user(users, }\StringTok{\textquotesingle{}Charlie\textquotesingle{}}\NormalTok{)}
\BuiltInTok{print}\NormalTok{(users)   }\CommentTok{\# [\textquotesingle{}Alice\textquotesingle{}, \textquotesingle{}Bob\textquotesingle{}, \textquotesingle{}Charlie\textquotesingle{}]}
\end{Highlighting}
\end{Shaded}

To prevent modification, pass a copy:

\begin{Shaded}
\begin{Highlighting}[]
\NormalTok{add\_new\_user(users[:], }\StringTok{\textquotesingle{}Charlie\textquotesingle{}}\NormalTok{)   }\CommentTok{\# Original unchanged}
\end{Highlighting}
\end{Shaded}

\subsubsection{Skills Checklist}\label{skills-checklist-30}

\begin{itemize}
\tightlist
\item[$\square$]
  I can pass a list or dictionary as an argument to a function.
\item[$\square$]
  I understand that modifying a list inside a function affects the
  original.
\item[$\square$]
  I can create a copy of a list to avoid unintended modifications.
\end{itemize}

\begin{center}\rule{0.5\linewidth}{0.5pt}\end{center}

\subsection{Lesson 6.5: Arbitrary Number of
Arguments}\label{lesson-6.5-arbitrary-number-of-arguments}

\subsubsection{\texorpdfstring{\texttt{*args} --- Arbitrary Positional
Arguments}{*args --- Arbitrary Positional Arguments}}\label{args-arbitrary-positional-arguments}

\begin{Shaded}
\begin{Highlighting}[]
\KeywordTok{def}\NormalTok{ make\_pizza(size, }\OperatorTok{*}\NormalTok{toppings):}
    \BuiltInTok{print}\NormalTok{(}\SpecialStringTok{f"Making a }\SpecialCharTok{\{}\NormalTok{size}\SpecialCharTok{\}}\SpecialStringTok{{-}inch pizza with:"}\NormalTok{)}
    \ControlFlowTok{for}\NormalTok{ topping }\KeywordTok{in}\NormalTok{ toppings:}
        \BuiltInTok{print}\NormalTok{(}\SpecialStringTok{f"  {-} }\SpecialCharTok{\{}\NormalTok{topping}\SpecialCharTok{\}}\SpecialStringTok{"}\NormalTok{)}

\NormalTok{make\_pizza(}\DecValTok{12}\NormalTok{, }\StringTok{\textquotesingle{}pepperoni\textquotesingle{}}\NormalTok{)}
\NormalTok{make\_pizza(}\DecValTok{16}\NormalTok{, }\StringTok{\textquotesingle{}mushrooms\textquotesingle{}}\NormalTok{, }\StringTok{\textquotesingle{}green peppers\textquotesingle{}}\NormalTok{, }\StringTok{\textquotesingle{}extra cheese\textquotesingle{}}\NormalTok{)}
\end{Highlighting}
\end{Shaded}

\subsubsection{\texorpdfstring{\texttt{**kwargs} --- Arbitrary Keyword
Arguments}{**kwargs --- Arbitrary Keyword Arguments}}\label{kwargs-arbitrary-keyword-arguments}

\begin{Shaded}
\begin{Highlighting}[]
\KeywordTok{def}\NormalTok{ build\_profile(first, last, }\OperatorTok{**}\NormalTok{user\_info):}
\NormalTok{    user\_info[}\StringTok{\textquotesingle{}first\_name\textquotesingle{}}\NormalTok{] }\OperatorTok{=}\NormalTok{ first}
\NormalTok{    user\_info[}\StringTok{\textquotesingle{}last\_name\textquotesingle{}}\NormalTok{] }\OperatorTok{=}\NormalTok{ last}
    \ControlFlowTok{return}\NormalTok{ user\_info}

\NormalTok{profile }\OperatorTok{=}\NormalTok{ build\_profile(}\StringTok{\textquotesingle{}albert\textquotesingle{}}\NormalTok{, }\StringTok{\textquotesingle{}einstein\textquotesingle{}}\NormalTok{,}
\NormalTok{                        location}\OperatorTok{=}\StringTok{\textquotesingle{}princeton\textquotesingle{}}\NormalTok{,}
\NormalTok{                        field}\OperatorTok{=}\StringTok{\textquotesingle{}physics\textquotesingle{}}\NormalTok{)}
\BuiltInTok{print}\NormalTok{(profile)}
\end{Highlighting}
\end{Shaded}

\subsubsection{Skills Checklist}\label{skills-checklist-31}

\begin{itemize}
\tightlist
\item[$\square$]
  I can define a function using \texttt{*args} for arbitrary positional
  arguments.
\item[$\square$]
  I can define a function using \texttt{**kwargs} for arbitrary keyword
  arguments.
\end{itemize}

\newpage

\section{Unit 7: Classes and Object-Oriented
Programming}\label{unit-7-classes-and-object-oriented-programming-1}

\subsection{Overview}\label{overview-7}

This unit introduces the fundamental concepts of Object-Oriented
Programming (OOP) in Python, focusing on defining and using classes,
attributes, methods, and inheritance.

\textbf{Objectives:}

\begin{itemize}
\tightlist
\item
  Define and instantiate classes.
\item
  Work with attributes and methods.
\item
  Understand inheritance and composition.
\end{itemize}

\begin{center}\rule{0.5\linewidth}{0.5pt}\end{center}

\subsection{Lesson 7.1: Why Learn About
Classes?}\label{lesson-7.1-why-learn-about-classes}

\subsubsection{Reading References}\label{reading-references-19}

\begin{itemize}
\tightlist
\item
  Zelle: Chapter 4, pages 65--67; Chapter 10, pages 247--248
\item
  Matthes: Chapter 9, page 157
\end{itemize}

\subsubsection{Classes as Blueprints}\label{classes-as-blueprints}

A \textbf{class} is a blueprint or template for creating objects. An
\textbf{object} is a concrete instance created from that blueprint.

\textbf{Analogy:} A ``Car'' class defines that every car has a
\texttt{color} and a \texttt{speed} (attributes) and can
\texttt{accelerate()} and \texttt{brake()} (methods). Your red sedan and
your friend's blue pickup are two separate \textbf{objects} of the same
``Car'' class.

\subsubsection{Key OOP Concepts}\label{key-oop-concepts}

\begin{itemize}
\tightlist
\item
  \textbf{Encapsulation} --- Bundling data and methods into a single
  unit.
\item
  \textbf{Abstraction} --- Hiding complexity from the outside.
\item
  \textbf{Inheritance} --- New class reuses functionality of existing
  class.
\item
  \textbf{Polymorphism} --- Same method name behaves differently for
  different classes.
\end{itemize}

\subsubsection{Skills Checklist}\label{skills-checklist-32}

\begin{itemize}
\tightlist
\item[$\square$]
  I can explain the basic idea of Object-Oriented Programming.
\item[$\square$]
  I can describe a class as a blueprint and an object as an instance.
\end{itemize}

\begin{center}\rule{0.5\linewidth}{0.5pt}\end{center}

\subsection{Lesson 7.2: Creating and Using a
Class}\label{lesson-7.2-creating-and-using-a-class}

\subsubsection{Reading References}\label{reading-references-20}

\begin{itemize}
\tightlist
\item
  Zelle: Chapter 10, pages 248--254
\item
  Matthes: Chapter 9, pages 158--162
\end{itemize}

\subsubsection{Class Syntax}\label{class-syntax}

\begin{Shaded}
\begin{Highlighting}[]
\KeywordTok{class}\NormalTok{ Dog:}
    \CommentTok{"""A simple attempt to model a dog."""}

    \KeywordTok{def} \FunctionTok{\_\_init\_\_}\NormalTok{(}\VariableTok{self}\NormalTok{, name, age):}
        \CommentTok{"""Initialize name and age attributes."""}
        \VariableTok{self}\NormalTok{.name }\OperatorTok{=}\NormalTok{ name}
        \VariableTok{self}\NormalTok{.age }\OperatorTok{=}\NormalTok{ age}
        \BuiltInTok{print}\NormalTok{(}\SpecialStringTok{f"A new dog named }\SpecialCharTok{\{}\VariableTok{self}\SpecialCharTok{.}\NormalTok{name}\SpecialCharTok{\}}\SpecialStringTok{ has been created!"}\NormalTok{)}

\CommentTok{\# Create instances}
\NormalTok{my\_dog }\OperatorTok{=}\NormalTok{ Dog(}\StringTok{"Buddy"}\NormalTok{, }\DecValTok{3}\NormalTok{)}
\NormalTok{your\_dog }\OperatorTok{=}\NormalTok{ Dog(}\StringTok{"Lucy"}\NormalTok{, }\DecValTok{1}\NormalTok{)}

\BuiltInTok{print}\NormalTok{(}\SpecialStringTok{f"My dog\textquotesingle{}s name is }\SpecialCharTok{\{}\NormalTok{my\_dog}\SpecialCharTok{.}\NormalTok{name}\SpecialCharTok{\}}\SpecialStringTok{."}\NormalTok{)}
\BuiltInTok{print}\NormalTok{(}\SpecialStringTok{f"Your dog\textquotesingle{}s name is }\SpecialCharTok{\{}\NormalTok{your\_dog}\SpecialCharTok{.}\NormalTok{name}\SpecialCharTok{\}}\SpecialStringTok{."}\NormalTok{)}
\end{Highlighting}
\end{Shaded}

The \texttt{\_\_init\_\_()} method runs automatically when a new object
is created. \texttt{self} refers to the specific instance being created.

\subsubsection{Skills Checklist}\label{skills-checklist-33}

\begin{itemize}
\tightlist
\item[$\square$]
  I can define a basic Python class.
\item[$\square$]
  I can define the \texttt{\_\_init\_\_} method within a class.
\item[$\square$]
  I understand the purpose of the \texttt{self} parameter.
\item[$\square$]
  I can create an instance (object) of a class.
\end{itemize}

\begin{center}\rule{0.5\linewidth}{0.5pt}\end{center}

\subsection{Lesson 7.3: Working with Attributes and
Methods}\label{lesson-7.3-working-with-attributes-and-methods}

\subsubsection{Reading References}\label{reading-references-21}

\begin{itemize}
\tightlist
\item
  Matthes: Chapter 9, pages 162--166
\end{itemize}

\subsubsection{Calling Methods}\label{calling-methods}

\begin{Shaded}
\begin{Highlighting}[]
\KeywordTok{class}\NormalTok{ Dog:}
    \KeywordTok{def} \FunctionTok{\_\_init\_\_}\NormalTok{(}\VariableTok{self}\NormalTok{, name, age):}
        \VariableTok{self}\NormalTok{.name }\OperatorTok{=}\NormalTok{ name}
        \VariableTok{self}\NormalTok{.age }\OperatorTok{=}\NormalTok{ age}

    \KeywordTok{def}\NormalTok{ sit(}\VariableTok{self}\NormalTok{):}
        \CommentTok{"""Simulate a dog sitting."""}
        \BuiltInTok{print}\NormalTok{(}\SpecialStringTok{f"}\SpecialCharTok{\{}\VariableTok{self}\SpecialCharTok{.}\NormalTok{name}\SpecialCharTok{\}}\SpecialStringTok{ is now sitting."}\NormalTok{)}

    \KeywordTok{def}\NormalTok{ roll\_over(}\VariableTok{self}\NormalTok{):}
        \CommentTok{"""Simulate a dog rolling over."""}
        \BuiltInTok{print}\NormalTok{(}\SpecialStringTok{f"}\SpecialCharTok{\{}\VariableTok{self}\SpecialCharTok{.}\NormalTok{name}\SpecialCharTok{\}}\SpecialStringTok{ rolled over!"}\NormalTok{)}

    \KeywordTok{def}\NormalTok{ get\_older(}\VariableTok{self}\NormalTok{, years\_to\_add):}
        \CommentTok{"""Add years to the dog\textquotesingle{}s age."""}
        \VariableTok{self}\NormalTok{.age }\OperatorTok{+=}\NormalTok{ years\_to\_add}
        \BuiltInTok{print}\NormalTok{(}\SpecialStringTok{f"}\SpecialCharTok{\{}\VariableTok{self}\SpecialCharTok{.}\NormalTok{name}\SpecialCharTok{\}}\SpecialStringTok{ is now }\SpecialCharTok{\{}\VariableTok{self}\SpecialCharTok{.}\NormalTok{age}\SpecialCharTok{\}}\SpecialStringTok{ years old."}\NormalTok{)}

\NormalTok{my\_dog }\OperatorTok{=}\NormalTok{ Dog(}\StringTok{"Buddy"}\NormalTok{, }\DecValTok{3}\NormalTok{)}
\NormalTok{my\_dog.sit()}
\NormalTok{my\_dog.roll\_over()}
\NormalTok{my\_dog.age }\OperatorTok{=} \DecValTok{4}          \CommentTok{\# Direct attribute modification}
\NormalTok{my\_dog.get\_older(}\DecValTok{2}\NormalTok{)     }\CommentTok{\# Method{-}based modification}
\end{Highlighting}
\end{Shaded}

\subsubsection{Skills Checklist}\label{skills-checklist-34}

\begin{itemize}
\tightlist
\item[$\square$]
  I can access an object's attributes using dot notation.
\item[$\square$]
  I can define and call methods within a class.
\item[$\square$]
  I can use a method to change an object's attribute.
\end{itemize}

\begin{center}\rule{0.5\linewidth}{0.5pt}\end{center}

\subsection{Lesson 7.4: Inheritance}\label{lesson-7.4-inheritance}

\subsubsection{Reading References}\label{reading-references-22}

\begin{itemize}
\tightlist
\item
  Zelle: Chapter 12, pages 351--352
\item
  Matthes: Chapter 9, pages 167--170
\end{itemize}

\subsubsection{Creating a Child Class}\label{creating-a-child-class}

\begin{Shaded}
\begin{Highlighting}[]
\KeywordTok{class}\NormalTok{ Car:}
    \KeywordTok{def} \FunctionTok{\_\_init\_\_}\NormalTok{(}\VariableTok{self}\NormalTok{, make, model, year):}
        \VariableTok{self}\NormalTok{.make }\OperatorTok{=}\NormalTok{ make}
        \VariableTok{self}\NormalTok{.model }\OperatorTok{=}\NormalTok{ model}
        \VariableTok{self}\NormalTok{.year }\OperatorTok{=}\NormalTok{ year}
        \VariableTok{self}\NormalTok{.odometer\_reading }\OperatorTok{=} \DecValTok{0}

    \KeywordTok{def}\NormalTok{ get\_descriptive\_name(}\VariableTok{self}\NormalTok{):}
        \ControlFlowTok{return} \SpecialStringTok{f"}\SpecialCharTok{\{}\VariableTok{self}\SpecialCharTok{.}\NormalTok{year}\SpecialCharTok{\}}\SpecialStringTok{ }\SpecialCharTok{\{}\VariableTok{self}\SpecialCharTok{.}\NormalTok{make}\SpecialCharTok{\}}\SpecialStringTok{ }\SpecialCharTok{\{}\VariableTok{self}\SpecialCharTok{.}\NormalTok{model}\SpecialCharTok{\}}\SpecialStringTok{"}\NormalTok{.title()}

    \KeywordTok{def}\NormalTok{ read\_odometer(}\VariableTok{self}\NormalTok{):}
        \BuiltInTok{print}\NormalTok{(}\SpecialStringTok{f"This car has }\SpecialCharTok{\{}\VariableTok{self}\SpecialCharTok{.}\NormalTok{odometer\_reading}\SpecialCharTok{\}}\SpecialStringTok{ miles on it."}\NormalTok{)}


\KeywordTok{class}\NormalTok{ ElectricCar(Car):}
    \KeywordTok{def} \FunctionTok{\_\_init\_\_}\NormalTok{(}\VariableTok{self}\NormalTok{, make, model, year):}
        \BuiltInTok{super}\NormalTok{().}\FunctionTok{\_\_init\_\_}\NormalTok{(make, model, year)   }\CommentTok{\# Call parent\textquotesingle{}s \_\_init\_\_}
        \VariableTok{self}\NormalTok{.battery\_size }\OperatorTok{=} \DecValTok{75}

    \KeywordTok{def}\NormalTok{ describe\_battery(}\VariableTok{self}\NormalTok{):}
        \BuiltInTok{print}\NormalTok{(}\SpecialStringTok{f"This car has a }\SpecialCharTok{\{}\VariableTok{self}\SpecialCharTok{.}\NormalTok{battery\_size}\SpecialCharTok{\}}\SpecialStringTok{{-}kWh battery."}\NormalTok{)}


\NormalTok{my\_ev }\OperatorTok{=}\NormalTok{ ElectricCar(}\StringTok{\textquotesingle{}ford\textquotesingle{}}\NormalTok{, }\StringTok{\textquotesingle{}mustang mach e\textquotesingle{}}\NormalTok{, }\DecValTok{2023}\NormalTok{)}
\BuiltInTok{print}\NormalTok{(my\_ev.get\_descriptive\_name())   }\CommentTok{\# Inherited method}
\NormalTok{my\_ev.describe\_battery()              }\CommentTok{\# New method}
\NormalTok{my\_ev.read\_odometer()                 }\CommentTok{\# Inherited method}
\end{Highlighting}
\end{Shaded}

\subsubsection{\texorpdfstring{\texttt{super()}
Function}{super() Function}}\label{super-function}

\texttt{super().\_\_init\_\_()} calls the parent class's
\texttt{\_\_init\_\_()} method, ensuring inherited attributes are
properly initialized.

\subsubsection{Skills Checklist}\label{skills-checklist-35}

\begin{itemize}
\tightlist
\item[$\square$]
  I can define a child class that inherits from a parent class.
\item[$\square$]
  I can use \texttt{super().\_\_init\_\_()} to properly initialize
  inherited attributes.
\item[$\square$]
  I can add new attributes and methods to a child class.
\item[$\square$]
  I can override a parent method in a child class.
\end{itemize}

\begin{center}\rule{0.5\linewidth}{0.5pt}\end{center}

\subsection{Lesson 7.5: Instances as Attributes
(Composition)}\label{lesson-7.5-instances-as-attributes-composition}

\subsubsection{Composition --- ``Has-A''
Relationship}\label{composition-has-a-relationship}

Instead of inheriting, a class \textbf{contains an instance of another
class} as an attribute.

\begin{Shaded}
\begin{Highlighting}[]
\KeywordTok{class}\NormalTok{ Battery:}
    \KeywordTok{def} \FunctionTok{\_\_init\_\_}\NormalTok{(}\VariableTok{self}\NormalTok{, battery\_size}\OperatorTok{=}\DecValTok{75}\NormalTok{):}
        \VariableTok{self}\NormalTok{.battery\_size }\OperatorTok{=}\NormalTok{ battery\_size}

    \KeywordTok{def}\NormalTok{ describe\_battery(}\VariableTok{self}\NormalTok{):}
        \BuiltInTok{print}\NormalTok{(}\SpecialStringTok{f"This car has a }\SpecialCharTok{\{}\VariableTok{self}\SpecialCharTok{.}\NormalTok{battery\_size}\SpecialCharTok{\}}\SpecialStringTok{{-}kWh battery."}\NormalTok{)}

    \KeywordTok{def}\NormalTok{ get\_range(}\VariableTok{self}\NormalTok{):}
        \ControlFlowTok{if} \VariableTok{self}\NormalTok{.battery\_size }\OperatorTok{==} \DecValTok{75}\NormalTok{:}
\NormalTok{            range\_miles }\OperatorTok{=} \DecValTok{260}
        \ControlFlowTok{else}\NormalTok{:}
\NormalTok{            range\_miles }\OperatorTok{=} \DecValTok{310}
        \BuiltInTok{print}\NormalTok{(}\SpecialStringTok{f"About }\SpecialCharTok{\{}\NormalTok{range\_miles}\SpecialCharTok{\}}\SpecialStringTok{ miles of range."}\NormalTok{)}


\KeywordTok{class}\NormalTok{ ElectricCar:}
    \KeywordTok{def} \FunctionTok{\_\_init\_\_}\NormalTok{(}\VariableTok{self}\NormalTok{, make, model, year):}
        \VariableTok{self}\NormalTok{.make }\OperatorTok{=}\NormalTok{ make}
        \VariableTok{self}\NormalTok{.model }\OperatorTok{=}\NormalTok{ model}
        \VariableTok{self}\NormalTok{.year }\OperatorTok{=}\NormalTok{ year}
        \VariableTok{self}\NormalTok{.battery }\OperatorTok{=}\NormalTok{ Battery()   }\CommentTok{\# Composition: ElectricCar HAS A Battery}


\NormalTok{my\_ev }\OperatorTok{=}\NormalTok{ ElectricCar(}\StringTok{\textquotesingle{}ford\textquotesingle{}}\NormalTok{, }\StringTok{\textquotesingle{}mustang mach e\textquotesingle{}}\NormalTok{, }\DecValTok{2023}\NormalTok{)}
\NormalTok{my\_ev.battery.describe\_battery()}
\NormalTok{my\_ev.battery.get\_range()}
\end{Highlighting}
\end{Shaded}

\subsubsection{Skills Checklist}\label{skills-checklist-36}

\begin{itemize}
\tightlist
\item[$\square$]
  I can explain the concept of composition in OOP.
\item[$\square$]
  I can create an instance of one class and assign it as an attribute of
  another.
\item[$\square$]
  I can access attributes and methods of a nested object.
\end{itemize}

\begin{center}\rule{0.5\linewidth}{0.5pt}\end{center}

\section{Unit 8: Files and
Exceptions}\label{unit-8-files-and-exceptions-1}

\subsection{Overview}\label{overview-8}

This unit covers interacting with text files for reading and writing
data, handling common errors gracefully using exceptions, and storing
data using JSON.

\textbf{Objectives:}

\begin{itemize}
\tightlist
\item
  Read from and write to files.
\item
  Handle exceptions with \texttt{try-except} blocks.
\item
  Store and load data using JSON.
\end{itemize}

\begin{center}\rule{0.5\linewidth}{0.5pt}\end{center}

\subsection{Lesson 8.1: Reading from a
File}\label{lesson-8.1-reading-from-a-file}

\subsubsection{Reading References}\label{reading-references-23}

\begin{itemize}
\tightlist
\item
  Zelle: Chapter 5, pages 125--127
\item
  Matthes: Chapter 10, pages 184--189
\end{itemize}

\subsubsection{Reading Files}\label{reading-files}

\begin{Shaded}
\begin{Highlighting}[]
\CommentTok{\# Reading entire file}
\ControlFlowTok{with} \BuiltInTok{open}\NormalTok{(}\StringTok{\textquotesingle{}my\_notes.txt\textquotesingle{}}\NormalTok{, }\StringTok{\textquotesingle{}r\textquotesingle{}}\NormalTok{) }\ImportTok{as}\NormalTok{ file\_object:}
\NormalTok{    contents }\OperatorTok{=}\NormalTok{ file\_object.read()}
\BuiltInTok{print}\NormalTok{(contents)}

\CommentTok{\# Reading line by line}
\ControlFlowTok{with} \BuiltInTok{open}\NormalTok{(}\StringTok{\textquotesingle{}my\_notes.txt\textquotesingle{}}\NormalTok{, }\StringTok{\textquotesingle{}r\textquotesingle{}}\NormalTok{) }\ImportTok{as}\NormalTok{ file\_object:}
    \ControlFlowTok{for}\NormalTok{ line }\KeywordTok{in}\NormalTok{ file\_object:}
        \BuiltInTok{print}\NormalTok{(line.strip())}

\CommentTok{\# Reading all lines into a list}
\ControlFlowTok{with} \BuiltInTok{open}\NormalTok{(}\StringTok{\textquotesingle{}my\_notes.txt\textquotesingle{}}\NormalTok{, }\StringTok{\textquotesingle{}r\textquotesingle{}}\NormalTok{) }\ImportTok{as}\NormalTok{ file\_object:}
\NormalTok{    lines }\OperatorTok{=}\NormalTok{ file\_object.readlines()}
\end{Highlighting}
\end{Shaded}

The \texttt{with} statement \textbf{automatically closes the file} even
if an error occurs.

\subsubsection{File Modes}\label{file-modes}

{\def\LTcaptype{none} % do not increment counter
\begin{longtable}[]{@{}ll@{}}
\toprule\noalign{}
Mode & Description \\
\midrule\noalign{}
\endhead
\bottomrule\noalign{}
\endlastfoot
\texttt{\textquotesingle{}r\textquotesingle{}} & Read (default) \\
\texttt{\textquotesingle{}w\textquotesingle{}} & Write (creates new or
overwrites) \\
\texttt{\textquotesingle{}a\textquotesingle{}} & Append (adds to end) \\
\end{longtable}
}

\subsubsection{Skills Checklist}\label{skills-checklist-37}

\begin{itemize}
\tightlist
\item[$\square$]
  I can open a file for reading using \texttt{with\ open()}.
\item[$\square$]
  I can read the entire content of a file into a string.
\item[$\square$]
  I can read a file line by line using a \texttt{for} loop.
\item[$\square$]
  I can remove whitespace/newlines using \texttt{.strip()}.
\end{itemize}

\begin{center}\rule{0.5\linewidth}{0.5pt}\end{center}

\subsection{Lesson 8.2: Writing to a
File}\label{lesson-8.2-writing-to-a-file}

\subsubsection{Reading References}\label{reading-references-24}

\begin{itemize}
\tightlist
\item
  Zelle: Chapter 5, pages 127--128
\item
  Matthes: Chapter 10, pages 190--192
\end{itemize}

\begin{Shaded}
\begin{Highlighting}[]
\CommentTok{\# Write mode (overwrites existing content)}
\ControlFlowTok{with} \BuiltInTok{open}\NormalTok{(}\StringTok{\textquotesingle{}programming\_message.txt\textquotesingle{}}\NormalTok{, }\StringTok{\textquotesingle{}w\textquotesingle{}}\NormalTok{) }\ImportTok{as}\NormalTok{ f:}
\NormalTok{    f.write(}\StringTok{"I love programming in Python!}\CharTok{\textbackslash{}n}\StringTok{"}\NormalTok{)}

\CommentTok{\# Append mode (adds to end)}
\ControlFlowTok{with} \BuiltInTok{open}\NormalTok{(}\StringTok{\textquotesingle{}programming\_message.txt\textquotesingle{}}\NormalTok{, }\StringTok{\textquotesingle{}a\textquotesingle{}}\NormalTok{) }\ImportTok{as}\NormalTok{ f:}
\NormalTok{    f.write(}\StringTok{"It\textquotesingle{}s a very versatile language.}\CharTok{\textbackslash{}n}\StringTok{"}\NormalTok{)}
\NormalTok{    f.write(}\StringTok{"I\textquotesingle{}m excited to learn more!}\CharTok{\textbackslash{}n}\StringTok{"}\NormalTok{)}
\end{Highlighting}
\end{Shaded}

\textbf{Important:} The \texttt{.write()} method does not add newlines
automatically. Use \texttt{\textbackslash{}n}.

\subsubsection{Skills Checklist}\label{skills-checklist-38}

\begin{itemize}
\tightlist
\item[$\square$]
  I can open a file for writing using \texttt{"w"} mode.
\item[$\square$]
  I can open a file for appending using \texttt{"a"} mode.
\item[$\square$]
  I can write a string to a file.
\item[$\square$]
  I can add newlines to a file using \texttt{\textbackslash{}n}.
\end{itemize}

\begin{center}\rule{0.5\linewidth}{0.5pt}\end{center}

\subsection{Lesson 8.3: Exceptions --- Handling
Errors}\label{lesson-8.3-exceptions-handling-errors}

\subsubsection{Reading References}\label{reading-references-25}

\begin{itemize}
\tightlist
\item
  Zelle: Chapter 7, pages 177--179
\item
  Matthes: Chapter 10, pages 192--199
\end{itemize}

\subsubsection{\texorpdfstring{\texttt{try-except}
Block}{try-except Block}}\label{try-except-block}

\begin{Shaded}
\begin{Highlighting}[]
\ControlFlowTok{try}\NormalTok{:}
\NormalTok{    result }\OperatorTok{=} \DecValTok{10} \OperatorTok{/} \DecValTok{0}
\ControlFlowTok{except} \PreprocessorTok{ZeroDivisionError}\NormalTok{:}
    \BuiltInTok{print}\NormalTok{(}\StringTok{"Error: You cannot divide by zero!"}\NormalTok{)}
\BuiltInTok{print}\NormalTok{(}\StringTok{"Program continues after error handling."}\NormalTok{)}
\end{Highlighting}
\end{Shaded}

\subsubsection{Common Exception Types}\label{common-exception-types}

{\def\LTcaptype{none} % do not increment counter
\begin{longtable}[]{@{}ll@{}}
\toprule\noalign{}
Exception & When it occurs \\
\midrule\noalign{}
\endhead
\bottomrule\noalign{}
\endlastfoot
\texttt{ValueError} & Invalid value (e.g., \texttt{int("abc")}) \\
\texttt{ZeroDivisionError} & Division by zero \\
\texttt{FileNotFoundError} & File doesn't exist \\
\texttt{IndexError} & List index out of range \\
\texttt{KeyError} & Dictionary key doesn't exist \\
\texttt{TypeError} & Wrong type (e.g., \texttt{"5"\ +\ 5}) \\
\end{longtable}
}

\subsubsection{Examples}\label{examples-13}

\textbf{Example 8.3.1: Handling ZeroDivisionError}

\begin{Shaded}
\begin{Highlighting}[]
\ControlFlowTok{try}\NormalTok{:}
\NormalTok{    result }\OperatorTok{=} \DecValTok{10} \OperatorTok{/} \DecValTok{0}
\ControlFlowTok{except} \PreprocessorTok{ZeroDivisionError}\NormalTok{:}
    \BuiltInTok{print}\NormalTok{(}\StringTok{"Error: Cannot divide by zero!"}\NormalTok{)}
\end{Highlighting}
\end{Shaded}

\textbf{Example 8.3.2: Handling ValueError}

\begin{Shaded}
\begin{Highlighting}[]
\ControlFlowTok{try}\NormalTok{:}
\NormalTok{    num }\OperatorTok{=} \BuiltInTok{int}\NormalTok{(}\BuiltInTok{input}\NormalTok{(}\StringTok{"Enter a number: "}\NormalTok{))}
    \BuiltInTok{print}\NormalTok{(}\SpecialStringTok{f"You entered: }\SpecialCharTok{\{}\NormalTok{num}\SpecialCharTok{\}}\SpecialStringTok{"}\NormalTok{)}
\ControlFlowTok{except} \PreprocessorTok{ValueError}\NormalTok{:}
    \BuiltInTok{print}\NormalTok{(}\StringTok{"Error: That was not a valid number."}\NormalTok{)}
\end{Highlighting}
\end{Shaded}

\subsubsection{Skills Checklist}\label{skills-checklist-39}

\begin{itemize}
\tightlist
\item[$\square$]
  I can use a \texttt{try-except} block to handle errors.
\item[$\square$]
  I can name common exception types and what causes them.
\item[$\square$]
  I can catch specific exception types.
\end{itemize}

\begin{center}\rule{0.5\linewidth}{0.5pt}\end{center}

\subsection{\texorpdfstring{Lesson 8.4: The \texttt{else} and
\texttt{finally}
Blocks}{Lesson 8.4: The else and finally Blocks}}\label{lesson-8.4-the-else-and-finally-blocks}

\subsubsection{\texorpdfstring{\texttt{else} Block --- Runs only on
success}{else Block --- Runs only on success}}\label{else-block-runs-only-on-success}

\begin{Shaded}
\begin{Highlighting}[]
\ControlFlowTok{try}\NormalTok{:}
\NormalTok{    num }\OperatorTok{=} \BuiltInTok{int}\NormalTok{(}\BuiltInTok{input}\NormalTok{(}\StringTok{"Enter a number: "}\NormalTok{))}
\ControlFlowTok{except} \PreprocessorTok{ValueError}\NormalTok{:}
    \BuiltInTok{print}\NormalTok{(}\StringTok{"That\textquotesingle{}s not a number."}\NormalTok{)}
\ControlFlowTok{else}\NormalTok{:}
    \BuiltInTok{print}\NormalTok{(}\SpecialStringTok{f"You entered }\SpecialCharTok{\{}\NormalTok{num}\SpecialCharTok{\}}\SpecialStringTok{."}\NormalTok{)   }\CommentTok{\# Only runs if no exception}
\end{Highlighting}
\end{Shaded}

\subsubsection{\texorpdfstring{\texttt{finally} Block --- Always runs
(cleanup)}{finally Block --- Always runs (cleanup)}}\label{finally-block-always-runs-cleanup}

\begin{Shaded}
\begin{Highlighting}[]
\NormalTok{file\_handle }\OperatorTok{=} \VariableTok{None}
\ControlFlowTok{try}\NormalTok{:}
\NormalTok{    file\_handle }\OperatorTok{=} \BuiltInTok{open}\NormalTok{(}\StringTok{\textquotesingle{}data.txt\textquotesingle{}}\NormalTok{, }\StringTok{\textquotesingle{}w\textquotesingle{}}\NormalTok{)}
\NormalTok{    file\_handle.write(}\StringTok{"Some data."}\NormalTok{)}
\ControlFlowTok{except} \PreprocessorTok{TypeError}\NormalTok{:}
    \BuiltInTok{print}\NormalTok{(}\StringTok{"A type error occurred."}\NormalTok{)}
\ControlFlowTok{finally}\NormalTok{:}
    \ControlFlowTok{if}\NormalTok{ file\_handle:}
\NormalTok{        file\_handle.close()     }\CommentTok{\# Always closes the file}
        \BuiltInTok{print}\NormalTok{(}\StringTok{"File closed."}\NormalTok{)}
\end{Highlighting}
\end{Shaded}

\subsubsection{Skills Checklist}\label{skills-checklist-40}

\begin{itemize}
\tightlist
\item[$\square$]
  I can use an \texttt{else} block for code that runs only on success.
\item[$\square$]
  I can use a \texttt{finally} block for code that always runs.
\end{itemize}

\begin{center}\rule{0.5\linewidth}{0.5pt}\end{center}

\subsection{Lesson 8.5: Storing Data with
JSON}\label{lesson-8.5-storing-data-with-json}

\subsubsection{Reading References}\label{reading-references-26}

\begin{itemize}
\tightlist
\item
  Matthes: Chapter 10, pages 200--204
\end{itemize}

\subsubsection{\texorpdfstring{The \texttt{json}
Module}{The json Module}}\label{the-json-module}

JSON (JavaScript Object Notation) is a lightweight format for storing
and exchanging data. Python's \texttt{json} module converts between
Python objects and JSON.

\begin{Shaded}
\begin{Highlighting}[]
\ImportTok{import}\NormalTok{ json}

\CommentTok{\# Save data}
\NormalTok{user\_settings }\OperatorTok{=}\NormalTok{ \{}
    \StringTok{"username"}\NormalTok{: }\StringTok{"coder\_kid"}\NormalTok{,}
    \StringTok{"theme"}\NormalTok{: }\StringTok{"dark"}\NormalTok{,}
    \StringTok{"notifications\_on"}\NormalTok{: }\VariableTok{True}\NormalTok{,}
    \StringTok{"favorite\_numbers"}\NormalTok{: [}\DecValTok{7}\NormalTok{, }\DecValTok{13}\NormalTok{, }\DecValTok{42}\NormalTok{]}
\NormalTok{\}}

\ControlFlowTok{with} \BuiltInTok{open}\NormalTok{(}\StringTok{\textquotesingle{}user\_settings.json\textquotesingle{}}\NormalTok{, }\StringTok{\textquotesingle{}w\textquotesingle{}}\NormalTok{) }\ImportTok{as}\NormalTok{ f:}
\NormalTok{    json.dump(user\_settings, f, indent}\OperatorTok{=}\DecValTok{4}\NormalTok{)}

\CommentTok{\# Load data}
\ControlFlowTok{try}\NormalTok{:}
    \ControlFlowTok{with} \BuiltInTok{open}\NormalTok{(}\StringTok{\textquotesingle{}user\_settings.json\textquotesingle{}}\NormalTok{, }\StringTok{\textquotesingle{}r\textquotesingle{}}\NormalTok{) }\ImportTok{as}\NormalTok{ f:}
\NormalTok{        loaded\_settings }\OperatorTok{=}\NormalTok{ json.load(f)}
    \BuiltInTok{print}\NormalTok{(}\SpecialStringTok{f"Username: }\SpecialCharTok{\{}\NormalTok{loaded\_settings[}\StringTok{\textquotesingle{}username\textquotesingle{}}\NormalTok{]}\SpecialCharTok{\}}\SpecialStringTok{"}\NormalTok{)}
    \BuiltInTok{print}\NormalTok{(}\SpecialStringTok{f"Theme: }\SpecialCharTok{\{}\NormalTok{loaded\_settings[}\StringTok{\textquotesingle{}theme\textquotesingle{}}\NormalTok{]}\SpecialCharTok{\}}\SpecialStringTok{"}\NormalTok{)}
\ControlFlowTok{except} \PreprocessorTok{FileNotFoundError}\NormalTok{:}
    \BuiltInTok{print}\NormalTok{(}\StringTok{"Settings file not found."}\NormalTok{)}
\ControlFlowTok{except}\NormalTok{ json.JSONDecodeError:}
    \BuiltInTok{print}\NormalTok{(}\StringTok{"Error: Could not decode settings file."}\NormalTok{)}
\end{Highlighting}
\end{Shaded}

\subsubsection{JSON\textless-\textgreater Python
Mapping}\label{json-python-mapping}

{\def\LTcaptype{none} % do not increment counter
\begin{longtable}[]{@{}ll@{}}
\toprule\noalign{}
JSON & Python \\
\midrule\noalign{}
\endhead
\bottomrule\noalign{}
\endlastfoot
Object \texttt{\{\}} & Dictionary \texttt{\{\}} \\
Array \texttt{{[}{]}} & List \texttt{{[}{]}} \\
String & \texttt{str} \\
Number & \texttt{int} or \texttt{float} \\
Boolean & \texttt{bool} \\
null & \texttt{None} \\
\end{longtable}
}

\subsubsection{Skills Checklist}\label{skills-checklist-41}

\begin{itemize}
\tightlist
\item[$\square$]
  I can explain what JSON is and why it's used.
\item[$\square$]
  I can use \texttt{json.dump()} to save Python data to a JSON file.
\item[$\square$]
  I can use \texttt{json.load()} to load JSON data from a file.
\item[$\square$]
  I can save and load lists of dictionaries using JSON.
\end{itemize}

\begin{center}\rule{0.5\linewidth}{0.5pt}\end{center}

\section{Unit 9: Testing and Data
Analysis}\label{unit-9-testing-and-data-analysis-1}

\subsection{Overview}\label{overview-9}

This unit introduces testing code for correctness and explores basic
text-based data analysis.

\begin{center}\rule{0.5\linewidth}{0.5pt}\end{center}

\subsection{Lesson 9.1: Testing a
Function}\label{lesson-9.1-testing-a-function}

\subsubsection{Reading References}\label{reading-references-27}

\begin{itemize}
\tightlist
\item
  Matthes: Chapter 11, pages 209--217
\end{itemize}

\subsubsection{\texorpdfstring{Unit Testing with
\texttt{unittest}}{Unit Testing with unittest}}\label{unit-testing-with-unittest}

\begin{Shaded}
\begin{Highlighting}[]
\CommentTok{\# my\_functions.py}
\KeywordTok{def}\NormalTok{ add\_two\_numbers(a, b):}
    \ControlFlowTok{return}\NormalTok{ a }\OperatorTok{+}\NormalTok{ b}
\end{Highlighting}
\end{Shaded}

\begin{Shaded}
\begin{Highlighting}[]
\CommentTok{\# test\_my\_functions.py}
\ImportTok{import}\NormalTok{ unittest}
\ImportTok{from}\NormalTok{ my\_functions }\ImportTok{import}\NormalTok{ add\_two\_numbers}

\KeywordTok{class}\NormalTok{ TestAddNumbers(unittest.TestCase):}
    \KeywordTok{def}\NormalTok{ test\_positive\_numbers(}\VariableTok{self}\NormalTok{):}
\NormalTok{        result }\OperatorTok{=}\NormalTok{ add\_two\_numbers(}\DecValTok{5}\NormalTok{, }\DecValTok{3}\NormalTok{)}
        \VariableTok{self}\NormalTok{.assertEqual(result, }\DecValTok{8}\NormalTok{)}

    \KeywordTok{def}\NormalTok{ test\_negative\_numbers(}\VariableTok{self}\NormalTok{):}
\NormalTok{        result }\OperatorTok{=}\NormalTok{ add\_two\_numbers(}\OperatorTok{{-}}\DecValTok{5}\NormalTok{, }\OperatorTok{{-}}\DecValTok{3}\NormalTok{)}
        \VariableTok{self}\NormalTok{.assertEqual(result, }\OperatorTok{{-}}\DecValTok{8}\NormalTok{)}

    \KeywordTok{def}\NormalTok{ test\_mixed\_numbers(}\VariableTok{self}\NormalTok{):}
\NormalTok{        result }\OperatorTok{=}\NormalTok{ add\_two\_numbers(}\DecValTok{10}\NormalTok{, }\OperatorTok{{-}}\DecValTok{4}\NormalTok{)}
        \VariableTok{self}\NormalTok{.assertEqual(result, }\DecValTok{6}\NormalTok{)}

\ControlFlowTok{if} \VariableTok{\_\_name\_\_} \OperatorTok{==} \StringTok{\textquotesingle{}\_\_main\_\_\textquotesingle{}}\NormalTok{:}
\NormalTok{    unittest.main()}
\end{Highlighting}
\end{Shaded}

\subsubsection{Common Assertion Methods}\label{common-assertion-methods}

{\def\LTcaptype{none} % do not increment counter
\begin{longtable}[]{@{}ll@{}}
\toprule\noalign{}
Method & Description \\
\midrule\noalign{}
\endhead
\bottomrule\noalign{}
\endlastfoot
\texttt{assertEqual(a,\ b)} & Check a == b \\
\texttt{assertTrue(x)} & Check x is True \\
\texttt{assertFalse(x)} & Check x is False \\
\texttt{assertIn(a,\ b)} & Check a in b \\
\texttt{assertIsNone(x)} & Check x is None \\
\end{longtable}
}

\subsubsection{Skills Checklist}\label{skills-checklist-42}

\begin{itemize}
\tightlist
\item[$\square$]
  I can create a basic test class using \texttt{unittest.TestCase}.
\item[$\square$]
  I can define test methods within a test class.
\item[$\square$]
  I can use \texttt{assertEqual()} to check for expected results.
\item[$\square$]
  I can run unit tests from the terminal.
\end{itemize}

\begin{center}\rule{0.5\linewidth}{0.5pt}\end{center}

\subsection{Lesson 9.2: Testing a
Class}\label{lesson-9.2-testing-a-class}

\subsubsection{\texorpdfstring{The \texttt{setUp()}
Method}{The setUp() Method}}\label{the-setup-method}

\texttt{setUp()} runs \textbf{before every test method}, providing a
fresh instance each time:

\begin{Shaded}
\begin{Highlighting}[]
\KeywordTok{class}\NormalTok{ TestEmployee(unittest.TestCase):}
    \KeywordTok{def}\NormalTok{ setUp(}\VariableTok{self}\NormalTok{):}
        \CommentTok{"""Create an employee instance for all test methods."""}
        \VariableTok{self}\NormalTok{.employee }\OperatorTok{=}\NormalTok{ Employee(}\StringTok{\textquotesingle{}janis\textquotesingle{}}\NormalTok{, }\StringTok{\textquotesingle{}joplin\textquotesingle{}}\NormalTok{, }\DecValTok{60000}\NormalTok{)}

    \KeywordTok{def}\NormalTok{ test\_give\_default\_raise(}\VariableTok{self}\NormalTok{):}
        \VariableTok{self}\NormalTok{.employee.give\_raise()}
        \VariableTok{self}\NormalTok{.assertEqual(}\VariableTok{self}\NormalTok{.employee.salary, }\DecValTok{65000}\NormalTok{)}

    \KeywordTok{def}\NormalTok{ test\_give\_custom\_raise(}\VariableTok{self}\NormalTok{):}
        \VariableTok{self}\NormalTok{.employee.give\_raise(}\DecValTok{10000}\NormalTok{)}
        \VariableTok{self}\NormalTok{.assertEqual(}\VariableTok{self}\NormalTok{.employee.salary, }\DecValTok{70000}\NormalTok{)}
\end{Highlighting}
\end{Shaded}

\subsubsection{Skills Checklist}\label{skills-checklist-43}

\begin{itemize}
\tightlist
\item[$\square$]
  I can write tests for a Python class.
\item[$\square$]
  I can use \texttt{setUp()} to prepare my tests.
\end{itemize}

\begin{center}\rule{0.5\linewidth}{0.5pt}\end{center}

\subsection{Lesson 9.3: Introduction to Random
Numbers}\label{lesson-9.3-introduction-to-random-numbers}

\begin{Shaded}
\begin{Highlighting}[]
\ImportTok{import}\NormalTok{ random}

\CommentTok{\# Random integer (inclusive)}
\NormalTok{die\_roll }\OperatorTok{=}\NormalTok{ random.randint(}\DecValTok{1}\NormalTok{, }\DecValTok{6}\NormalTok{)}
\BuiltInTok{print}\NormalTok{(}\SpecialStringTok{f"You rolled: }\SpecialCharTok{\{}\NormalTok{die\_roll}\SpecialCharTok{\}}\SpecialStringTok{"}\NormalTok{)}

\CommentTok{\# Random float between 0.0 and 1.0}
\NormalTok{rand\_float }\OperatorTok{=}\NormalTok{ random.random()}

\CommentTok{\# Random choice from list}
\NormalTok{colors }\OperatorTok{=}\NormalTok{ [}\StringTok{"red"}\NormalTok{, }\StringTok{"blue"}\NormalTok{, }\StringTok{"green"}\NormalTok{]}
\NormalTok{chosen }\OperatorTok{=}\NormalTok{ random.choice(colors)}
\BuiltInTok{print}\NormalTok{(}\SpecialStringTok{f"Chosen color: }\SpecialCharTok{\{}\NormalTok{chosen}\SpecialCharTok{\}}\SpecialStringTok{"}\NormalTok{)}
\end{Highlighting}
\end{Shaded}

\subsubsection{Skills Checklist}\label{skills-checklist-44}

\begin{itemize}
\tightlist
\item[$\square$]
  I can import the \texttt{random} module.
\item[$\square$]
  I can generate a random integer within a range.
\item[$\square$]
  I can select a random item from a list.
\end{itemize}

\begin{center}\rule{0.5\linewidth}{0.5pt}\end{center}

\subsection{Lesson 9.4: Analyzing Text
Data}\label{lesson-9.4-analyzing-text-data}

\begin{Shaded}
\begin{Highlighting}[]
\NormalTok{sentence }\OperatorTok{=} \StringTok{"the quick brown fox jumps over the lazy dog the quick brown fox"}
\NormalTok{words }\OperatorTok{=}\NormalTok{ sentence.lower().split()}

\NormalTok{word\_frequencies }\OperatorTok{=}\NormalTok{ \{\}}
\ControlFlowTok{for}\NormalTok{ word }\KeywordTok{in}\NormalTok{ words:}
    \ControlFlowTok{if}\NormalTok{ word }\KeywordTok{in}\NormalTok{ word\_frequencies:}
\NormalTok{        word\_frequencies[word] }\OperatorTok{+=} \DecValTok{1}
    \ControlFlowTok{else}\NormalTok{:}
\NormalTok{        word\_frequencies[word] }\OperatorTok{=} \DecValTok{1}

\BuiltInTok{print}\NormalTok{(}\StringTok{"Word Frequencies:"}\NormalTok{)}
\ControlFlowTok{for}\NormalTok{ word, count }\KeywordTok{in}\NormalTok{ word\_frequencies.items():}
    \BuiltInTok{print}\NormalTok{(}\SpecialStringTok{f"\textquotesingle{}}\SpecialCharTok{\{}\NormalTok{word}\SpecialCharTok{\}}\SpecialStringTok{\textquotesingle{}: }\SpecialCharTok{\{}\NormalTok{count}\SpecialCharTok{\}}\SpecialStringTok{"}\NormalTok{)}
\end{Highlighting}
\end{Shaded}

\subsubsection{Skills Checklist}\label{skills-checklist-45}

\begin{itemize}
\tightlist
\item[$\square$]
  I can read text content from a file.
\item[$\square$]
  I can split a string of text into individual words.
\item[$\square$]
  I can use a dictionary to count the frequency of words.
\end{itemize}

\begin{center}\rule{0.5\linewidth}{0.5pt}\end{center}

\subsection{Lesson 9.5: Working with CSV
Data}\label{lesson-9.5-working-with-csv-data}

\begin{Shaded}
\begin{Highlighting}[]
\ImportTok{import}\NormalTok{ csv}

\NormalTok{filename }\OperatorTok{=} \StringTok{\textquotesingle{}students.csv\textquotesingle{}}
\ControlFlowTok{try}\NormalTok{:}
    \ControlFlowTok{with} \BuiltInTok{open}\NormalTok{(filename, }\StringTok{\textquotesingle{}r\textquotesingle{}}\NormalTok{) }\ImportTok{as}\NormalTok{ file\_object:}
\NormalTok{        reader }\OperatorTok{=}\NormalTok{ csv.reader(file\_object)}
\NormalTok{        header\_row }\OperatorTok{=} \BuiltInTok{next}\NormalTok{(reader)}
        \BuiltInTok{print}\NormalTok{(}\SpecialStringTok{f"Columns: }\SpecialCharTok{\{}\NormalTok{header\_row}\SpecialCharTok{\}}\SpecialStringTok{"}\NormalTok{)}
        \ControlFlowTok{for}\NormalTok{ row }\KeywordTok{in}\NormalTok{ reader:}
            \BuiltInTok{print}\NormalTok{(row)}
\ControlFlowTok{except} \PreprocessorTok{FileNotFoundError}\NormalTok{:}
    \BuiltInTok{print}\NormalTok{(}\SpecialStringTok{f"Error: The file }\SpecialCharTok{\{}\NormalTok{filename}\SpecialCharTok{\}}\SpecialStringTok{ was not found."}\NormalTok{)}
\end{Highlighting}
\end{Shaded}

\subsubsection{Manual CSV Parsing}\label{manual-csv-parsing}

\begin{Shaded}
\begin{Highlighting}[]
\ControlFlowTok{with} \BuiltInTok{open}\NormalTok{(}\StringTok{\textquotesingle{}students.csv\textquotesingle{}}\NormalTok{, }\StringTok{\textquotesingle{}r\textquotesingle{}}\NormalTok{) }\ImportTok{as}\NormalTok{ f:}
\NormalTok{    header }\OperatorTok{=}\NormalTok{ f.readline().strip()}
    \ControlFlowTok{for}\NormalTok{ line }\KeywordTok{in}\NormalTok{ f:}
\NormalTok{        data }\OperatorTok{=}\NormalTok{ line.strip().split(}\StringTok{\textquotesingle{},\textquotesingle{}}\NormalTok{)}
\NormalTok{        name }\OperatorTok{=}\NormalTok{ data[}\DecValTok{0}\NormalTok{]}
\NormalTok{        age }\OperatorTok{=}\NormalTok{ data[}\DecValTok{1}\NormalTok{]}
\NormalTok{        grade }\OperatorTok{=}\NormalTok{ data[}\DecValTok{2}\NormalTok{]}
        \BuiltInTok{print}\NormalTok{(}\SpecialStringTok{f"Name: }\SpecialCharTok{\{}\NormalTok{name}\SpecialCharTok{\}}\SpecialStringTok{, Age: }\SpecialCharTok{\{}\NormalTok{age}\SpecialCharTok{\}}\SpecialStringTok{, Grade: }\SpecialCharTok{\{}\NormalTok{grade}\SpecialCharTok{\}}\SpecialStringTok{"}\NormalTok{)}
\end{Highlighting}
\end{Shaded}

\subsubsection{Skills Checklist}\label{skills-checklist-46}

\begin{itemize}
\tightlist
\item[$\square$]
  I can explain what CSV data is.
\item[$\square$]
  I can read a CSV file line by line.
\item[$\square$]
  I can split a line of CSV data into individual values.
\end{itemize}

\newpage

\section{Unit 10: Project-Based Learning --- Building Text-Based
Applications}\label{unit-10-project-based-learning-building-text-based-applications}

\subsection{Overview}\label{overview-10}

This unit guides students through designing and building larger,
interactive text-based applications, integrating all previously learned
concepts.

\textbf{Objectives:}

\begin{itemize}
\tightlist
\item
  Practice modular programming and code organization.
\item
  Develop a text-based interactive program.
\item
  Apply top-down design principles.
\end{itemize}

\begin{center}\rule{0.5\linewidth}{0.5pt}\end{center}

\subsection{Lesson 10.1: Designing Your
Project}\label{lesson-10.1-designing-your-project}

\subsubsection{Reading References}\label{reading-references-28}

\begin{itemize}
\tightlist
\item
  Zelle: Chapter 9, pages 225--235
\item
  Matthes: Part II: Projects, pages 225--226
\end{itemize}

\subsubsection{Top-Down Design}\label{top-down-design}

\textbf{Top-down design} (stepwise refinement) breaks a complex problem
into smaller, more manageable subproblems.

\textbf{The Design Process:}

\textbf{1. Analyze the Problem}

\begin{itemize}
\tightlist
\item
  Identify \textbf{inputs} (information the program needs).
\item
  Identify \textbf{outputs} (results it must produce).
\item
  Define the scope of the project.
\end{itemize}

\textbf{2. Determine the Major Components} Write a high-level algorithm
in plain English:

\begin{verbatim}
show_greeting()
get_user_input()
calculate_cost()
print_ticket()
\end{verbatim}

\textbf{3. Refine Each Component} Repeat the design process for each
component, adding detail until simple enough to code.

\subsubsection{Benefits of Top-Down
Design}\label{benefits-of-top-down-design}

\begin{itemize}
\tightlist
\item
  \textbf{Reduces Complexity} --- Focus on one piece at a time.
\item
  \textbf{Encourages Reusability} --- Functions written for one part may
  be useful elsewhere.
\item
  \textbf{Simplifies Debugging} --- Test each component in isolation.
\item
  \textbf{Improves Collaboration} --- Clear plan for assigning work.
\end{itemize}

\subsubsection{Accessibility Considerations for I/O
Design}\label{accessibility-considerations-for-io-design}

\begin{itemize}
\tightlist
\item
  Provide \textbf{clear, descriptive prompts} for all input.
\item
  Provide \textbf{explicit, labeled output} --- avoid relying on visual
  formatting.
\item
  Print \textbf{confirmation messages} after every action.
\item
  Use consistent \textbf{output separators} like \texttt{-\/-\/-} to
  group information.
\end{itemize}

\subsubsection{Skills Checklist}\label{skills-checklist-47}

\begin{itemize}
\tightlist
\item[$\square$]
  I can break down a larger problem into smaller, manageable parts.
\item[$\square$]
  I can identify the key features required for a project.
\item[$\square$]
  I can plan which data structures will be most suitable for my
  project's data.
\item[$\square$]
  I can outline the main steps or flow of my program.
\end{itemize}

\begin{center}\rule{0.5\linewidth}{0.5pt}\end{center}

\subsection{Lesson 10.2: Building a Text-Based Adventure Game (Part
1)}\label{lesson-10.2-building-a-text-based-adventure-game-part-1}

\subsubsection{Reading References}\label{reading-references-29}

\begin{itemize}
\tightlist
\item
  Matthes: Chapter 6, pages 105--111
\end{itemize}

\subsubsection{Room Definition and
Movement}\label{room-definition-and-movement}

\begin{Shaded}
\begin{Highlighting}[]
\NormalTok{room\_descriptions }\OperatorTok{=}\NormalTok{ \{}
    \StringTok{"start\_room"}\NormalTok{: \{}
        \StringTok{"description"}\NormalTok{: }\StringTok{"You are in a dimly lit cave. There\textquotesingle{}s a narrow passage to the NORTH."}\NormalTok{,}
        \StringTok{"exits"}\NormalTok{: \{}\StringTok{"north"}\NormalTok{: }\StringTok{"tunnel"}\NormalTok{\}}
\NormalTok{    \},}
    \StringTok{"tunnel"}\NormalTok{: \{}
        \StringTok{"description"}\NormalTok{: }\StringTok{"The tunnel is damp and cold. You can go SOUTH back to the cave or EAST to a clearing."}\NormalTok{,}
        \StringTok{"exits"}\NormalTok{: \{}\StringTok{"south"}\NormalTok{: }\StringTok{"start\_room"}\NormalTok{, }\StringTok{"east"}\NormalTok{: }\StringTok{"clearing"}\NormalTok{\}}
\NormalTok{    \},}
    \StringTok{"clearing"}\NormalTok{: \{}
        \StringTok{"description"}\NormalTok{: }\StringTok{"You\textquotesingle{}ve reached a small, sunny clearing. A faint path leads WEST."}\NormalTok{,}
        \StringTok{"exits"}\NormalTok{: \{}\StringTok{"west"}\NormalTok{: }\StringTok{"tunnel"}\NormalTok{\}}
\NormalTok{    \}}
\NormalTok{\}}

\NormalTok{current\_room }\OperatorTok{=} \StringTok{"start\_room"}

\ControlFlowTok{while} \VariableTok{True}\NormalTok{:}
    \BuiltInTok{print}\NormalTok{(}\StringTok{"{-}{-}{-}"}\NormalTok{)}
    \BuiltInTok{print}\NormalTok{(room\_descriptions[current\_room][}\StringTok{"description"}\NormalTok{])}
\NormalTok{    command }\OperatorTok{=} \BuiltInTok{input}\NormalTok{(}\StringTok{"What do you want to do? "}\NormalTok{).lower()}

    \ControlFlowTok{if}\NormalTok{ command }\KeywordTok{in}\NormalTok{ room\_descriptions[current\_room][}\StringTok{"exits"}\NormalTok{]:}
\NormalTok{        current\_room }\OperatorTok{=}\NormalTok{ room\_descriptions[current\_room][}\StringTok{"exits"}\NormalTok{][command]}
    \ControlFlowTok{elif}\NormalTok{ command }\OperatorTok{==} \StringTok{"quit"}\NormalTok{:}
        \BuiltInTok{print}\NormalTok{(}\StringTok{"Exiting game. Goodbye!"}\NormalTok{)}
        \ControlFlowTok{break}
    \ControlFlowTok{else}\NormalTok{:}
        \BuiltInTok{print}\NormalTok{(}\StringTok{"I don\textquotesingle{}t understand that command or you can\textquotesingle{}t go that way."}\NormalTok{)}
\end{Highlighting}
\end{Shaded}

\subsubsection{Skills Checklist}\label{skills-checklist-48}

\begin{itemize}
\tightlist
\item[$\square$]
  I can design a way to represent different locations in a text-based
  game.
\item[$\square$]
  I can implement player movement based on user commands.
\item[$\square$]
  I can update and track the player's current location.
\item[$\square$]
  I can parse simple commands from user input.
\end{itemize}

\begin{center}\rule{0.5\linewidth}{0.5pt}\end{center}

\subsection{Lesson 10.3: Building a Text-Based Adventure Game (Part
2)}\label{lesson-10.3-building-a-text-based-adventure-game-part-2}

\subsubsection{Adding Items, Inventory, and
Puzzles}\label{adding-items-inventory-and-puzzles}

\begin{Shaded}
\begin{Highlighting}[]
\NormalTok{rooms }\OperatorTok{=}\NormalTok{ \{}
    \StringTok{"start"}\NormalTok{: \{}
        \StringTok{"description"}\NormalTok{: }\StringTok{"A dark cave. There\textquotesingle{}s a SHINY KEY on the ground."}\NormalTok{,}
        \StringTok{"exits"}\NormalTok{: \{}\StringTok{"north"}\NormalTok{: }\StringTok{"locked\_door\_room"}\NormalTok{\},}
        \StringTok{"items"}\NormalTok{: [}\StringTok{"shiny key"}\NormalTok{]}
\NormalTok{    \},}
    \StringTok{"locked\_door\_room"}\NormalTok{: \{}
        \StringTok{"description"}\NormalTok{: }\StringTok{"A small chamber with a heavy locked door to the NORTH."}\NormalTok{,}
        \StringTok{"exits"}\NormalTok{: \{}\StringTok{"south"}\NormalTok{: }\StringTok{"start"}\NormalTok{\},}
        \StringTok{"items"}\NormalTok{: []}
\NormalTok{    \}}
\NormalTok{\}}

\NormalTok{current\_room }\OperatorTok{=} \StringTok{"start"}
\NormalTok{player\_inventory }\OperatorTok{=}\NormalTok{ []}
\NormalTok{door\_locked }\OperatorTok{=} \VariableTok{True}

\ControlFlowTok{while} \VariableTok{True}\NormalTok{:}
    \BuiltInTok{print}\NormalTok{(}\StringTok{"}\CharTok{\textbackslash{}n}\StringTok{"} \OperatorTok{+}\NormalTok{ rooms[current\_room][}\StringTok{"description"}\NormalTok{])}
    \ControlFlowTok{if}\NormalTok{ rooms[current\_room][}\StringTok{"items"}\NormalTok{]:}
        \BuiltInTok{print}\NormalTok{(}\SpecialStringTok{f"You see: }\SpecialCharTok{\{}\StringTok{\textquotesingle{}, \textquotesingle{}}\SpecialCharTok{.}\NormalTok{join(rooms[current\_room][}\StringTok{\textquotesingle{}items\textquotesingle{}}\NormalTok{])}\SpecialCharTok{\}}\SpecialStringTok{"}\NormalTok{)}
    \BuiltInTok{print}\NormalTok{(}\SpecialStringTok{f"Your inventory: }\SpecialCharTok{\{}\StringTok{\textquotesingle{}, \textquotesingle{}}\SpecialCharTok{.}\NormalTok{join(player\_inventory) }\ControlFlowTok{if}\NormalTok{ player\_inventory }\ControlFlowTok{else} \StringTok{\textquotesingle{}empty\textquotesingle{}}\SpecialCharTok{\}}\SpecialStringTok{"}\NormalTok{)}

\NormalTok{    command }\OperatorTok{=} \BuiltInTok{input}\NormalTok{(}\StringTok{"What do you want to do? "}\NormalTok{).lower()}

    \ControlFlowTok{if}\NormalTok{ command }\OperatorTok{==} \StringTok{"take shiny key"} \KeywordTok{and} \StringTok{"shiny key"} \KeywordTok{in}\NormalTok{ rooms[current\_room][}\StringTok{"items"}\NormalTok{]:}
\NormalTok{        player\_inventory.append(}\StringTok{"shiny key"}\NormalTok{)}
\NormalTok{        rooms[current\_room][}\StringTok{"items"}\NormalTok{].remove(}\StringTok{"shiny key"}\NormalTok{)}
        \BuiltInTok{print}\NormalTok{(}\StringTok{"You picked up the shiny key."}\NormalTok{)}
    \ControlFlowTok{elif}\NormalTok{ command }\OperatorTok{==} \StringTok{"use shiny key"} \KeywordTok{and}\NormalTok{ current\_room }\OperatorTok{==} \StringTok{"locked\_door\_room"} \KeywordTok{and} \StringTok{"shiny key"} \KeywordTok{in}\NormalTok{ player\_inventory }\KeywordTok{and}\NormalTok{ door\_locked:}
\NormalTok{        door\_locked }\OperatorTok{=} \VariableTok{False}
        \BuiltInTok{print}\NormalTok{(}\StringTok{"You unlocked the door!"}\NormalTok{)}
\NormalTok{        rooms[}\StringTok{"locked\_door\_room"}\NormalTok{][}\StringTok{"exits"}\NormalTok{][}\StringTok{"north"}\NormalTok{] }\OperatorTok{=} \StringTok{"treasure\_room"}
    \ControlFlowTok{elif}\NormalTok{ command }\KeywordTok{in}\NormalTok{ rooms[current\_room][}\StringTok{"exits"}\NormalTok{]:}
        \ControlFlowTok{if}\NormalTok{ command }\OperatorTok{==} \StringTok{"north"} \KeywordTok{and}\NormalTok{ current\_room }\OperatorTok{==} \StringTok{"locked\_door\_room"} \KeywordTok{and}\NormalTok{ door\_locked:}
            \BuiltInTok{print}\NormalTok{(}\StringTok{"The door is locked! You need a key."}\NormalTok{)}
        \ControlFlowTok{else}\NormalTok{:}
\NormalTok{            current\_room }\OperatorTok{=}\NormalTok{ rooms[current\_room][}\StringTok{"exits"}\NormalTok{][command]}
    \ControlFlowTok{elif}\NormalTok{ command }\OperatorTok{==} \StringTok{"quit"}\NormalTok{:}
        \ControlFlowTok{break}
    \ControlFlowTok{else}\NormalTok{:}
        \BuiltInTok{print}\NormalTok{(}\StringTok{"I don\textquotesingle{}t understand that command."}\NormalTok{)}
\end{Highlighting}
\end{Shaded}

\subsubsection{Skills Checklist}\label{skills-checklist-49}

\begin{itemize}
\tightlist
\item[$\square$]
  I can add items to locations and manage a player's inventory.
\item[$\square$]
  I can implement commands for picking up and dropping items.
\item[$\square$]
  I can design and implement a simple puzzle.
\item[$\square$]
  I can define win/lose conditions for the game.
\end{itemize}

\begin{center}\rule{0.5\linewidth}{0.5pt}\end{center}

\subsection{Lesson 10.4: Developing a Command-Line
Tool}\label{lesson-10.4-developing-a-command-line-tool}

\subsubsection{Simple To-Do List Tool}\label{simple-to-do-list-tool}

\begin{Shaded}
\begin{Highlighting}[]
\ImportTok{import}\NormalTok{ json}

\NormalTok{TODO\_FILE }\OperatorTok{=} \StringTok{\textquotesingle{}todos.json\textquotesingle{}}
\NormalTok{todos }\OperatorTok{=}\NormalTok{ []}


\KeywordTok{def}\NormalTok{ load\_todos():}
    \KeywordTok{global}\NormalTok{ todos}
    \ControlFlowTok{try}\NormalTok{:}
        \ControlFlowTok{with} \BuiltInTok{open}\NormalTok{(TODO\_FILE, }\StringTok{\textquotesingle{}r\textquotesingle{}}\NormalTok{) }\ImportTok{as}\NormalTok{ f:}
\NormalTok{            todos }\OperatorTok{=}\NormalTok{ json.load(f)}
        \BuiltInTok{print}\NormalTok{(}\StringTok{"Tasks loaded."}\NormalTok{)}
    \ControlFlowTok{except} \PreprocessorTok{FileNotFoundError}\NormalTok{:}
        \BuiltInTok{print}\NormalTok{(}\StringTok{"No existing tasks found."}\NormalTok{)}
\NormalTok{        todos }\OperatorTok{=}\NormalTok{ []}


\KeywordTok{def}\NormalTok{ save\_todos():}
    \ControlFlowTok{with} \BuiltInTok{open}\NormalTok{(TODO\_FILE, }\StringTok{\textquotesingle{}w\textquotesingle{}}\NormalTok{) }\ImportTok{as}\NormalTok{ f:}
\NormalTok{        json.dump(todos, f, indent}\OperatorTok{=}\DecValTok{4}\NormalTok{)}
    \BuiltInTok{print}\NormalTok{(}\StringTok{"Tasks saved."}\NormalTok{)}


\KeywordTok{def}\NormalTok{ add\_task(task\_description):}
\NormalTok{    todos.append(\{}\StringTok{"description"}\NormalTok{: task\_description, }\StringTok{"completed"}\NormalTok{: }\VariableTok{False}\NormalTok{\})}
    \BuiltInTok{print}\NormalTok{(}\SpecialStringTok{f"Task \textquotesingle{}}\SpecialCharTok{\{}\NormalTok{task\_description}\SpecialCharTok{\}}\SpecialStringTok{\textquotesingle{} added."}\NormalTok{)}


\KeywordTok{def}\NormalTok{ view\_tasks():}
    \ControlFlowTok{if} \KeywordTok{not}\NormalTok{ todos:}
        \BuiltInTok{print}\NormalTok{(}\StringTok{"No tasks in your list."}\NormalTok{)}
        \ControlFlowTok{return}
    \BuiltInTok{print}\NormalTok{(}\StringTok{"{-}{-}{-} Your To{-}Do List {-}{-}{-}"}\NormalTok{)}
    \ControlFlowTok{for}\NormalTok{ i, task }\KeywordTok{in} \BuiltInTok{enumerate}\NormalTok{(todos):}
\NormalTok{        status }\OperatorTok{=} \StringTok{"[X]"} \ControlFlowTok{if}\NormalTok{ task[}\StringTok{"completed"}\NormalTok{] }\ControlFlowTok{else} \StringTok{"[ ]"}
        \BuiltInTok{print}\NormalTok{(}\SpecialStringTok{f"}\SpecialCharTok{\{}\NormalTok{i }\OperatorTok{+} \DecValTok{1}\SpecialCharTok{\}}\SpecialStringTok{. }\SpecialCharTok{\{}\NormalTok{status}\SpecialCharTok{\}}\SpecialStringTok{ }\SpecialCharTok{\{}\NormalTok{task[}\StringTok{\textquotesingle{}description\textquotesingle{}}\NormalTok{]}\SpecialCharTok{\}}\SpecialStringTok{"}\NormalTok{)}
    \BuiltInTok{print}\NormalTok{(}\StringTok{"{-}{-}{-}{-}{-}{-}{-}{-}{-}{-}{-}{-}{-}{-}{-}{-}{-}{-}{-}{-}{-}{-}{-}"}\NormalTok{)}


\KeywordTok{def}\NormalTok{ mark\_completed(task\_number):}
    \ControlFlowTok{try}\NormalTok{:}
\NormalTok{        index }\OperatorTok{=} \BuiltInTok{int}\NormalTok{(task\_number) }\OperatorTok{{-}} \DecValTok{1}
        \ControlFlowTok{if} \DecValTok{0} \OperatorTok{\textless{}=}\NormalTok{ index }\OperatorTok{\textless{}} \BuiltInTok{len}\NormalTok{(todos):}
\NormalTok{            todos[index][}\StringTok{"completed"}\NormalTok{] }\OperatorTok{=} \VariableTok{True}
            \BuiltInTok{print}\NormalTok{(}\SpecialStringTok{f"Task \textquotesingle{}}\SpecialCharTok{\{}\NormalTok{todos[index][}\StringTok{\textquotesingle{}description\textquotesingle{}}\NormalTok{]}\SpecialCharTok{\}}\SpecialStringTok{\textquotesingle{} marked as complete."}\NormalTok{)}
        \ControlFlowTok{else}\NormalTok{:}
            \BuiltInTok{print}\NormalTok{(}\StringTok{"Invalid task number."}\NormalTok{)}
    \ControlFlowTok{except} \PreprocessorTok{ValueError}\NormalTok{:}
        \BuiltInTok{print}\NormalTok{(}\StringTok{"Please enter a valid number."}\NormalTok{)}


\KeywordTok{def}\NormalTok{ main\_menu():}
    \BuiltInTok{print}\NormalTok{(}\StringTok{"}\CharTok{\textbackslash{}n}\StringTok{{-}{-}{-} To{-}Do List Menu {-}{-}{-}"}\NormalTok{)}
    \BuiltInTok{print}\NormalTok{(}\StringTok{"1. Add Task"}\NormalTok{)}
    \BuiltInTok{print}\NormalTok{(}\StringTok{"2. View Tasks"}\NormalTok{)}
    \BuiltInTok{print}\NormalTok{(}\StringTok{"3. Mark Task Complete"}\NormalTok{)}
    \BuiltInTok{print}\NormalTok{(}\StringTok{"4. Save Tasks"}\NormalTok{)}
    \BuiltInTok{print}\NormalTok{(}\StringTok{"5. Load Tasks"}\NormalTok{)}
    \BuiltInTok{print}\NormalTok{(}\StringTok{"6. Quit"}\NormalTok{)}
    \BuiltInTok{print}\NormalTok{(}\StringTok{"{-}{-}{-}{-}{-}{-}{-}{-}{-}{-}{-}{-}{-}{-}{-}{-}{-}{-}{-}{-}{-}{-}{-}"}\NormalTok{)}


\NormalTok{load\_todos()}

\ControlFlowTok{while} \VariableTok{True}\NormalTok{:}
\NormalTok{    main\_menu()}
\NormalTok{    choice }\OperatorTok{=} \BuiltInTok{input}\NormalTok{(}\StringTok{"Enter your choice (1{-}6): "}\NormalTok{)}

    \ControlFlowTok{if}\NormalTok{ choice }\OperatorTok{==} \StringTok{\textquotesingle{}1\textquotesingle{}}\NormalTok{:}
\NormalTok{        desc }\OperatorTok{=} \BuiltInTok{input}\NormalTok{(}\StringTok{"Enter task description: "}\NormalTok{)}
\NormalTok{        add\_task(desc)}
    \ControlFlowTok{elif}\NormalTok{ choice }\OperatorTok{==} \StringTok{\textquotesingle{}2\textquotesingle{}}\NormalTok{:}
\NormalTok{        view\_tasks()}
    \ControlFlowTok{elif}\NormalTok{ choice }\OperatorTok{==} \StringTok{\textquotesingle{}3\textquotesingle{}}\NormalTok{:}
\NormalTok{        view\_tasks()}
\NormalTok{        task\_num }\OperatorTok{=} \BuiltInTok{input}\NormalTok{(}\StringTok{"Enter task number to mark complete: "}\NormalTok{)}
\NormalTok{        mark\_completed(task\_num)}
    \ControlFlowTok{elif}\NormalTok{ choice }\OperatorTok{==} \StringTok{\textquotesingle{}4\textquotesingle{}}\NormalTok{:}
\NormalTok{        save\_todos()}
    \ControlFlowTok{elif}\NormalTok{ choice }\OperatorTok{==} \StringTok{\textquotesingle{}5\textquotesingle{}}\NormalTok{:}
\NormalTok{        load\_todos()}
    \ControlFlowTok{elif}\NormalTok{ choice }\OperatorTok{==} \StringTok{\textquotesingle{}6\textquotesingle{}}\NormalTok{:}
\NormalTok{        save\_todos()}
        \BuiltInTok{print}\NormalTok{(}\StringTok{"Goodbye!"}\NormalTok{)}
        \ControlFlowTok{break}
    \ControlFlowTok{else}\NormalTok{:}
        \BuiltInTok{print}\NormalTok{(}\StringTok{"Invalid choice. Please try again."}\NormalTok{)}
\end{Highlighting}
\end{Shaded}

\subsubsection{Skills Checklist}\label{skills-checklist-50}

\begin{itemize}
\tightlist
\item[$\square$]
  I can design an interactive menu system for a tool.
\item[$\square$]
  I can use a \texttt{while} loop to maintain continuous user
  interaction.
\item[$\square$]
  I can break down the tool's functionality into separate functions.
\item[$\square$]
  I can implement file I/O to save and load data.
\end{itemize}

\begin{center}\rule{0.5\linewidth}{0.5pt}\end{center}

\subsection{Lesson 10.5: Refinement and User
Experience}\label{lesson-10.5-refinement-and-user-experience}

\subsubsection{Principles for Accessible Text
Programs}\label{principles-for-accessible-text-programs}

\begin{enumerate}
\def\labelenumi{\arabic{enumi}.}
\tightlist
\item
  \textbf{Clear prompts} --- Tell the user exactly what to enter and in
  what format.
\item
  \textbf{Consistent output} --- Use the same format for similar
  information.
\item
  \textbf{Graceful error handling} --- Never let the program crash on
  bad input.
\item
  \textbf{Helpful error messages} --- Tell the user what went wrong and
  how to fix it.
\item
  \textbf{``Help'' commands} --- Provide a way to list available
  commands.
\item
  \textbf{Confirmation messages} --- Confirm every action taken.
\end{enumerate}

\subsubsection{Robust Input Validation}\label{robust-input-validation}

\begin{Shaded}
\begin{Highlighting}[]
\KeywordTok{def}\NormalTok{ get\_positive\_number(prompt):}
    \CommentTok{"""Prompts user until a positive integer is entered."""}
    \ControlFlowTok{while} \VariableTok{True}\NormalTok{:}
        \ControlFlowTok{try}\NormalTok{:}
\NormalTok{            num }\OperatorTok{=} \BuiltInTok{int}\NormalTok{(}\BuiltInTok{input}\NormalTok{(prompt))}
            \ControlFlowTok{if}\NormalTok{ num }\OperatorTok{\textgreater{}} \DecValTok{0}\NormalTok{:}
                \ControlFlowTok{return}\NormalTok{ num}
            \ControlFlowTok{else}\NormalTok{:}
                \BuiltInTok{print}\NormalTok{(}\StringTok{"Please enter a positive number."}\NormalTok{)}
        \ControlFlowTok{except} \PreprocessorTok{ValueError}\NormalTok{:}
            \BuiltInTok{print}\NormalTok{(}\StringTok{"Invalid input. Please enter a whole number."}\NormalTok{)}


\NormalTok{age }\OperatorTok{=}\NormalTok{ get\_positive\_number(}\StringTok{"Enter your age: "}\NormalTok{)}
\BuiltInTok{print}\NormalTok{(}\SpecialStringTok{f"You entered age: }\SpecialCharTok{\{}\NormalTok{age}\SpecialCharTok{\}}\SpecialStringTok{"}\NormalTok{)}
\end{Highlighting}
\end{Shaded}

\subsubsection{Skills Checklist}\label{skills-checklist-51}

\begin{itemize}
\tightlist
\item[$\square$]
  I can write clear, descriptive prompts for user input.
\item[$\square$]
  I can handle invalid input gracefully using \texttt{try-except}.
\item[$\square$]
  I can provide helpful error messages.
\end{itemize}

\begin{center}\rule{0.5\linewidth}{0.5pt}\end{center}

\section{Unit 11: Capstone Project --- Accessible Personal Data
Manager}\label{unit-11-capstone-project-accessible-personal-data-manager}

\subsection{Overview}\label{overview-11}

Students will develop a fully functional, screen-reader-accessible
graphical application using NiceGUI that allows users to manage a
collection of structured personal data.

\textbf{Project Goal:} Build a simple application to manage a list of
structured items.

\textbf{Project Options:}

\begin{itemize}
\tightlist
\item
  Contact Manager: name, phone, email, notes
\item
  Book/Movie Tracker: title, author/director, genre, status
\item
  Simple To-Do List: task description, due date, status
\item
  Simple Inventory: item name, quantity, location, notes
\end{itemize}

\begin{center}\rule{0.5\linewidth}{0.5pt}\end{center}

\subsection{Core Functionality
Requirements}\label{core-functionality-requirements}

\begin{enumerate}
\def\labelenumi{\arabic{enumi}.}
\tightlist
\item
  \textbf{Add New Item} --- A form to input details for a new item.
\item
  \textbf{View All Items} --- Display items in a clear, accessible list
  format.
\item
  \textbf{Edit Item} --- Select an existing item and modify its details.
\item
  \textbf{Delete Item} --- Remove an item from the collection.
\item
  \textbf{Save/Load Data} --- Persist the data to a JSON file.
\end{enumerate}

\begin{center}\rule{0.5\linewidth}{0.5pt}\end{center}

\subsection{Setup}\label{setup}

\begin{Shaded}
\begin{Highlighting}[]
\ExtensionTok{pip}\NormalTok{ install nicegui}
\end{Highlighting}
\end{Shaded}

\begin{center}\rule{0.5\linewidth}{0.5pt}\end{center}

\subsection{Guided Implementation}\label{guided-implementation}

\subsubsection{Lesson 11.1: Introduction to
NiceGUI}\label{lesson-11.1-introduction-to-nicegui}

\textbf{Task 1:} Run the \texttt{app\_skeleton.py} file and navigate
using the keyboard and screen reader to understand the existing
structure.

\textbf{Task 2:} Define your data model class:

\begin{Shaded}
\begin{Highlighting}[]
\KeywordTok{class}\NormalTok{ Contact:}
    \KeywordTok{def} \FunctionTok{\_\_init\_\_}\NormalTok{(}\VariableTok{self}\NormalTok{, name, phone, email, notes):}
        \VariableTok{self}\NormalTok{.name }\OperatorTok{=}\NormalTok{ name}
        \VariableTok{self}\NormalTok{.phone }\OperatorTok{=}\NormalTok{ phone}
        \VariableTok{self}\NormalTok{.email }\OperatorTok{=}\NormalTok{ email}
        \VariableTok{self}\NormalTok{.notes }\OperatorTok{=}\NormalTok{ notes}
\end{Highlighting}
\end{Shaded}

\textbf{Task 3:} Implement data persistence using the \texttt{json}
module within the \texttt{load\_data()} and \texttt{save\_data()}
functions.

\subsubsection{Lesson 11.2: Data Model and
Storage}\label{lesson-11.2-data-model-and-storage}

\textbf{Task 1:} Complete the \texttt{add\_item\_logic()} function:

\begin{enumerate}
\def\labelenumi{\arabic{enumi}.}
\tightlist
\item
  Get values from input fields using \texttt{.value} attribute.
\item
  Validate (e.g., name cannot be empty).
\item
  Provide accessible feedback:
  \texttt{ui.notify("Name\ is\ required!",\ type=\textquotesingle{}negative\textquotesingle{})}.
\end{enumerate}

\textbf{Task 2:} Create a new item object and add it to
\texttt{all\_items}.

\textbf{Task 3:} Clear input fields and provide success notification.

\subsubsection{Lesson 11.3--11.5: Displaying and Managing
Items}\label{lesson-11.311.5-displaying-and-managing-items}

\textbf{Implement \texttt{update\_item\_display()}:}

\begin{enumerate}
\def\labelenumi{\arabic{enumi}.}
\tightlist
\item
  Clear the display: \texttt{items\_display\_column.clear()}.
\item
  Loop through \texttt{all\_items} and create \texttt{ui.card} for each.
\item
  Update \texttt{item\_selection\_dropdown.options}.
\end{enumerate}

\textbf{Implement \texttt{update\_item\_logic()}:}

\begin{enumerate}
\def\labelenumi{\arabic{enumi}.}
\tightlist
\item
  Find the item in \texttt{all\_items}.
\item
  Update attributes from edit form fields.
\item
  Call \texttt{update\_item\_display()} and \texttt{save\_data()}.
\end{enumerate}

\textbf{Implement \texttt{delete\_selected\_item\_logic()}:}

\begin{enumerate}
\def\labelenumi{\arabic{enumi}.}
\tightlist
\item
  Find and remove the item from \texttt{all\_items}.
\item
  Call \texttt{update\_item\_display()}.
\item
  Provide accessible feedback.
\end{enumerate}

\begin{center}\rule{0.5\linewidth}{0.5pt}\end{center}

\subsection{Assessment Criteria}\label{assessment-criteria}

\subsubsection{Functionality (60\%)}\label{functionality-60}

\begin{itemize}
\tightlist
\item[$\square$]
  Ability to add new items with all specified attributes.
\item[$\square$]
  Correct display of all current items.
\item[$\square$]
  Correct editing of existing items.
\item[$\square$]
  Correct deletion of items.
\item[$\square$]
  Accurate saving and loading to/from \texttt{data.json}.
\item[$\square$]
  No major bugs or crashes during typical usage.
\end{itemize}

\subsubsection{Code Quality \& Structure
(20\%)}\label{code-quality-structure-20}

\begin{itemize}
\tightlist
\item[$\square$]
  Clear and concise Python code.
\item[$\square$]
  Appropriate use of functions and/or classes.
\item[$\square$]
  Meaningful variable and function names.
\item[$\square$]
  Adequate comments where complex logic exists.
\item[$\square$]
  Proper indentation.
\end{itemize}

\subsubsection{Accessibility \& User Experience
(15\%)}\label{accessibility-user-experience-15}

\begin{itemize}
\tightlist
\item[$\square$]
  All interactive elements are correctly labeled for screen readers.
\item[$\square$]
  Clear and consistent user feedback via \texttt{ui.notify()}.
\item[$\square$]
  Logical flow and presentation of information.
\item[$\square$]
  No visual-only cues for critical information.
\end{itemize}

\subsubsection{README and Reflection
(5\%)}\label{readme-and-reflection-5}

\begin{itemize}
\tightlist
\item[$\square$]
  Complete and clear README.
\item[$\square$]
  Thoughtful reflection on development and accessibility challenges.
\end{itemize}

\begin{center}\rule{0.5\linewidth}{0.5pt}\end{center}

\subsection{Deliverables}\label{deliverables}

\begin{enumerate}
\def\labelenumi{\arabic{enumi}.}
\tightlist
\item
  \texttt{app.py} --- Complete NiceGUI application code.
\item
  \texttt{data.json} --- Example JSON file demonstrating saved data.
\item
  \texttt{README.md} --- Including:

  \begin{itemize}
  \tightlist
  \item
    Project title and student name.
  \item
    Description of the application.
  \item
    Instructions on how to run: \texttt{python\ app.py}.
  \item
    Instructions for interacting with a screen reader.
  \item
    Reflection on challenges faced.
  \end{itemize}
\end{enumerate}

\begin{center}\rule{0.5\linewidth}{0.5pt}\end{center}

\subsection{Supplemental: AI Tools and Version
Control}\label{supplemental-ai-tools-and-version-control}

\subsubsection{AI-Powered Coding Tools}\label{ai-powered-coding-tools}

\textbf{Benefits:}

\begin{itemize}
\tightlist
\item
  Increased productivity and speed.
\item
  Syntax and idiom discovery.
\item
  Rapid prototyping.
\end{itemize}

\textbf{Critical Warning:} AI tools can produce code that looks correct
but contains subtle bugs. \textbf{Always:}

\begin{enumerate}
\def\labelenumi{\arabic{enumi}.}
\tightlist
\item
  Read every line of AI-generated code.
\item
  Verify the logic and intent.
\item
  Look for familiar error types (indentation, name errors, type errors).
\item
  Run and test incrementally.
\item
  Add your own comments and docstrings.
\end{enumerate}

\textbf{The Golden Rule:} You are responsible for your code. Treat AI
suggestions as a first draft that needs meticulous editing.

\subsubsection{Git and Version Control}\label{git-and-version-control}

\textbf{Git} is the most widely used distributed version control system,
invented in 2005 by Linus Torvalds for the Linux kernel project.

\textbf{Core workflow:}

\begin{enumerate}
\def\labelenumi{\arabic{enumi}.}
\tightlist
\item
  \textbf{Working directory} --- Make changes to files.
\item
  \textbf{Staging area} --- \texttt{git\ add} --- Stage changes.
\item
  \textbf{Repository} --- \texttt{git\ commit} --- Save a permanent
  snapshot.
\item
  \textbf{Remote} --- \texttt{git\ push} / \texttt{git\ pull} --- Sync
  with GitHub.
\end{enumerate}

\textbf{Why Git is especially valuable with AI code:}

\begin{itemize}
\tightlist
\item
  \textbf{Rollback} --- If AI-generated code breaks your app, revert
  instantly.
\item
  \textbf{\texttt{git\ blame}} --- See exactly which commit introduced a
  bug.
\item
  \textbf{Branches} --- Test AI experiments in isolation without
  affecting main code.
\item
  \textbf{Pull requests} --- Review AI-generated code before merging.
\end{itemize}

\textbf{Basic commands:}

\begin{Shaded}
\begin{Highlighting}[]
\FunctionTok{git}\NormalTok{ init                    }\CommentTok{\# Initialize a repository}
\FunctionTok{git}\NormalTok{ add .                   }\CommentTok{\# Stage all changes}
\FunctionTok{git}\NormalTok{ commit }\AttributeTok{{-}m} \StringTok{"Add feature"} \CommentTok{\# Save a snapshot}
\FunctionTok{git}\NormalTok{ log                     }\CommentTok{\# View history}
\FunctionTok{git}\NormalTok{ diff                    }\CommentTok{\# See what changed}
\FunctionTok{git}\NormalTok{ checkout }\AttributeTok{{-}{-}}\NormalTok{ file.py     }\CommentTok{\# Discard changes to a file}
\end{Highlighting}
\end{Shaded}


\end{document}